\documentclass[11pt,a4paper]{amsart}

\usepackage[margin=1in]{geometry}
\usepackage[T1]{fontenc}
\usepackage[utf8]{inputenc}
\usepackage{lmodern}
\usepackage[protrusion=true,expansion=false]{microtype}
\emergencystretch=3em
\usepackage{amsmath,amssymb,amsthm,mathtools}
\usepackage{enumitem}
\usepackage{float}
\usepackage{aliascnt}
\usepackage{booktabs}
\usepackage{array}
\usepackage{tabularx}
\usepackage[colorlinks=true,linkcolor=blue,citecolor=blue,urlcolor=blue]{hyperref}
\usepackage[capitalize,nameinlink]{cleveref}

\numberwithin{equation}{section}

\newtheorem{theorem}{Theorem}[section]
\newaliascnt{proposition}{theorem}
\newtheorem{proposition}[proposition]{Proposition}
\aliascntresetthe{proposition}
\newaliascnt{lemma}{theorem}
\newtheorem{lemma}[lemma]{Lemma}
\aliascntresetthe{lemma}
\newaliascnt{corollary}{theorem}
\newtheorem{corollary}[corollary]{Corollary}
\aliascntresetthe{corollary}
\newaliascnt{definition}{theorem}
\newtheorem{definition}[definition]{Definition}
\aliascntresetthe{definition}
\newaliascnt{assumption}{theorem}
\newtheorem{assumption}[assumption]{Assumption}
\aliascntresetthe{assumption}
\theoremstyle{remark}
\newaliascnt{remark}{theorem}
\newtheorem{remark}[remark]{Remark}
\aliascntresetthe{remark}

\crefname{theorem}{Theorem}{Theorems}
\Crefname{theorem}{Theorem}{Theorems}
\crefname{proposition}{Proposition}{Propositions}
\Crefname{proposition}{Proposition}{Propositions}
\crefname{lemma}{Lemma}{Lemmas}
\Crefname{lemma}{Lemma}{Lemmas}
\crefname{corollary}{Corollary}{Corollaries}
\Crefname{corollary}{Corollary}{Corollaries}
\crefname{definition}{Definition}{Definitions}
\Crefname{definition}{Definition}{Definitions}
\crefname{assumption}{Assumption}{Assumptions}
\Crefname{assumption}{Assumption}{Assumptions}
\crefname{remark}{Remark}{Remarks}
\Crefname{remark}{Remark}{Remarks}

\DeclareMathOperator{\RS}{RS}
\DeclareMathOperator{\MCA}{MCA}
\DeclareMathOperator{\CA}{CA}
\DeclareMathOperator{\rank}{rank}
\DeclareMathOperator{\Span}{Span}
\DeclareMathOperator{\Spec}{Spec}
\DeclareMathOperator{\Hom}{Hom}
\DeclareMathOperator{\im}{im}
\DeclareMathOperator{\codim}{codim}
\DeclareMathOperator{\supp}{supp}
\DeclareMathOperator{\Fib}{Fib}
\DeclareMathOperator{\Sym}{Sym}
\DeclareMathOperator{\Gr}{Gr}
\DeclareMathOperator{\Ent}{H}
\DeclareMathOperator{\ord}{ord}
\DeclareMathOperator{\Tr}{Tr}
\DeclareMathOperator{\Gal}{Gal}
\DeclareMathOperator{\Res}{Res}
\DeclareMathOperator{\Aut}{Aut}
\DeclareMathOperator{\diag}{diag}
\DeclareMathOperator{\wt}{wt}

\newcommand{\F}{\mathbb F}
\newcommand{\B}{\mathbb B}
\newcommand{\K}{\mathbb K}
\newcommand{\Z}{\mathbb Z}
\newcommand{\A}{\mathbb A}
\newcommand{\Pj}{\mathbb P}
\newcommand{\C}{\mathbb C}
\newcommand{\Prob}{\mathbb P}
\newcommand{\E}{\mathbb E}
\newcommand{\eps}{\varepsilon}
\newcommand{\Omc}{\Omega^{\circ}}
\newcommand{\Ccal}{\mathcal C}
\newcommand{\Dcal}{\mathcal D}
\newcommand{\Ecal}{\mathcal E}
\newcommand{\Fcal}{\mathcal F}
\newcommand{\Gcal}{\mathcal G}
\newcommand{\Hcal}{\mathcal H}
\newcommand{\Ical}{\mathcal I}
\newcommand{\Jcal}{\mathcal J}
\newcommand{\Lcal}{\mathcal L}
\newcommand{\Mcal}{\mathcal M}
\newcommand{\Ncal}{\mathcal N}
\newcommand{\Pcal}{\mathcal P}
\newcommand{\Qcal}{\mathcal Q}
\newcommand{\Rcal}{\mathcal R}
\newcommand{\Scal}{\mathcal S}
\newcommand{\Tcal}{\mathcal T}
\newcommand{\Ucal}{\mathcal U}
\newcommand{\Vcal}{\mathcal V}
\newcommand{\Wcal}{\mathcal W}
\newcommand{\Xcal}{\mathcal X}
\newcommand{\Ycal}{\mathcal Y}
\newcommand{\Zcal}{\mathcal Z}
\newcommand{\Gord}{\Gamma^{\rm ord}}
\newcommand{\Gsid}{\Gamma^{\rm sid}}
\newcommand{\Gprim}{\Gamma^{\rm prim}}
\newcommand{\barN}{\overline N}
\newcommand{\Renyi}{R\'enyi}
\newcommand{\Hb}{H_2}
\newcommand{\UU}{\mathbb U}
\newcommand{\abs}[1]{\lvert #1\rvert}
\newcommand{\norm}[1]{\lVert #1\rVert}
\newcommand{\floor}[1]{\left\lfloor #1\right\rfloor}
\newcommand{\ceil}[1]{\left\lceil #1\right\rceil}
\newcommand{\defeq}{\coloneqq}
\newcommand{\one}{\mathbf 1}
\newcommand{\Fq}{\mathbb F_q}
\newcolumntype{P}[1]{>{\raggedright\arraybackslash}p{#1}}
\newcolumntype{Y}{>{\raggedright\arraybackslash}X}


\title{Shortening Bounds for Reed--Solomon MCA}
\author{Przemek Chojecki}
\thanks{Affiliation: Ulam.ai. Public profile: \url{https://github.com/przchojecki}.}
\address{Ulam.ai}
\email{prz.chojecki@gmail.com}
\urladdr{https://github.com/przchojecki}
\hypersetup{pdftitle={Shortening Bounds for Reed--Solomon MCA},
pdfauthor={Przemek Chojecki},
pdfsubject={Finite and asymptotic bounds for Reed--Solomon mutual correlated agreement},
pdfkeywords={Reed-Solomon codes, mutual correlated agreement, Proximity Prize, agreement-set shortening, MDS codes, circle codes}}
\subjclass[2020]{94B35, 11T71, 05B40}
\keywords{Reed--Solomon codes, mutual correlated agreement, Proximity Prize,
agreement-set shortening, MDS codes, soundness gaps, circle codes}
\date{July 17, 2026}

\begin{document}
\raggedbottom

\begin{abstract}
We derive an explicit exponent \(\Psi_\rho\) that bounds the Reed--Solomon MCA bad-slope numerator at every fixed relative radius between Johnson and capacity.  The resulting positive-relative-radius exponential-budget safe-frontier certificate strictly improves the smallest-test MDS exponent and gives a constant post-Johnson radius for every positive usable budget exponent.  An all-test-size MDS circuit-incidence envelope gives exact large-field capacity plateaux, improved adjacent thresholds, and four explicit length-\(512\), \(2^{-128}\)-secure smooth multiplicative certificates beyond Johnson.  An exact CA--MCA decomposition gives challenge-restricted, endpoint-exact linear-budget thresholds through asymptotically half the minimum distance, and monomial equivalence transfers the applicable bounds to circle presentations.  A Gowers--cube argument proves primitive max-fiber flatness from an image-normalized Sidon payment at an accessible moment order.  The unrestricted subexponential-budget smooth/circle frontier remains open because shortening has positive exponential cost and the analytic payment, residual ray compiler, profile add-back, and matching attacks are not yet available.
\end{abstract}

\maketitle
\setcounter{tocdepth}{1}
\tableofcontents
\clearpage

\section{Introduction and main results}\label{sec:introduction}

By a finite Reed--Solomon \emph{row} we mean data
\((\F,D,k)\), where \(\F\) is a finite field,
\(D\subseteq\F\) is an \emph{evaluation domain} of size \(n\), and
\(1\le k<n\).  The associated code
\(C=\RS_\F(D,k)\subseteq\F^D\) consists of the evaluations on \(D\) of
polynomials of degree less than \(k\).  Thus a support is simply a subset
of the coordinate set \(D\).  A word \(u\in\F^D\) is
\emph{explained on a support} \(S\subseteq D\) if its restriction to
\(S\) agrees with one such degree-less-than-\(k\) polynomial.
When the evaluation points lie in a designated coefficient subfield, we
write \(D\subseteq\B\subseteq\F\); \(\B\) controls locator coefficients
and domain structure, while \(\F\) controls codeword coefficients and
slopes.
We call a linear \([n,k]\) code \emph{maximum-distance separable} (MDS)
when its minimum distance is \(n-k+1\), equivalently when puncturing to
every \(k\)-set is an isomorphism onto \(\F^k\).  Such a \(k\)-set is an
\emph{information set}.  The \emph{support atlas} at agreement \(a\ge k+1\) is the
injection obtained by assigning to each bad slope a noncommon exact-\(a\)
support.  Given a punctured generator matrix and one additional word, the
\emph{augmented-word matroid} is the column matroid of the matrix obtained by
appending that word as one further row.  A rank-\(s\) matroid is \emph{paving}
when every set of at most \(s-1\) elements is independent.  The
\emph{tangent construction} is the explicit received-line construction that
produces at least \(\min\{\abs\Gamma,n-a+1\}\) bad slopes at agreement
\(a\); it is proved in \cref{prop:universal-tangent-floor}.  A
\emph{profile} is a declared witness mechanism, and a \emph{first-match
atlas} is a disjoint assignment of witnesses to the first applicable
profile.  Such an atlas is \emph{witness-exhaustive} only if every bad slope
is assigned to a profile.  A \emph{ray} is a projective direction, whereas
MCA counts distinct affine slopes in the declared challenge set; a
\emph{ray compiler} is a theorem that bounds this projection.  Finally, the
\emph{soundness gap} in bits is the interval between a rigorously proved
security lower bound and an available attack upper bound.  These descriptive
terms supply no bound unless the corresponding properties are proved.
Throughout,
\(h(x)=-x\log x-(1-x)\log(1-x)\) denotes binary entropy in natural
units, and \(H_2(x)=h(x)/\log 2\) denotes binary entropy in bits, with the
usual convention \(0\log0=0\).

A received pair \(r=(r_0,r_1)\in(\F^D)^2\) determines the
\emph{received line} \(\{r_0+\gamma r_1:\gamma\in\F\}\).  The pair is
simultaneously explained on \(S\) if each of the two received words has an explanation
there, possibly by different polynomials.  A finite slope \(\gamma\) is
\emph{support-wise MCA-bad at agreement \(a\)} if some support
\(S\subseteq D\), \(\abs S\ge a\), explains
\(r_0+\gamma r_1\) but does not simultaneously explain the pair.  The
same support is used in both tests; this is the support-wise convention.
A tuple \((\gamma,S,h)\), where \(\deg h<k\) and
\(h|_S=(r_0+\gamma r_1)|_S\), is a \emph{witness} to that bad slope when
the pair is not simultaneously explained on \(S\), or equivalently is
\emph{support-wise nontrivial} there.  Finally,
\(B_C^{\rm MCA}(a)\) is the maximum, over all received pairs, of the
number of distinct bad slopes.  It counts slopes, not witnesses or
supports.

For a general linear code \(C\subseteq\F^D\), the phrase ``\(u\) is
explained on \(S\)'' means \(u|_S\in C|_S\); the MCA definition above
then applies verbatim.  For a nonempty slope set \(\Gamma\subseteq\F\),
its ordinary correlated-agreement numerator
\(B_{C,\Gamma}^{\rm CA}(a)\) is the same bad-slope maximum restricted to
received pairs that are not simultaneously explained on any support of
size at least \(a\).  These conventions make the MDS statement below
well-defined before the detailed decomposition in \cref{sec:witnesses}.

This is the affine-line definition used by the Proximity Prize
\cite{ProximityPrize} and by Arnon--Boneh--Fenzi
\cite[Definition~4.3]{ABF26};
\cref{lem:mca-definition-bridge} records its exact maximal-domain
reformulation.  For a real radius \(\delta\), a witnessing
support has size at least \((1-\delta)n\).  Since support size is integral,
put
\[
 r(\delta)=\floor{\delta n},\qquad a(\delta)=n-r(\delta).
\]
Thus the error is a right-continuous integer staircase.  In particular, a
statement about a ``largest'' real radius needs an endpoint convention;
in the exact rows below the safe set is half open and has a supremum but no
largest element.

Here is the asymptotic formulation used throughout.  If
\(\varnothing\ne\Gamma\subseteq\F\) is the \emph{challenge set}, meaning
the set of slopes allowed in the test, let
\(B_{C,\Gamma}^{\rm MCA}(a)\) denote the same maximum with
\(\gamma\in\Gamma\), and for every \(0\le\delta\le1\) put
\[
 e_{\rm MCA,\Gamma}(C,\delta)
   =\frac{B_{C,\Gamma}^{\rm MCA}
      (n-\floor{\delta n})}{\abs\Gamma}.
 \tag{1.1}\label{eq:intro-normalized-mca}
\]
For the standard full-field challenge we abbreviate
\(e_{\rm MCA}=e_{\rm MCA,\F}\).  We keep two endpoint conventions separate.
The customary \emph{subcapacity} quantities are
\[
 \begin{split}
 \delta_{C,\mathrm{sub,sup}}^*(\eps)
   &=\sup\{0\le\delta<1-k/n:e_{\rm MCA}(C,\delta)\le\eps\},\\
 \delta_{C,\mathrm{sub,grid}}^*(\eps)
   &=\max\left\{\frac rn:0\le r\le n-k-1,
          e_{\rm MCA}(C,r/n)\le\eps\right\}.
 \end{split}                                      \tag{SUB}
 \label{eq:intro-subcapacity-thresholds}
\]
The official Prize interval is \([0,1]\), so we also put
\[
 \begin{split}
 \delta_{C,\mathrm{off,sup}}^*(\eps)
   &=\sup\{0\le\delta\le1:e_{\rm MCA}(C,\delta)\le\eps\},\\
 \delta_{C,\mathrm{off,grid}}^*(\eps)
   &=\max\left\{\frac rn:0\le r\le n,
          e_{\rm MCA}(C,r/n)\le\eps\right\}.
 \end{split}                                      \tag{OFF}
\label{eq:intro-official-thresholds}
\]
For a restricted challenge set we insert \(\Gamma\) after \(C\) in this
notation and replace \(e_{\rm MCA}\) by \(e_{\rm MCA,\Gamma}\); omission of
\(\Gamma\) always means the full-field challenge.
For an empty safe set we assign the supremum value zero and leave the grid
maximum undefined.  Otherwise, if the
first unsafe integer radius is \(r_0\) before capacity, the grid value is
\((r_0-1)/n\), the real supremum is \(r_0/n\), and the latter endpoint is
unsafe.  The exceptional all-radius case, in which the official maximum is
attained at one, is isolated below.
By a \emph{row sequence} we mean a family
\(C_n=\RS_{\F_n}(D_n,k_n)\) indexed so that \(\abs{D_n}=n\); the fields,
domains, and dimensions may all vary with \(n\).  Every claimed upper
bound is uniform over received pairs, whereas a lower construction need
only exhibit the stated received pair.  For such a sequence,
challenge sets
\(\Gamma_n\), and targets \(0\le\eps_n\le1\), define
\[
 \begin{split}
 B_n^*&=\floor{\eps_n\abs{\Gamma_n}},\\
 a_n^*&=\min\{a\in\{k_n+1,\ldots,n\}:
              B_{C_n,\Gamma_n}^{\rm MCA}(a)\le B_n^*\},\\
 \delta_{n,\mathrm{sub,grid}}^*&=1-\frac{a_n^*}{n},\qquad
 \delta_{n,\mathrm{sub,sup}}^*=\frac{n-a_n^*+1}{n},
 \end{split}
 \tag{1.2}\label{eq:intro-asymptotic-threshold}
\]
whenever the set defining \(a_n^*\) is nonempty.  Hereafter we write
\(\delta_{n,\mathrm{sub,grid}}^*=1-a_n^*/n\) for the largest safe
subcapacity grid radius; the corresponding continuous
supremum is one grid step to the right whenever \(a_n^*-1\) is the first
unsafe agreement.  Every suppressed symbol \(\delta_{n,*}^*\) is relative
to the displayed pair \((\Gamma_n,\eps_n)\).  Thus the asymptotic RS--MCA problem is to determine
\(\delta_{n,\mathrm{sub,grid}}^*\), or equivalently the first
\emph{safe agreement}: agreement \(a\) is safe when its bad-slope
numerator is at most \(B_n^*\), and unsafe otherwise.  Since the numerator
is nonincreasing in \(a\), the resulting safe/unsafe boundary is the
\emph{MCA frontier}.  This formulation retains the target and challenge
size; replacing them by a zero-exponent slogan loses real information when
the target is small.

\paragraph{Comparison with the literature and organization.}
For arbitrary prescribed Reed--Solomon domains, the proved positive theory
has a linear bad-slope numerator on every compact subinterval below the
Johnson radius \(J_\rho=1-\sqrt\rho\).  We use the list-correlated-agreement
theorem of BCHKS \cite[Thm.~4.6]{BCHKS25}, specialized to two received
coordinates and reindexed from degree at most \(k-1\) to dimension \(k\).
The source states this theorem and sketches the generalization following
its statement; we use only its displayed \(M=1\) specialization and mark it
as an imported input.
Hab\"ock's earlier note \cite{Habock25} outlines the quadratic predecessor
but is not load-bearing here.  Jo \cite{Jo26} proves an all-test-size
circuit-incidence envelope for MDS codes and an exact agreement-set
shortening theorem for Reed--Solomon codes.  The former contains both the
support-atlas and paving bounds.  Jo also combines shortening with BCHKS to
prove polynomial numerators at every fixed number of integer steps beyond
Johnson for each fixed reciprocal-integer rate; we use his transfer below
to derive the positive-relative-radius exponential-budget exponent
\(\Psi_\rho\).  Below half
the relative minimum distance, the BCHKS two-radius theorem, the exact
CA--MCA decomposition, and the tangent construction already imply the
\(r+1\) staircase.  Our contribution in that range is the finite,
challenge-restricted, endpoint-exact formulation, not a new constant-rate
phenomenon.  We retain an independent augmented-word matroid interpretation
of the smallest-test specialization of Jo's envelope and derive the
positive-relative-radius exponential-budget safe-frontier certificate in
the notation relevant to the Prize.  A separate
parameter-retained interpolation audit is confined to
\cref{app:retained-interpolation}; it is conditional on one explicit lift
assumption and is not used by any unconditional main theorem.

Capacity results for subspace-design codes and random-evaluation ensembles
do not settle arbitrary smooth or circle rows
\cite{GoyalGuruswami25,GoyalGuruswamiSunWootters26}.  Conversely, explicit
near-capacity obstructions for structured domains
\cite{CS25,KKH26,Kambire26} prevent a one-sided profile calculation from
being promoted to a complete frontier.  On the attack side, Gao, Yang, Xu,
and Kan \cite{GaoYangXuKan26} constructively convert list-decoding
counterexamples into MCA lower bounds for related linear codes, with an
RS specialization that changes at most one evaluation point; this is
complementary to our fixed-code finite-MDS upper bounds.  In particular, the identity equation
\(H_2(\rho+g)-\beta g=\tau\) proposed in the source repository
\cite{RSMCARepo} is an attack crossing
unless a witness-exhaustive profile atlas, its Fourier or Sidon payments,
the residual ray compiler, and a matching safe reserve have all been proved.
One exact implication inside that program is established here.  On a
Boolean support leaf, an image-normalized Sidon payment at an accessible
moment order \(q_N\), meaning \(\log L_N/q_N=o(N)\), forces any positive-rate
max-fiber excess into a high-energy fiber; the Balog--Szemer\'edi--Gowers
theorem and Boolean-cube difference growth then rule it out.  This
primitive inverse step is proved in
\cref{sec:primitive-inverse}, but its Sidon premise is not a consequence of
smooth or circle folding alone and its conclusion is a support-fiber bound,
not a distinct-slope bound.
We therefore replace that conditional equality by two unconditional
statements on their natural scales: an endpoint-exact linear-budget formula
and a one-sided positive-relative-radius exponential-budget safe-frontier
certificate.

Polynomial-generator preservation and syndrome-space results concern other
sampler or code classes \cite{BordageChiesaGuanManzur25,YuanZhu26}, while
WHIR supplies a protocol-level use of MCA \cite{WHIR}; none determines the
arbitrary-domain finite MDS thresholds studied here.

\begin{table}[tbp]
\centering
\caption{Proof status and the role of this paper.}
\label{tab:novelty-status}
\small
\begin{tabularx}{\textwidth}{@{}P{0.28\textwidth}P{0.20\textwidth}Y@{}}
\toprule
Statement & Status & Attribution and role here \\
\midrule
Positive-relative-radius exponential-budget safe-frontier certificate &
derived here & Jo's shortening transfer and the BCHKS Johnson estimate give
the exponent \(\Psi_\rho\); this is one-sided and not a matching
near-capacity threshold. \\
All-test-size MDS circuit envelope & imported and reproduced &
Jo's rejecting-test proof is reproduced here; at \(b=k+1\) it has the
augmented-word matroid interpretation recorded in
\cref{rem:paving-saving}. \\
RS staircase below half the minimum distance & prior upper input; exact
bookkeeping here & BCHKS plus the CA--MCA conversion supply the asymptotic
phenomenon; this paper retains restricted challenges, integer floors, and
open endpoints. \\
Linear MCA bound below Johnson & prior & BCHKS Theorem~4.6 supplies the
Johnson-range input used in the exponential-budget certificate; Hab\"ock
outlines the earlier quadratic route. \\
Primitive inverse from a Sidon payment at accessible moment order &
proved implication & Given \(\log L_N/q_N=o(N)\), BSG and Boolean-cube
difference growth eliminate the high-energy Boolean branch; proving the
image-normalized Sidon payment and the residual ray compiler remain separate
inputs. \\
Parameter-retained BCIKS--BCHKS calculation & conditional & The four
candidate KoalaBear rows depend on the factor-lifting assumption isolated in
\cref{app:retained-interpolation}; no unconditional theorem uses it. \\
Unrestricted smooth/circle near-capacity frontier & open & The repository's
profile, analytic-payment, ray, add-back, and reserve interfaces are not
assumed in any theorem. \\
\bottomrule
\end{tabularx}
\end{table}

Here \emph{smooth} means possession of specified complete equal-fiber
folding maps, not differentiability.  A \emph{circle presentation} means one
of the circle evaluation codes defined in \cref{sec:smooth-circle-domains}
and transferred only through a proved agreement-preserving monomial
equivalence.  All linear-algebraic, MDS, paving, syndrome, and ordinary
Reed--Solomon conclusions are characteristic-independent unless stated
otherwise.  Circle uniformizations using \(x^2+y^2=1\), stereographic
projection, or a norm-one torus require odd characteristic.  In
characteristic two we claim only the explicitly stated Reed--Solomon results
on additive domains in extension fields; ``binary'' describes
\(\F_2\)-linear domain structure inside \(\F_{2^m}\), not an \(n\)-point
domain contained in the two-element field.

\paragraph{Scope of the unconditional claims.}
The exact linear-budget theorem, Jo's MDS circuit envelope, the one-sided
shortening certificate, the all-radius completion, and every explicitly
verified finite row are unconditional subject only to the displayed BCHKS
and Jo inputs.  The parameter-retained appendix is
conditional on its displayed factor-lifting assumption.  The primitive
inverse theorem is unconditional as an implication from its stated Sidon
payment at an accessible moment order, but it does not establish that
premise on every smooth or circle leaf.  No theorem in this paper asserts the unrestricted
constant-rate, positive-relative-radius smooth/circle profile frontier.

\paragraph{AI-use disclosure.}
OpenAI ChatGPT/Codex, a GPT-5-based system, was used substantively to draft
and restructure prose, propose and audit proof steps, check calculations,
and edit \LaTeX\ throughout this manuscript.  The human author supplied the
research program, source manuscripts, repository materials, mathematical
choices, and corrections, and retains sole responsibility for correctness
and originality.  A more detailed use statement accompanies the source
release as \path{AI_USE_v9.2.md}.

\paragraph{Principal original bound.}
Put \(J_\rho=1-\sqrt\rho\).  For
\(J_\rho\le\delta\le1-\rho\), set \(\alpha=1-\delta\) and
\[
 \tau_\rho(\delta)=\frac{\rho-\alpha^2}{1+\rho-2\alpha},\qquad
 \Psi_\rho(\delta)=H_2(\tau_\rho(\delta))
 -\alpha H_2\!\left(\frac{\tau_\rho(\delta)}{\alpha}\right), \tag{SF1}
 \label{eq:intro-Psi}
\]
and extend \(\Psi_\rho\) by zero on \([0,J_\rho]\).  For comparison, put
\[
 \Phi_\rho(\delta)=H_2(\rho)
 -(1-\delta)H_2\!\left(\frac{\rho}{1-\delta}\right),
 \qquad 0\le\delta\le1-\rho.                       \tag{SF2}
 \label{eq:intro-Phi}
\]

\begin{theorem}[Positive-relative-radius shortening exponent]
\label{thm:intro-shortening-exponent}
Let \(C_n=\RS_{\F_n}(D_n,k_n)\) with \(k_n/n\to\rho\in(0,1)\), and let
\(\varnothing\ne\Gamma_n\subseteq\F_n\).  Fix
\(J_\rho<\delta<1-\rho\) and put
\(a_n=n-\floor{\delta n}\).  Then
\[
 \limsup_{n\to\infty}\frac1n\log_2
 \bigl(1+B_{C_n,\Gamma_n}^{\rm MCA}(a_n)\bigr)
 \le\Psi_\rho(\delta).                             \tag{SF3}
 \label{eq:shortening-exponent-bound}
\]
Moreover, \(0<\tau_\rho(\delta)<\rho\) and
\[
 \Psi_\rho(\delta)<\Phi_\rho(\delta)              \tag{SF4}
 \label{eq:shortening-beats-circuit}
\]
strictly between Johnson and capacity.  The exponent extends continuously
with \(\Psi_\rho(J_\rho)=0\) and
\(\Psi_\rho(1-\rho)=H_2(\rho)\).
\end{theorem}

The proof is given in \cref{sec:agreement-shortening}.  The optimization
leading to \(\Psi_\rho\) was carried out after the author learned of Jo's
agreement-set shortening theorem \cite{Jo26}: Jo proves the transfer and the
fixed-integer-step consequence, while the positive-relative-radius entropy
optimization and its Prize-budget formulation are derived here.

\begin{theorem}[Imported circuit-incidence MDS envelope, Jo
  {\cite[Thm.~4.2]{Jo26}}]\label{thm:intro-paving}
Let \(C\subseteq\F^n\) be a linear \([n,k]\) MDS code,
let \(\varnothing\ne\Gamma\subseteq\F\), and let \(k+1\le a\le n\).
Then
\[
 B_{C,\Gamma}^{\rm CA}(a)\le B_{C,\Gamma}^{\rm MCA}(a)
 \le U_{n,k,\Gamma}(a)
 \defeq\min\left\{\abs\Gamma,
 \min_{k+1\le b\le a}
 \floor{\frac{\binom nb}{\binom{a-1}{b-1}}}\right\}.       \tag{1.3}
 \label{eq:intro-paving}
\]
The full-field estimate is due to Jo; the restricted-challenge form follows
by intersecting the bad-slope set with \(\Gamma\).  Its endpoint test sizes
\(b=k+1\) and \(b=a\) recover, respectively, the paving and support-atlas
bounds used below.
The bound is invariant under coordinate permutations and nonzero diagonal
scalings, and hence applies without loss to generalized Reed--Solomon codes
and to the circle presentations proved monomially equivalent to them.
\end{theorem}

\begin{theorem}[Imported BCHKS Johnson-range estimate, \(M=1\)]
\label{thm:bchks-johnson-general}
Let \(N,d,E\) be integers with \(0<d<N\), \(0<E<N\), and
\((N-E)^2>Nd\).  Put
\[
 \rho_0=\frac dN,\qquad \gamma_0=\frac EN,\qquad
 \eta_0=1-\sqrt{\rho_0}-\gamma_0,
\]
\[
 m_B=\max\left\{\left\lceil
       \frac{\sqrt{\rho_0}}{\eta_0}\right\rceil,3\right\},
 \qquad x_B=m_B+\frac12,
\]
and
\[
 \Jcal(N,d,E)=
 \left(\frac{2x_B^5+3x_B\gamma_0\rho_0}
 {3\rho_0^{3/2}}\right)N+\frac{x_B}{\sqrt{\rho_0}}. \tag{B1}
 \label{eq:bchks-johnson-general}
\]
For every Reed--Solomon code \(C=\RS_\F(L,d+1)\) of length \(N\) and
every nonempty \(\Gamma\subseteq\F\),
\[
 B_{C,\Gamma}^{\rm MCA}(N-E)
 \le B_C^{\rm MCA}(N-E)
 \le \Jcal(N,d,E).                                 \tag{B2}
 \label{eq:bchks-johnson-bound}
\]
\end{theorem}

\begin{proof}
This is the affine-line \(M=1\) specialization of
BCHKS \cite[Thm.~4.6]{BCHKS25}, in the explicit normalization of
Jo \cite[Thm.~2.3 and Eqs.~(6)--(9)]{Jo26}.  There the polynomial degree is
at most \(d\), hence the code dimension is \(d+1\), and the strict
Johnson condition is precisely \((N-E)^2>Nd\).  Restricting the slope set
from \(\F\) to \(\Gamma\) can only decrease the numerator.
\end{proof}

\begin{theorem}[Official endpoint and all-radius saturation]
\label{thm:official-saturation}
Let \(C\subseteq\F^n\) be a linear \([n,k]\) MDS code with \(1\le k<n\), let
\(\varnothing\ne\Gamma\subseteq\F\), let \(0\le\eps\le1\), and extend the
threshold notation to \(0\le a\le n\).  Then
\[
 B_{C,\Gamma}^{\rm MCA}(a)=B_{C,\Gamma}^{\rm MCA}(k+1)
 \qquad(0\le a\le k+1).                           \tag{OFF1}
 \label{eq:subcapacity-constant}
\]
Put \(B^*=\floor{\eps\abs\Gamma}\), and, when the defining set is
nonempty, put
\[
 a^*=\min\{a\in\{k+1,\ldots,n\}:
               B_{C,\Gamma}^{\rm MCA}(a)\le B^*\}.
\]
If
\(B_{C,\Gamma}^{\rm MCA}(k+1)\le B^*\), the complete official safe set is
\([0,1]\) and
\[
 \delta_{C,\Gamma,\mathrm{off,grid}}^*(\eps)
 =\delta_{C,\Gamma,\mathrm{off,sup}}^*(\eps)=1.   \tag{OFF2}
 \label{eq:official-saturation}
\]
In particular, the sufficient condition
\[
 B^*\ge U_{n,k,\Gamma}(k+1)
     =\min\left\{\abs\Gamma,\binom n{k+1}\right\} \tag{OFF3}
 \label{eq:official-saturation-budget}
\]
implies \eqref{eq:official-saturation}.  For the full-field target
\(2^{-128}\), it is enough that
\(\floor{2^{-128}\abs\F}\ge\binom n{k+1}\), equivalently
\(\abs\F\ge 2^{128}\binom n{k+1}\).

If instead \(B_{C,\Gamma}^{\rm MCA}(k+1)>B^*\) and \(a^*\) exists, then
necessarily \(a^*>k+1\), and the official and subcapacity transitions agree:
the largest safe grid radius is \(1-a^*/n\), the complete real safe set is
\([0,(n-a^*+1)/n)\), and its right endpoint is unsafe.
\end{theorem}

\begin{proof}
For every support \(S\) with \(\abs S\le k\), extend \(S\) to a
\(k\)-element information set \(T\).  The puncturing map
\(C\to\F^T\) is an isomorphism, so its further projection
\(C\to\F^S\) is surjective.  Hence both
received words are explained on \(S\), so no support-wise nontrivial witness
has size at most \(k\).  This proves \eqref{eq:subcapacity-constant}.  The
first conclusion follows immediately.  At \(a=k+1\),
\cref{thm:intro-paving} gives
\(B_{C,\Gamma}^{\rm MCA}(k+1)\le
U_{n,k,\Gamma}(k+1)\), proving the sufficient condition.  In the
nonsaturated case, monotonicity and the identity
\(a(\delta)=n-\floor{\delta n}\) give the stated grid point and half-open
real interval.
\end{proof}

\begin{corollary}[Four exact-plateau saturated rows]
\label{cor:intro-four-saturated-rows}
Let
\[
 p=(2^{128}-255)2^{128}+1
\]
and let \(D\le\F_p^\times\) be the subgroup of order \(128\).  For
\[
 C_\rho=\RS_{\F_p}(D,128\rho),
 \qquad \rho\in\{1/2,1/4,1/8,1/16\},
\]
the complete official safe set at target \(2^{-128}\) is \([0,1]\), so
both official grid and supremum endpoints equal \(1\) for all four rates.
Moreover, with \(M_\rho=\binom{128}{128\rho+1}\),
\[
 B_{C_\rho}^{\rm MCA}(a)=M_\rho
 \qquad(0\le a\le128\rho+1).
\]
The field has fewer than \(2^{256}\) elements and the evaluation domain has
power-of-two order.
\end{corollary}

\begin{proof}
The Proth certificate in \cref{prop:paving-finite-rows} proves that \(p\)
is prime; since \(128\mid p-1\), the subgroup exists.  The target budget is
\(2^{128}-255\), while
\[
\begin{array}{c|r}
k&\binom{128}{k+1}\\ \hline
64&23582666872052266206656578733667004800\\
32&4299074680733907393985381161600\\
16&614965786737727286400\\
8 &19062702032000
\end{array}
\]
and every entry is smaller than that budget.  Apply
\cref{thm:official-saturation}.  The largest \(M_\rho\) is less than
\(2^{125}\), while \(p>2^{255}\), so
\(p>\binom{M_\rho}{2}\) in every row.  The exact plateau follows from
\cref{thm:jo-rs-endpoint}.
\end{proof}

\begin{remark}[The field-size quantifier]
\label{rem:official-field-size-quantifier}
For fixed \(n\) and \(k\), arbitrary scalar-field enlargement eventually
forces \eqref{eq:official-saturation-budget}; the resulting endpoint \(1\) is a
budget-saturation statement, not a near-capacity decoding theorem.  If the
phrase \emph{sufficiently large} is intended to retain a prescribed field
ceiling or a growing-length relation between \(q\) and \(n\), the supplied row
must instead be evaluated by the finite or asymptotic certificates below.  In
particular, selected saturated rows do not solve the uniform given-\(C\)
problem.
\end{remark}

The endpoint test sizes in \eqref{eq:intro-paving} recover the support-atlas
and paving bounds, while intermediate sizes can improve both at finite
length.  For Reed--Solomon rows and their monomially equivalent GRS or
circle presentations---indeed, after \cref{prop:universal-tangent-floor},
for every linear MDS code---put
\(L_{n,\Gamma}(a)=\min\{\abs\Gamma,n-a+1\}\), the tangent construction and
\cref{thm:intro-paving} give the finite security interval
\[
 \log_2\frac{\abs\Gamma}{U_{n,k,\Gamma}(a)}
 \le -\log_2\frac{B_{C,\Gamma}^{\rm MCA}(a)}{\abs\Gamma}
 \le\log_2\frac{\abs\Gamma}{L_{n,\Gamma}(a)}.       \tag{1.4}
 \label{eq:intro-security-interval}
\]
Thus replacing one rigorous upper numerator \(U_0\) by a smaller \(U_1\)
earns exactly \(\log_2(U_0/U_1)\) proved bits.  The remaining interval has
width at most \(\log_2(U_1/L_{n,\Gamma}(a))\); this is the auditable
pay-per-bit quantity, independent of any eventual prize-allocation rule.
For rigorous upper numerators \(U_0\ge U_1>0\) and attack lower numerators
\(0<L_0\le L_1\), the exact uncertainty contraction is
\[
 \log_2\frac{U_0}{L_0}-\log_2\frac{U_1}{L_1}
 =\log_2\frac{U_0}{U_1}+\log_2\frac{L_1}{L_0}.     \tag{SG}
\]
This identity keeps proof-side and falsification-side gains separate and
prevents the same closed bit from being counted twice.

\begin{theorem}[Exact linear-budget threshold]
\label{thm:intro-unconditional-asymptotic-deep}
Let \(C_n=\RS_{\F_n}(D_n,k_n)\), let \(k_n/n\to\rho\in(0,1)\), let
\(\varnothing\ne\Gamma_n\subseteq\F_n\), let \(0\le\eps_n\le1\), and put
\(B_n^*=\floor{\eps_n\abs{\Gamma_n}}\).
If \(B_n^*=0\), no agreement is safe; if
\(B_n^*=\abs{\Gamma_n}\), then \(a_n^*=k_n+1\),
\(\delta_{n,\mathrm{sub,grid}}^*=1-(k_n+1)/n\),
\(\delta_{n,\mathrm{sub,sup}}^*=1-k_n/n\), and
\(\delta_{C_n,\Gamma_n,\mathrm{off,grid}}^*
=\delta_{C_n,\Gamma_n,\mathrm{off,sup}}^*=1\).
Otherwise suppose that, for some \(\eta>0\), eventually
\[
 1\le B_n^*<\abs{\Gamma_n},\qquad
 B_n^*-1\le\left(\frac{1-\rho}{2}-\eta\right)n.    \tag{1.5}
 \label{eq:intro-linear-budget-range}
\]
Then
\[
 a_n^*=n-B_n^*+1,\qquad
 \delta_{n,\mathrm{sub,grid}}^*=\frac{B_n^*-1}{n},\tag{1.6}
 \label{eq:intro-linear-budget-formula}
\]
and, under the continuous convention \(r(\delta)=\floor{\delta n}\), the
complete official safe set is \([0,B_n^*/n)\); its endpoint is unsafe.
Thus the official and subcapacity grid values agree and the official
supremum is \(B_n^*/n\).
Under these otherwise hypotheses, if \(B_n^*/n\to b<(1-\rho)/2\), then
\(\delta_{n,\mathrm{sub,grid}}^*\to b\) and
\(\delta_{C_n,\Gamma_n,\mathrm{off,sup}}^*(\eps_n)\to b\).
\end{theorem}

\begin{proof}
The zero-budget case follows from the tangent numerator one at agreement
\(n\); the full-budget case follows from
\cref{thm:official-saturation}.  In the remaining case,
the uniform below-half staircase of \cref{cor:positive-radius-density},
at radius \(B_n^*-1\), gives numerator \(B_n^*\).  At the next radius
\(B_n^*\), the tangent construction gives
\(\min\{\abs{\Gamma_n},B_n^*+1\}=B_n^*+1\).
The range assumption also implies \(B_n^*\le n-k_n-1\) eventually, so both
radii belong to the admissible grid.  Monotonicity proves the first-safe
agreement and \cref{lem:appendix-endpoint} gives the half-open real interval.
\end{proof}

For the prize-facing fixed-bit scaling there is no entropy inversion to
perform.  If the challenge is the whole field, the target is
\(2^{-\lambda}\), \(q_n/n\to\kappa\), and
\(b=\kappa2^{-\lambda}<(1-\rho)/2\), then, whenever the endpoint cases are
excluded,
\[
 a_n^*=n-\floor{q_n2^{-\lambda}}+1,\qquad
 \delta_{n,\mathrm{sub,grid}}^*\longrightarrow b, \tag{1.7}
 \label{eq:intro-prize-linear-formula}
\]
This is an exact asymptotic formula for \(q_n=\Theta(n)\).  If instead
\(q_n=2^{o(n)}\) and the fixed-bit budget is \(o(n)\), the exact grid
threshold still has the form \eqref{eq:intro-linear-budget-formula} and
therefore tends to radius zero.

The shortening exponent of \cref{thm:intro-shortening-exponent} converts
directly into a target-budget certificate.  For \(\theta>0\), define
\[
 D_{\rho,\mathrm{short}}^{\rm sub}(\theta)
 =\sup\{0\le\delta<1-\rho:\Psi_\rho(\delta)<\theta\},
 \quad
 D_{\rho,\mathrm{short}}^{\rm off}(\theta)=
 \begin{cases}
 D_{\rho,\mathrm{short}}^{\rm sub}(\theta),&
                    \theta\le H_2(\rho),\\
 1,&\theta>H_2(\rho).
 \end{cases}                                      \tag{1.10}
 \label{eq:intro-D-short}
\]
In particular, every \(\theta>0\) moves the certified subcapacity radius
strictly beyond Johnson.

\begin{theorem}[Positive-relative-radius exponential-budget safe-frontier certificate]
\label{thm:intro-exponential-frontier}
Let \(C_n=\RS_{\F_n}(D_n,k_n)\), let \(k_n/n\to\rho\in(0,1)\),
let \(\varnothing\ne\Gamma_n\subseteq\F_n\), let \(0\le\eps_n\le1\), and put
\(B_n^*=\floor{\eps_n\abs{\Gamma_n}}\).  Suppose
\[
 \frac1n\log_2(1+B_n^*)\longrightarrow\theta>0.    \tag{1.11}
 \label{eq:intro-budget-exponent}
\]
Then \(a_n^*\) is defined for all sufficiently large \(n\), and
\[
 \liminf_{n\to\infty}\delta_{n,\mathrm{sub,grid}}^*
 \ge D_{\rho,\mathrm{short}}^{\rm sub}(\theta),\qquad
 \liminf_{n\to\infty}\delta_{C_n,\Gamma_n,\mathrm{off,grid}}^*(\eps_n)
 \ge D_{\rho,\mathrm{short}}^{\rm off}(\theta).  \tag{1.12}
 \label{eq:intro-certified-frontier}
\]
If \(\theta\ge H_2(\rho)\), then
\(\delta_{n,\mathrm{sub,grid}}^*\to1-\rho\).  If the inequality is
strict, then \(a_n^*=k_n+1\) eventually and the official safe set is
\([0,1]\), so
\(\delta_{C_n,\Gamma_n,\mathrm{off,grid}}^*(\eps_n)
=\delta_{C_n,\Gamma_n,\mathrm{off,sup}}^*(\eps_n)=1\).
At equality the subcapacity limiting radius is still capacity, but finite
constants decide whether the official endpoint saturates at one.  For
\(0<\theta<H_2(\rho)\),
\eqref{eq:intro-certified-frontier} is a one-sided certificate, not an
exact threshold.
\end{theorem}

\begin{proof}
The positive budget exponent gives \(B_n^*\ge1\) eventually, while
\cref{thm:intro-paving} gives
\(B_{C_n,\Gamma_n}^{\rm MCA}(n)\le1\); hence \(a_n^*\) exists.
For every fixed \(0<\delta<J_\rho\), apply
\cref{thm:bchks-johnson-general} with
\[
 N=n,\qquad d=k_n-1,\qquad E=\floor{\delta n}.
\]
Here \(d/N\to\rho\) and
\(1-\sqrt{d/N}-E/N\to J_\rho-\delta>0\), so the integer \(m_B\) in
\eqref{eq:bchks-johnson-general} is bounded and
\(\Jcal(N,d,E)=O_{\rho,\delta}(n)\).  At \(\delta=0\), the same
subexponential conclusion
follows already from \cref{thm:intro-paving}.  Restricting from \(\F_n\)
to \(\Gamma_n\) can only decrease the numerator, and every such numerator
is eventually below a budget of positive exponent.  Beyond Johnson,
\cref{thm:intro-shortening-exponent} shows that every fixed radius with
\(\Psi_\rho(\delta)<\theta\) is safe.  At Johnson itself, choose a nearby
larger radius with the same strict inequality and use monotonicity.
Taking suprema gives \eqref{eq:intro-certified-frontier}; the official
bound follows because it contains the subcapacity safe grid, with the strict
saturation branch supplied by \cref{thm:official-saturation}.
Every subcapacity grid radius is at most \(1-(k_n+1)/n\), proving the
matching subcapacity limit when \(\theta\ge H_2(\rho)\).  If
\(\theta>H_2(\rho)\), the envelope at \(a=k_n+1\) has exponent at most
\(H_2(\rho)\), so the first agreement is safe; then
\cref{thm:official-saturation} gives the official endpoint one.
\end{proof}

\begin{remark}[Imported Johnson-range input]
\label{rem:bchks-theorem46-provenance}
The load-bearing Johnson-range input is BCHKS Theorem~4.6 in the
standalone \((N,d,E)\) normalization of
\cref{thm:bchks-johnson-general}.  The source states the displayed theorem
and sketches the list-correlated-agreement generalization.  Thus the
below-Johnson part of
\cref{thm:intro-exponential-frontier,cor:intro-security-exponents} depend on
that published statement.  Jo's shortening transfer and rejecting-test
proof of the all-test-size MDS envelope are reproduced in
\cref{thm:jo-shortening-transfer,lem:rejecting-test-abundance}.
\end{remark}

\begin{corollary}[Security and soundness exponents]
\label{cor:intro-security-exponents}
Let \(C_n=\RS_{\F_n}(D_n,k_n)\) with \(k_n/n\to\rho\in(0,1)\), put
\(q_n=\abs{\F_n}\), assume the full-field challenge, and suppose
\(\log_2 q_n/n\to\phi\).  For a fixed radius \(\delta\), let
\(s_n(\delta)=-\log_2 e_{\rm MCA}(C_n,\delta)\).  Then
\[
 \frac{s_n(\delta)}n\longrightarrow\phi
 \quad(0\le\delta<J_\rho),\qquad
 \liminf_{n\to\infty}\frac{s_n(\delta)}n
 \ge\bigl(\phi-\Psi_\rho(\delta)\bigr)_+
 \quad(J_\rho\le\delta<1-\rho).                   \tag{1.15}
 \label{eq:intro-security-exponents}
\]
The tangent attack gives \(\limsup s_n(\delta)/n\le\phi\), so the
remaining soundness-gap rate beyond Johnson is at most
\(\min\{\phi,\Psi_\rho(\delta)\}\).  Relative to the smallest-test MDS
exponent, shortening closes
\[
 \min\{\phi,\Phi_\rho(\delta)\}
 -\min\{\phi,\Psi_\rho(\delta)\}
\]
bits per coordinate; this equals
\(\Phi_\rho(\delta)-\Psi_\rho(\delta)\) only when
\(\phi\ge\Phi_\rho(\delta)\).  For a target
\(2^{-\lambda_n}\) with \(\lambda_n\ge0\),
\(\lambda_n/n\to\lambda\), and
\(\phi-\lambda>0\), the certified limiting frontier is
\(D_{\rho,\mathrm{short}}^{\rm sub}(\phi-\lambda)\): every fixed radius
strictly below it is
eventually safe.  If \(\phi-\lambda>H_2(\rho)\), the official safe set is
\([0,1]\) eventually.  When \(\phi>0\), a fixed 128-bit target has
\(\lambda=0\); for \(q_n=\Theta(n)\), the exact linear-budget formula
\eqref{eq:intro-prize-linear-formula} is the relevant statement instead.
\end{corollary}

\begin{proof}
Below Johnson, the BCHKS linear upper numerator and the tangent polynomial
lower numerator have zero base-two exponent, which gives the limit.  For
\(\delta>J_\rho\), combine the field cap with
\cref{thm:intro-shortening-exponent}.  At \(\delta=J_\rho\), monotonicity
gives, for every fixed \(\delta'>J_\rho\) and all large \(n\),
\[
 B_{C_n}^{\rm MCA}(n-\floor{J_\rho n})
 \le B_{C_n}^{\rm MCA}(n-\floor{\delta'n}).
\]
Apply that theorem and let
\(\delta'\downarrow J_\rho\).  The target statement follows because
\(\log_2(1+\floor{q_n2^{-\lambda_n}})/n\to\phi-\lambda\).
\end{proof}

The universal MDS envelope gives 128-bit multiplicative and circle rows
beyond Johnson at length \(512\); see
\cref{prop:paving-finite-rows,prop:paving-circle-certificate}.
The exact Proth rows of
\cref{sec:prize-rows} are instances of the below-half formula and are
included for endpoint and certificate verification, not as a new
literature consequence.  Fixed-length congruence classes give positive
prime density, but that density tends to zero along growing dyadic
prime-field lengths.  Circle claims below always pass through an explicit
agreement-preserving diagonal equivalence; no statement is made for an
unidentified bivariate presentation.
\section{MCA witnesses and exact support reduction}\label{sec:witnesses}

Throughout this section \(\F\) is a finite field, \(D\subseteq\F\) has
size \(n\), and \(C=\RS_\F(D,k)\).  Write
\(\F[X]_{<k}=\{f\in\F[X]:\deg f<k\}\), and let
\(\operatorname{ev}_S:\F[X]_{<k}\to\F^S\) be restriction of the
evaluation map.

\begin{definition}[Explanation and support-wise nontriviality]
\label{def:explanation}
For \(u\in\F^D\) and \(S\subseteq D\), an \emph{explanation of \(u\) on
\(S\)} is a polynomial \(f\in\F[X]_{<k}\) with
\(\operatorname{ev}_S(f)=u|_S\).  We write
\(\operatorname{Expl}_C(u;S)\) when such an \(f\) exists.  A pair
\(r=(r_0,r_1)\) is \emph{simultaneously explained on \(S\)} when
\(\operatorname{Expl}_C(r_0;S)\) and
\(\operatorname{Expl}_C(r_1;S)\) both hold, with no requirement that the
two explaining polynomials coincide.  It is \emph{support-wise
nontrivial on \(S\)} when it is not simultaneously explained there.
We denote this last condition by \(\operatorname{NT}_C(r;S)\).
\end{definition}

Thus support-wise nontriviality is a condition on the same support that
explains the affine combination.  It does not mean that both coordinates
are individually unexplained; failure of either explanation is enough.

\begin{definition}[MCA-bad slope]\label{def:mca-bad}
Let \(r=(r_0,r_1)\in(\F^D)^2\) and let \(a\) be an agreement threshold.
A finite slope \(\gamma\in\F\) is support-wise MCA-bad at agreement \(a\)
if there are \(S\subseteq D\), \(\abs S\ge a\), and
\(h\in\F[X]_{<k}\) such that
\(\operatorname{ev}_S(h)=(r_0+\gamma r_1)|_S\) and
\(\operatorname{NT}_C(r;S)\) holds.  Define
\[
        B_C^{\rm MCA}(a)
        =\max_{r\in(\F^D)^2}\abs{\{\gamma\in\F:
          \gamma\text{ is MCA-bad for }r\text{ at }a\}}.
\]
\end{definition}

\begin{lemma}[Support-wise and maximal-domain formulations]
\label{lem:mca-definition-bridge}
For a received pair \(r=(r_0,r_1)\) and an integer \(a\), let
\[
 \mathcal E_{C,a}(r)=\{S\subseteq D:\abs S\ge a,
     \ r_0|_S\in C|_S,\ r_1|_S\in C|_S\},
\]
and let \(\mathcal M_{C,a}(r)\) be its inclusion-maximal members.  Put
\(\operatorname{Agr}(u,c)=\{x\in D:u(x)=c(x)\}\).  A slope \(\gamma\)
is MCA-bad for \(r\) at agreement \(a\) if and only if there is \(c\in C\)
such that, for \(A=\operatorname{Agr}(r_0+\gamma r_1,c)\),
\[
 \abs A\ge a\quad\text{and}\quad
 A\nsubseteq M\quad\text{for every }M\in\mathcal M_{C,a}(r). \tag{MDB}
 \label{eq:maximal-domain-bridge}
\]
If \(a>n-d_{\min}(C)\), the last condition is equivalently
\(A\notin\mathcal M_{C,a}(r)\).  Consequently, for \(0<\delta<1\) and
\(a(\delta)=\ceil{(1-\delta)n}=n-\floor{\delta n}\),
\[
 \varepsilon_{\rm mca}(C,\delta)=\frac1{\abs\F}
 \max_r\#\{\gamma\in\F:\gamma\text{ satisfies \eqref{eq:maximal-domain-bridge}}\},
\]
which is the quantity in \cite[Definition~4.3]{ABF26}; our values at
\(\delta=0,1\) are its natural endpoint extensions.
\end{lemma}

\begin{proof}
For every \(S\subseteq D\) with \(\abs S\ge a\), simultaneous
explanation on \(S\) is equivalent, by finiteness and restriction, to
containment in some member of \(\mathcal M_{C,a}(r)\).  If \(S\) witnesses
that \(\gamma\) is bad and \(c\) explains the line point on \(S\), then
\(S\subseteq A=\operatorname{Agr}(r_0+\gamma r_1,c)\).  Were \(A\)
contained in a maximal common domain, restriction would simultaneously
explain the pair on \(S\), a contradiction.  Conversely,
\eqref{eq:maximal-domain-bridge} says that the pair is not simultaneously
explained on \(A\), so \(S=A\) is a support-wise witness.

Assume now \(a>n-d_{\min}(C)\) and that the pair is simultaneously
explained on \(A\).  The two explaining codewords and \(c\) are unique on
\(A\).  Extending \(A\) to a maximal common domain therefore uses the same
two codewords, whose affine combination is \(c\); the extension is contained
in \(A\) and hence equals it.  Thus \(A\) itself is maximal.  The probability
identity follows because the official sampler chooses \(\gamma\) uniformly
from \(\F\).
\end{proof}

The official MCA maximum mixes two logically different sources: received
pairs that are globally far from every codeword pair, and pairs that are
close to one codeword pair but fail on a slope-dependent support.  The
following definitions and theorem separate them exactly.

\begin{definition}[CA numerator and sparse mutual numerator]
\label{def:ca-sparse-numerators}
Assume \(0\le a\le n\) and put \(t=n-a\).  A received pair
\((r_0,r_1)\) is \emph{column-far at
agreement \(a\)} if it is not simultaneously explained on any support of
size at least \(a\).  Let \(B_C^{\rm CA}(a)\) be the maximum, over
column-far pairs, of the number of slopes \(\gamma\) for which
\(r_0+\gamma r_1\) is explained on some support of size at least \(a\).
If there is no column-far pair, this maximum is defined to be zero.
Define the \emph{sparse mutual numerator}
\[
 \sigma_C(a)=\max\left\{
 \#\{\gamma:\gamma\text{ is MCA-bad for }(e_0,e_1)\}:
 \abs{\supp(e_0)\cup\supp(e_1)}\le t\right\}.
 \tag{SP1}\label{eq:sparse-mutual-numerator}
\]
Thus \(B_C^{\rm CA}(a)/\abs\F\) is the ordinary affine-line correlated
agreement error on the integer grid, while \(\sigma_C(a)\) records the
purely mutual obstruction after translation by a nearby codeword pair.
\end{definition}

\begin{theorem}[Exact sparsification of the mutual layer]
\label{thm:exact-sparsification}
For every linear code \(C\subseteq\F^D\) and every agreement threshold
\(0\le a\le n\),
\[
 B_C^{\rm MCA}(a)=
 \max\{B_C^{\rm CA}(a),\sigma_C(a)\}.              \tag{SP2}
\label{eq:exact-sparsification}
\]
Equivalently, after division by \(\abs\F\), the official MCA error is the
maximum of the CA error and the sparse mutual error.
\end{theorem}

\begin{proof}
Both quantities on the right are obtained by maximizing the MCA-bad set
over restricted classes of pairs, so the left side is at least their
maximum.  Conversely fix a received pair.  If it is column-far, every
support that explains a line point necessarily fails to explain the pair;
its MCA-bad set is therefore its CA-bad set and has size at most
\(B_C^{\rm CA}(a)\).

Otherwise choose codewords \(c_0,c_1\in C\) and a common support of size
at least \(a\), and put \(e_i=r_i-c_i\).  The union of the supports of
\(e_0,e_1\) has size at most \(n-a=t\).  At each slope, translation sends
an explaining codeword \(h\) to
\(h-c_0-\gamma c_1\), and is a bijection between witness supports for
\((r_0,r_1)\) and for \((e_0,e_1)\).  It also preserves simultaneous
explanation on that support.  Hence the two pairs have exactly the same
MCA-bad slopes, and their number is at most \(\sigma_C(a)\).
\end{proof}

For a nonempty challenge set \(\Gamma\subseteq\F\), write
\(B_{C,\Gamma}^{\rm CA}(a)\) and \(\sigma_{C,\Gamma}(a)\) for the same
maxima with the bad slopes restricted to \(\Gamma\).  The proof just given
is unchanged after intersecting every bad-slope set with \(\Gamma\), and
therefore gives the exact challenge-restricted identity
\[
 B_{C,\Gamma}^{\rm MCA}(a)
 =\max\{B_{C,\Gamma}^{\rm CA}(a),\sigma_{C,\Gamma}(a)\}.
 \tag{SP3}\label{eq:challenge-sparsification}
\]

\begin{theorem}[Exact mutual layer below half the distance]
\label{thm:exact-half-distance-sparse}
Let \(C\subseteq\F^D\) be a linear code of minimum distance
\(d_{\min}\), let \(\varnothing\ne\Gamma\subseteq\F\), and let
\(0\le r\le n\) satisfy \(2r\le d_{\min}-1\).  Then
\[
 \sigma_{C,\Gamma}(n-r)=\min\{\abs\Gamma,r\},
 \qquad
 B_{C,\Gamma}^{\rm MCA}(n-r)
 =\max\left\{B_{C,\Gamma}^{\rm CA}(n-r),
              \min\{\abs\Gamma,r\}\right\}.       \tag{HD1}
\label{eq:exact-half-distance-sparse}
\]
For a Reed--Solomon row the hypothesis is \(2r\le n-k\).
\end{theorem}

\begin{proof}
Fix a sparse pair \((e_0,e_1)\) whose support union \(E\) has size at
most \(r\), and let \(\gamma\) be bad on a support \(S\) of size at least
\(n-r\), explained there by a codeword \(c\in C\).  On \(S\setminus E\)
the sparse line point is zero, so
\(\supp(c)\subseteq E\cup(D\setminus S)\) has size at most \(2r<d_{\min}\).
Thus \(c=0\).  If \(S\cap E\) contained no nonzero pair coordinate, the
zero codeword pair would simultaneously explain \((e_0,e_1)\) on \(S\),
contrary to badness.  Hence some \(x\in S\cap E\) satisfies
\[
 (e_0(x),e_1(x))\ne(0,0),\qquad e_0(x)+\gamma e_1(x)=0.
\]
Necessarily \(e_1(x)\ne0\), and this coordinate determines the single
slope \(-e_0(x)/e_1(x)\).  Therefore every sparse pair has at most
\(\min\{\abs\Gamma,\abs E\}\le\min\{\abs\Gamma,r\}\) bad slopes.

For the reverse bound put \(t=\min\{\abs\Gamma,r\}\), choose distinct
\(x_1,\ldots,x_t\in D\) and distinct
\(\gamma_1,\ldots,\gamma_t\in\Gamma\), and set
\(e_1(x_i)=1\), \(e_0(x_i)=-\gamma_i\), with both words zero elsewhere.
For \(S_i=D\setminus\{x_j:j\ne i\}\), the line point of slope
\(\gamma_i\) is zero on \(S_i\), and
\(\abs{S_i}=n-t+1\ge n-r\).  The restriction of \(e_1\) to \(S_i\)
cannot be explained by a codeword: such a codeword would be nonzero at
\(x_i\) and zero on \(n-t\) coordinates, hence would have weight at most
\(t<d_{\min}\).  Thus every \(\gamma_i\) is MCA-bad and equality holds.
The second formula is \eqref{eq:challenge-sparsification}; an RS code has
\(d_{\min}=n-k+1\).
\end{proof}

\begin{lemma}[Pair-interleaved Johnson bound]
\label{lem:pair-interleaved-johnson}
Let \(C=\RS_\F(D,k)\), let \(1\le A\le n\), and suppose
\(A^2>(k-1)n\).  For every received pair \(y\in(\F^D)^2\), the number
of codeword pairs \((p_0,p_1)\in C^2\) simultaneously agreeing with
\(y\) on at least \(A\) coordinates is at most
\[
 J_2(A)=\floor{
   \frac{n(A-k+1)}{A^2-(k-1)n}}.                  \tag{HD2}
 \label{eq:pair-johnson}
\]
\end{lemma}

\begin{proof}
Let the list have size \(L\), and for every member retain exactly \(A\)
agreement coordinates.  If \(d_x\) is the number of retained sets
containing \(x\), then \(\sum_xd_x=LA\).  Two distinct codeword pairs
can agree coordinatewise on at most \(k-1\) points, since at least one
nonzero difference polynomial has degree less than \(k\).  Therefore
\[
 \frac{L^2A^2/n-LA}{2}
 \le\sum_{x\in D}\binom{d_x}{2}
 \le\binom L2(k-1),
\]
where the first inequality is Cauchy--Schwarz.  Rearrangement gives
\(L(A^2-(k-1)n)\le n(A-k+1)\), and integrality gives
\eqref{eq:pair-johnson}.
\end{proof}

\begin{theorem}[Self-contained half-Johnson CA and MCA bound]
\label{thm:self-contained-half-johnson}
Let \(C=\RS_\F(D,k)\), let \(\varnothing\ne\Gamma\subseteq\F\), and
let \(r\) be an integer with \(0\le r\le n\).  Put \(A=n-2r\).  If
\(A>0\) and \(A^2>(k-1)n\), then
\[
 B_{C,\Gamma}^{\rm CA}(n-r)
 \le\min\{\abs\Gamma,1+(r+1)J_2(A)\}.             \tag{HD3}
 \label{eq:half-johnson-ca}
\]
If in addition \(2r\le n-k\), the same upper bound holds for
\(B_{C,\Gamma}^{\rm MCA}(n-r)\).
\end{theorem}

\begin{proof}
Fix a column-far received pair and let \(Z\subseteq\Gamma\) be its
CA-bad set.  The assertion is immediate when \(\abs Z\le1\), so fix an
anchor \(\gamma_0\in Z\).  For every
\(\gamma\in Z\setminus\{\gamma_0\}\), choose explaining codewords
\(c_{\gamma_0},c_\gamma\) and agreement sets of size at least \(n-r\),
and define
\[
 P_1=\frac{c_{\gamma_0}-c_\gamma}{\gamma_0-\gamma},
 \qquad P_0=c_{\gamma_0}-\gamma_0P_1.
\]
On the intersection of the two agreement sets, of size at least
\(n-2r=A\), the codeword pair \((P_0,P_1)\) agrees with the received
pair.  Hence all such pairs lie in a list of size at most \(J_2(A)\) by
\cref{lem:pair-interleaved-johnson}.

Fix one pair in that list and let \(E\) be its column-error set, of size
\(e\).  Column-farness gives \(e\ge r+1\).  If a slope \(\gamma\) rides
this pair, meaning its explaining codeword is \(P_0+\gamma P_1\), then
the nonzero affine coordinate errors vanish at at least \(e-r\)
coordinates of \(E\).  Each coordinate vanishes for at most one finite
slope, so double counting gives at most
\(e/(e-r)\le r+1\) riders.  Summing over the pair list and adding the
anchor proves \eqref{eq:half-johnson-ca}.  Under \(2r\le n-k\),
\cref{thm:exact-half-distance-sparse} says that the other branch of MCA
has only \(\min\{\abs\Gamma,r\}\) slopes.  Since
\(n(A-k+1)-(A^2-(k-1)n)=A(n-A)\ge0\), one has \(J_2(A)\ge1\), so this
tangent term is no larger than the displayed CA bound.
\end{proof}

\begin{remark}[The exact remaining prize certificate]
\label{rem:exact-prize-certificate}
At a proposed safe agreement \(a_+\), an unrestricted finite-row proof
must establish both
\[
 B_C^{\rm CA}(a_+)\le B^*,\qquad
 \sigma_C(a_+)\le B^*,\qquad
 B^*=\floor{2^{-128}\abs\F}.
\]
An unsafe construction at \(a_+-1\) does not prove either upper bound.
This identity is the shortest exact description of what remains for the
near-capacity rows not covered by
\cref{cor:half-distance-window-compiler}; the older quadratic compiler
\cref{cor:prize-window-compiler} is a self-contained subrange.
\end{remark}

The finite-slope restriction costs at most the point at infinity, which may
also be treated by exchanging the two coordinates.  The support-wise
condition is essential: explanation and nontriviality are tested on the
same set.

\begin{definition}[Exact-agreement witness incidence]
\label{def:exact-witness-incidence}
Assume \(k+1\le a\le n\).  For a received line \(r=(r_0,r_1)\), its
exact-\(a\) witness incidence is
\[
 \Wcal_a(r)=\left\{(\gamma,S,h):
 \begin{array}{l}
   \gamma\in\F,\quad S\in\binom Da,\quad h\in\F[X]_{<k},\\
   \operatorname{ev}_S(h)=(r_0+\gamma r_1)|_S,\quad
   \operatorname{NT}_C(r;S)
 \end{array}\right\}.
 \tag{2.1}\label{eq:exact-witness-incidence}
\]
Its slope projection is
\(\pi_\gamma(\gamma,S,h)=\gamma\), and its exact bad-slope set is
\(Z_a(r)=\pi_\gamma(\Wcal_a(r))\subseteq\F\).
\end{definition}

\begin{lemma}[Reduction to exact agreement]
\label{lem:exact-agreement-reduction}
For \(a\ge k+1\), the set \(Z_a(r)\) is precisely the set of slopes that
are MCA-bad at threshold \(a\) in \cref{def:mca-bad}.
\end{lemma}

\begin{proof}
An exact-\(a\) witness is a threshold-\(a\) witness.  Conversely, let
\((\gamma,S,h)\) witness badness with \(\abs S\ge a\).  At least one of
\(r_0|_S,r_1|_S\) is outside
\(\operatorname{ev}_S(\F[X]_{<k})\); call it \(u|_S\).  Interpolate \(u\)
on any fixed \(k\)-subset \(K\subseteq S\).  If the resulting polynomial
agreed with \(u\) at every point of \(S\setminus K\), it would explain
\(u\) on \(S\), a contradiction.  Hence some \(x\in S\setminus K\)
makes \(u\) unexplained on \(K\cup\{x\}\).  Extend this set to an
\(a\)-subset \(S'\subseteq S\).  Non-explanation persists on supersets,
while the same \(h\) explains \(r_0+\gamma r_1\) on \(S'\).
\end{proof}

\begin{proposition}[Exact support-atlas upper bound]
\label{prop:exact-support-upper}
For every Reed--Solomon row, every nonempty challenge set
\(\Gamma\subseteq\F\), and every \(k+1\le a\le n\),
\[
 B_{C,\Gamma}^{\rm MCA}(a)
 \le\min\left\{\abs\Gamma,\binom na\right\}.        \tag{2.2}
\]
This is a literal finite-row bound: it has no polynomial or
subexponential loss.
\end{proposition}

\begin{proof}
By \cref{lem:exact-agreement-reduction}, choose one noncommon exact-\(a\)
support for every bad slope.  One noncommon support carries at most one
slope, since two distinct explained slopes would separately explain the
two received words on that support.  The chosen-support map is therefore
injective.  Intersect with the challenge set.
\end{proof}

For an MDS code on a coordinate set \(S\), a \(b\)-coordinate
\emph{local test} is a subset \(T\in\binom Sb\); it accepts a word \(w\)
when \(w|_T\in C|_T\) and rejects it otherwise.
\begin{lemma}[Rejecting-test abundance, Jo
  {\cite[Lem.~4.1]{Jo26}}]\label{lem:rejecting-test-abundance}
Let \(C\subseteq\F^S\) be an \([m,k]\) MDS code, let \(w\notin C\), and let
\(k+1\le b\le m\).  Then at least
\[
                         \binom{m-1}{b-1}           \tag{2.3}
\]
sets \(T\in\binom Sb\) satisfy \(w|_T\notin C|_T\).  This number is
best possible.
\end{lemma}

\begin{proof}
We reproduce Jo's induction.  Use strong induction on \(m\), simultaneously
in \(k\).  For \(k=0\),
choose \(x\) with \(w(x)\ne0\); every \(b\)-set containing \(x\) rejects.
For \(k\ge1\), fix \(x\in S\) and put \(S'=S\setminus\{x\}\).
If \(w|_{S'}\in C|_{S'}\), choose \(c\in C\) agreeing with \(w\) on
\(S'\).  Since \(w\notin C\), the words differ at \(x\), and every
\(b\)-set containing \(x\) rejects by MDS uniqueness on its other
\(b-1\ge k\) coordinates.

Otherwise, if \(b<m\), induction gives \(\binom{m-2}{b-1}\) rejecting
\(b\)-sets avoiding \(x\); for \(b=m\) this contribution is zero.  Choose
\(c\in C\) with \(c(x)=w(x)\), put \(v=w-c\),
and shorten at \(x\):
\[
 C_x=\{d|_{S'}:d\in C,\ d(x)=0\}.
\]
This is \([m-1,k-1]\) MDS and \(v|_{S'}\notin C_x\).  A set
\(T\cup\{x\}\) accepts \(w\) precisely when \(T\) accepts \(v\) for
\(C_x\), so induction supplies \(\binom{m-2}{b-2}\) further rejecting
sets.  Pascal's identity proves the bound.  Sharpness is attained by
\(w=c+e_x\), where \(e_x\) is nonzero only at \(x\): exactly the sets
containing \(x\) reject.
\end{proof}

\begin{proof}[Proof of \cref{thm:intro-paving}]
Following Jo, the CA inequality is immediate.  Fix a received pair, let \(Z\subseteq\Gamma\)
be its bad-slope set, and, for every \(\gamma\in Z\), choose a witnessing support \(S_\gamma\) of
size at least \(a\).  The direction word is outside \(C|_{S_\gamma}\),
whereas the affine combination at \(\gamma\) is inside it.  For fixed
\(k+1\le b\le a\), \cref{lem:rejecting-test-abundance} gives at least
\(\binom{a-1}{b-1}\) subsets \(T\in\binom{S_\gamma}b\) on which the
direction rejects while the affine combination accepts.  One fixed \(T\)
cannot be owned by two distinct slopes, since subtraction would make the
direction accept on \(T\).  Hence
\[
                 \abs Z\binom{a-1}{b-1}\le\binom nb.       \tag{2.4}
\]
Take the floor, minimize over \(b\), and intersect with \(\Gamma\).
Coordinate permutations and invertible diagonal scalings preserve every
punctured membership test.
\end{proof}

\begin{remark}[Paving specialization]\label{rem:paving-saving}
At \(b=k+1\), the rejecting sets are the bases of the rank-\((k+1)\)
augmented-word paving matroid, giving the independent matroid
interpretation used in the earlier Chojecki manuscript \cite{RSMCARepo}.
Jo's strengthening is to use every test size.  The local lower bound is
sharp for a one-coordinate perturbation, so further universal progress
must exploit compatibility among different slopes.
\end{remark}


\section{Exact syndrome geometry}
\label{sec:syndrome-exact}

The \emph{syndrome} of a word is the result of applying a parity-check
matrix; it vanishes exactly on codewords.  Thus syndrome space forgets the
part of a received word already lying in the code and retains only its error
class.  This description separates possible error supports from distinct
line parameters.  Put \(R=n-k\), the redundancy (and syndrome-space
dimension), and choose the generalized Vandermonde parity-check matrix
\(H\) with columns
\[
 h_x=\lambda_x(1,x,\ldots,x^{R-1})^{\mathsf T},\qquad
 \lambda_x=\left(\prod_{y\in D\setminus\{x\}}(x-y)\right)^{-1}.
 \tag{3.1}\label{eq:parity-columns}
\]
Lagrange interpolation gives
\(\sum_{x\in D}\lambda_xx^j=0\) for \(0\le j\le n-2\), so this matrix
annihilates every degree-less-than-\(k\) evaluation.  Every set of at most
\(R\) columns is independent.  For \(E\subseteq D\), write
\(V_E=\Span\{h_x:x\in E\}\); it is exactly the set of syndromes of words
supported on \(E\).  We write \(s(u)=Hu^{\mathsf T}\) for the syndrome of
\(u\).

\begin{proposition}[Syndrome-line normal form]
\label{prop:syndrome-line-normal-form}
Let \(S=D\setminus E\).  A word \(u\) is explained on \(S\) if and only if
\(s(u)\in V_E\).  Consequently \(\gamma\) is MCA-bad on \(S\) for a pair
\((u_0,u_1)\) if and only if
\[
 s(u_0)+\gamma s(u_1)\in V_E,
 \qquad
 \{s(u_0),s(u_1)\}\not\subseteq V_E.
 \tag{3.2}\label{eq:syndrome-line}
\]
For fixed \(E\), there is at most one such finite slope.
\end{proposition}

\begin{proof}
If \(u=c+e\), where \(c\in C\) and \(e\) is supported on \(E\), then
\(s(u)=s(e)\in V_E\).  Conversely, if
\(s(u)=\sum_{x\in E}e_xh_x\), subtract the word supported on \(E\) with
values \(e_x\); the result has zero syndrome and belongs to \(C\).  Apply
this equivalence to the line point and to both coordinates.  Passing
\eqref{eq:syndrome-line} to \(\F^R/V_E\) shows that two distinct solutions
would put both projected syndromes at zero, contradicting badness.
This proof uses only that \(H\) is a full-rank parity-check matrix with
kernel \(C\); hence the normal form applies verbatim to every linear code
when \(H\) is replaced by any such parity check.
\end{proof}

Since \(\lambda_x\ne0\), one has
\([h_x]=[1:x:\cdots:x^{R-1}]\); these are the points of the rational
normal curve in projective syndrome space, while the weights
\(\lambda_x\) merely choose vector representatives.  A \emph{secant span} is the linear span
\(V_E\) of a selected set of those points.  We call the meeting of a
syndrome line with \(V_E\) \emph{transverse} when the entire line is not
contained in \(V_E\), equivalently when its two generators are not both in
\(V_E\).  The next result is a \emph{compiler} in the sense that it
translates the bad-slope problem, exactly and in both directions, into this
line--secant incidence problem.  A \emph{ray} means a one-dimensional
direction in syndrome or coefficient space; different raw witnesses may
represent the same ray and hence the same slope.

\begin{theorem}[Exact syndrome--secant compiler]
\label{thm:syndrome-secant-exact}
Put \(t=n-a\), and for a syndrome line
\(\ell_{y_0,y_1}=\{y_0+\gamma y_1:\gamma\in\F\}\) define
\[
 \Theta_t(y_0,y_1)=
 \abs{\{\gamma\in\F:\exists E\subseteq D,\ \abs E\le t,
 y_0+\gamma y_1\in V_E,
 \ \{y_0,y_1\}\not\subseteq V_E\}}.
 \tag{3.3}\label{eq:transverse-secant-count}
\]
Then
\[
 B_C^{\rm MCA}(a)=\max_{y_0,y_1\in\F^{n-k}}
                    \Theta_t(y_0,y_1).
 \tag{3.4}\label{eq:exact-secant-numerator}
\]
For each fixed \(E\), the affine line has at most one transverse
intersection parameter with \(V_E\).  For a Reed--Solomon code the
spaces \(V_E\) are the spans of at most \(t\) points of the weighted
rational normal curve in \eqref{eq:parity-columns}.  Thus the
higher-dimensional ray problem is exactly a transverse line--secant
incidence problem, with repeated intersection parameters deduplicated.
\end{theorem}

\begin{proof}
Apply \cref{prop:syndrome-line-normal-form} with \(S=D\setminus E\), and
take the union over \(\abs E\le t\).  The syndrome map is surjective, so
maximizing over received pairs is the same as maximizing over
\((y_0,y_1)\).  In the quotient \(\F^{n-k}/V_E\), a transverse
intersection solves
\(\overline y_0+\gamma\overline y_1=0\) with the two projected vectors
not both zero, and therefore has at most one finite solution.  The last
assertion is the Vandermonde description of the parity columns.
\end{proof}

A \emph{parity-column circuit} is a minimally linearly dependent set of
columns of \(H\).  The next lemma says that many distinct transverse rays
cannot all be supported inside an independent union of columns.

\begin{lemma}[Independent-union ray bound]
\label{lem:independent-union-rays}
Let every set of at most \(R\) parity columns be independent, put
\(t<R\), and fix a syndrome line.  For each slope in a transverse set
\(Z\), choose an error set \(E_\gamma\), \(\abs{E_\gamma}\le t\),
witnessing \eqref{eq:transverse-secant-count}.  If
\(U=\bigcup_{\gamma\in Z}E_\gamma\) has size at most \(R\), then
\(\abs Z\le t+1\).  Consequently, more than \(t+1\) transverse slopes
force the witness union to contain a parity-column circuit.
\end{lemma}

\begin{proof}
If \(\abs Z\le1\), there is nothing to prove; assume that \(Z\) contains
two distinct slopes.  Two distinct line points lie in \(V_U\), hence
\(y_0,y_1\in V_U\).
Column independence gives unique coefficient lifts
\(e_0,e_1\in\F^U\), and every chosen lift is
\(e_\gamma=e_0+\gamma e_1\).  Let \(U'\) be the set of \(s\) coordinates
where the pair \((e_0,e_1)\) is nonzero.  If \(s\le t\) and
\(E_\gamma\) contained all of \(U'\), then
\(y_0,y_1\in V_{U'}\subseteq V_{E_\gamma}\), contradicting the
transversality of this particular witness.  Hence some
\(x\in U'\setminus E_\gamma\) exists, and uniqueness of the lift gives
\(e_0(x)+\gamma e_1(x)=0\).  Each nonzero affine coordinate function has
at most one root, so \(\abs Z\le s\le t\).  If \(s>t\), every chosen lift of
weight at most \(t\) has at least \(s-t\) zero coordinates.  Double
counting the pairs \((\gamma,x)\) with
\(e_0(x)+\gamma e_1(x)=0\) gives
\(\abs Z(s-t)\le s\), whence
\(\abs Z\le\floor{s/(s-t)}\le t+1\).  If \(\abs U\le R\), independence
holds; the contrapositive proves the circuit assertion.
\end{proof}

\begin{theorem}[Bounded residual-kernel ray compiler]
\label{thm:bounded-residual-kernel-ray}
Let \(H\) be a parity-check matrix over \(\F\), let \(U\) be a set of
columns, and put
\[
 \kappa(U)=\dim_\F\ker(H_U:\F^U\longrightarrow\F^{n-k}).
\]
Fix a syndrome line \(y_0+\gamma y_1\), an integer \(t\ge0\), and a set
\(Z\) of finite transverse slopes.  Suppose that for every \(\gamma\in Z\)
there is \(E_\gamma\subseteq U\), \(\abs{E_\gamma}\le t\), such that
\[
 y_0+\gamma y_1\in V_{E_\gamma},\qquad
 \{y_0,y_1\}\not\subseteq V_{E_\gamma}.
\]
Then
\[
                  \abs Z\le(t+1)\abs{\F}^{\kappa(U)}. \tag{RC$_{\rm ker}$}
\]
For a Reed--Solomon parity check of redundancy \(R=n-k\),
\(\kappa(U)=\max\{0,\abs U-R\}\).  Hence a certified residual chart with
\((\abs U-R)_+\log\abs\F=o(n)\) has a direct subexponential ray payment
with the displayed exact finite loss.
\end{theorem}

\begin{proof}
The assertion is immediate when \(\abs Z\le1\).  Otherwise two line
points belong to \(V_U\), so \(y_0,y_1\in V_U\).  Choose lifts
\(b_0,b_1\in\F^U\) with \(H_Ub_i=y_i\).  For every \(\gamma\), choose a
lift \(c_\gamma\) supported in \(E_\gamma\) and put
\(z_\gamma=c_\gamma-b_0-\gamma b_1\in\ker H_U\).

Fix \(z\in\ker H_U\), and consider the slopes for which
\(z_\gamma=z\).  Put
\[
 U_z=\{x\in U:(b_0(x)+z(x),b_1(x))\ne(0,0)\},
 \qquad s=\abs{U_z}.
\]
Every active coordinate is a nonzero affine function of \(\gamma\), and
therefore vanishes for at most one slope.  Hence the total number of active
zero incidences in this kernel class is at most \(s\).  If \(s>t\), the
weight bound on \(c_\gamma\) forces at least \(s-t\) active zeros, so
\[
 \abs{Z_z}(s-t)\le s,\qquad
 \abs{Z_z}\le\floor{s/(s-t)}\le t+1.
\]
If \(s\le t\), transversality forces
at least one active zero: without one,
\(U_z=\supp c_\gamma\subseteq E_\gamma\), while both \(b_0+z\) and
\(b_1\) are supported in \(U_z\), putting \(y_0,y_1\) in
\(V_{E_\gamma}\).  Thus this kernel class contains at most \(s\le t\)
slopes.  Summing over the \(\abs{\F}^{\kappa(U)}\) kernel elements proves
the bound.  The MDS column property gives the Reed--Solomon formula for
\(\kappa(U)\).
\end{proof}

\begin{theorem}[Single MDS-circuit ray compiler]
\label{thm:single-mds-circuit-ray}
Let \(H\) be a Reed--Solomon parity check of redundancy \(R\), let
\(U\subseteq D\) have size \(R+1\), put \(t<R\), and fix a syndrome line
\(y_0+\gamma y_1\).  If every slope in a transverse set \(Z\) has a witnessing error set
\(E_\gamma\subseteq U\) of size at most \(t\), then
\[
                         \abs Z\le\binom{R+1}{2}.   \tag{RC$_{\rm circ}$}
\]
Thus one fixed MDS-circuit profile has a field-independent polynomial ray
payment.
\end{theorem}

\begin{proof}
There is nothing to prove when \(\abs Z\le1\).  Otherwise two distinct
line points lie in \(V_U\); subtracting them and solving for the constant
term gives \(y_0,y_1\in V_U\).  Choose lifts
\(b_0,b_1\in\F^U\).  The circuit kernel is generated by a vector
\(z\) whose coordinates are all nonzero, since a zero coordinate would
give a proper dependent subset of the minimal circuit \(U\).  For each slope choose
\(c_\gamma\in\F\) so that
\(e_\gamma=b_0+\gamma b_1+c_\gamma z\) is supported in
\(E_\gamma\).  For \(x\in U\), set
\[
             f_x(\gamma)=-(b_0(x)+\gamma b_1(x))/z(x).
\]
A zero coordinate is an equality \(c_\gamma=f_x(\gamma)\), and at least
\(R+1-t\ge2\) such equalities hold.

Every transverse slope makes two nonidentical functions \(f_x,f_{x'}\)
equal.  Indeed, otherwise all vanishing coordinates belong to one class
\(C\) of identical affine functions, say
\(f_x=\alpha+\beta\gamma\) on \(C\), and no function outside \(C\) has
that value at \(\gamma\).  Then
\(\supp e_\gamma=U\setminus C\subseteq E_\gamma\), while
\(b_0+\alpha z\) and \(b_1+\beta z\) are also supported in
\(U\setminus C\).  This puts both \(y_0,y_1\) in \(V_{E_\gamma}\), a
contradiction.  Two nonidentical affine functions agree at at most one
finite slope.  Charging each slope to one unordered pair proves the
displayed bound.
\end{proof}



\section{The mean-overlap regime}\label{sec:mean-overlap}

For a fixed received line, the exact bad-slope set is therefore the set of
unique cancelling line parameters contributed by the transverse quotient
equations.  For each \(E\), \(\abs E\le n-a\), project the two
syndromes to \(\F^R/V_E\).  If the second projection is nonzero and the
first is a scalar multiple of it, record the unique cancelling scalar;
otherwise record nothing.  Deduplicating these scalars gives the numerator.
In particular, the number of supports, the number of pairs of supports, and
the number of distinct slopes are three different quantities.

We call \(3(n-a)\le d-1\) the \emph{deep regime}: the allowed error
radius \(n-a\) is so small relative to the minimum codeword weight \(d\)
that the union of three permitted error supports still cannot contain a
nonzero codeword support.

\begin{theorem}[Deep-regime upper bound]\label{thm:deep-regime-upper}
Let \(C\subseteq\F^n\) be linear with minimum nonzero weight \(d\), let
\(0\le a\le n\), put \(r=n-a\), and assume \(3r\le d-1\).  Then
\[
 B_C^{\rm MCA}(a)\le r+1.
 \tag{3.5}\label{eq:deep-upper}
\]
\end{theorem}

\begin{proof}
Fix \((u_0,u_1)\), and let \(Z\) be its set of MCA-bad slopes at agreement
\(a\).  If \(\abs Z\le r+1\), there is nothing to prove.  Otherwise choose
distinct \(\gamma_1,\gamma_2\in Z\), together with witnessing codewords
\(p_1,p_2\) and error sets \(E_1,E_2\), each of size at most \(r\), such
that \(u_0+\gamma_i u_1=p_i\) off \(E_i\).  Solving the two affine
equations gives \(c_0,c_1\in C\) for which the error pair
\(e_i=u_i-c_i\) is supported on \(T=E_1\cup E_2\), \(\abs T\le2r\).

For any \(\gamma\in Z\), choose a witnessing explaining codeword \(p\).  The
codeword \(p-(c_0+\gamma c_1)\) is supported on at most \(3r\le d-1\)
positions and is zero.  Thus \(\wt(e_0+\gamma e_1)\le r\).  Let \(T'\) be
the support of the error pair.  If \(\abs{T'}\ge r+1\), every
\(\gamma\in Z\) makes at least \(\abs{T'}-r\) nonzero affine coordinate
functions \(e_{0,x}+\gamma e_{1,x}\) vanish.  Each coordinate has at most
one root, so \(\abs Z(\abs{T'}-r)\le\abs{T'}\) and \(\abs Z\le r+1\).

It remains that \(\abs{T'}\le r\).  If \(\gamma\) is MCA-bad on a set
\(S\), \(\abs S\ge n-r\), its explaining codeword equals
\(c_0+\gamma c_1\), since their difference would otherwise have weight at
most \(2r\le d-1\).  The set \(S\) must meet \(T'\), and at a point of
\(S\cap T'\) the nonzero pair satisfies
\(e_{0,x}+\gamma e_{1,x}=0\).  Such a point determines \(\gamma\), giving
at most \(\abs{T'}\le r\) bad slopes.
\end{proof}

Recall that for a nonempty challenge set \(\Gamma\subseteq\F\),
\(B_{C,\Gamma}^{\rm MCA}(a)\) denotes the same maximum with slopes restricted to
\(\Gamma\).

The following ``tangent floor'' is an elementary family of neighboring
coordinate-span intersections.  The name does not assert a hidden
differentiability or multiplicity condition.

\begin{proposition}[Universal MDS tangent floor]\label{prop:universal-tangent-floor}
For every linear \([n,k]\) MDS code, every \(a\ge k+1\), and every nonempty
\(\Gamma\subseteq\F\),
\[
 B_{C,\Gamma}^{\rm MCA}(a)
 \ge\min\{\abs\Gamma,n-a+1\}.
 \tag{3.6}\label{eq:tangent-floor}
\]
\end{proposition}

\begin{proof}
Choose a full-rank parity-check matrix \(H\) for \(C\).  The MDS property
is equivalent to independence of every set of at most \(n-k\) columns; the
syndrome-line normal form of \cref{prop:syndrome-line-normal-form} is linear
algebra and applies to this \(H\).  Put \(j=n-a\),
\(t=\min\{\abs\Gamma,j+1\}\), and choose distinct
\(\gamma_1,\ldots,\gamma_t\in\Gamma\) and independent parity columns
\(h_{x_1},\ldots,h_{x_t}\).  Define
\(s_1=\sum_i h_{x_i}\) and \(s_0=-\sum_i\gamma_i h_{x_i}\), and lift them
to received words.  For \(E_i=\{x_\ell:\ell\ne i\}\),
\(s_0+\gamma_i s_1\in V_{E_i}\), whereas \(s_1\notin V_{E_i}\).
\Cref{prop:syndrome-line-normal-form} makes every \(\gamma_i\) bad on
\(D\setminus E_i\), which has size at least \(a\).
\end{proof}

\begin{corollary}[Exact deep numerator]\label{cor:exact-deep-numerator}
If \(C\subseteq\F^n\) is a linear \([n,k]\) MDS code,
\(1\le k<n\), \(k+1\le a\le n\),
\(\varnothing\ne\Gamma\subseteq\F\), and \(3(n-a)\le n-k\), then
\[
 B_{C,\Gamma}^{\rm MCA}(a)
 =\min\{\abs\Gamma,n-a+1\}.
\]
\end{corollary}

\begin{proof}
The upper bound is \cref{thm:deep-regime-upper}, intersected with
\(\Gamma\).  The hypotheses imply \(a\ge k+1\), so
\cref{prop:universal-tangent-floor} gives equality.
\end{proof}

\begin{remark}[The zero-radius normalization]
Under the formal affine sampler, \cref{cor:exact-deep-numerator} at
\(a=n\) gives \(B_C^{\rm MCA}(n)=1\), and hence
\(e_{\rm MCA}(C,0)=1/\abs\F\) for every proper linear MDS code.  This also
follows directly: in \(\F^n/C\), a noncontained affine line meets
the zero coset in at most one parameter, and the bound is attained by a
line through the zero coset.  A value \(2/\abs\F\) belongs to a different
or loose convention and must not be imported into the affine definition
of \cite{ABF26}.
\end{remark}

The next lemma isolates the deterministic mechanism behind both extensions
of the deep range.

\begin{lemma}[Large-overlap collapse]
\label{lem:large-overlap-collapse}
Fix a received pair and an exact agreement \(a=n-r\ge k+1\).  Choose one
noncommon exact witness support \(S_z\) and explaining polynomial \(p_z\)
for every slope \(z\) in its bad set \(Z\).  If two distinct chosen
supports satisfy
\[
                    \abs{S_z\cap S_w}\ge k+r,       \tag{MO1}
\label{eq:large-overlap-collapse}
\]
then \(\abs Z\le r+1\).
\end{lemma}

\begin{proof}
Put
\[
 c_1=\frac{p_z-p_w}{z-w},\qquad c_0=p_z-zc_1.
\]
On \(S_z\cap S_w\), the two line equations imply that the received pair
equals \((c_0,c_1)\).  Let \(C_0\) be the full common support of this
codeword pair and write \(c=\abs{C_0}\ge k+r\).  For every \(u\in Z\),
\(\abs{S_u\cap C_0}\ge a+c-n\ge k\), so the explaining polynomial
\(p_u\) equals \(c_0+uc_1\).

Subtract this codeword pair from the received pair.  At every coordinate
outside \(C_0\) at least one of the two residual values is nonzero, and a
coordinate can cancel at at most one finite slope.  A noncommon witness
for \(u\) must contain at least \(\max\{1,a-c\}\) such cancellations:
the support-size condition gives \(a-c\), and one is still necessary when
\(a\le c\).  Double counting therefore gives
\[
 \abs Z\le
 \floor{\frac{n-c}{\max\{1,a-c\}}}\le r+1.        \tag{MO2}
\label{eq:large-overlap-count}
\]
Indeed, for \(c\ge a\) the numerator is at most \(r\); for \(c=a-1\)
it is \(r+1\); and, with \(x=a-c\ge2\), the quotient is
\((r+x)/x\le r+1\).
\end{proof}

\begin{theorem}[Quadratic mean-overlap theorem]
\label{thm:quadratic-mean-overlap}
Let \(C=\RS_\F(D,k)\), put \(R=n-k\), let
\(0\le r\le R-1\), and let \(\varnothing\ne\Gamma\subseteq\F\).  If
\[
 (n-r)^2\ge n(k+r),
 \qquad\text{equivalently}\qquad
 r^2\ge n(3r-R),                                  \tag{MO3}
\label{eq:quadratic-overlap-condition}
\]
then
\[
 B_{C,\Gamma}^{\rm MCA}(n-r)=\min\{\abs\Gamma,r+1\}. \tag{MO4}
\label{eq:quadratic-overlap-exact}
\]
The same upper bound holds for every linear \([n,k]\) MDS code.
\end{theorem}

\begin{proof}
The universal tangent construction gives the lower bound.  For the upper
bound, fix a received pair and its challenge-restricted bad set \(Z\).
If \(\abs Z\le r+1\) there is nothing to prove.  Otherwise choose
\(r+2\) slopes and distinct exact witness supports
\(S_1,\ldots,S_{r+2}\), each of size \(a=n-r\).  Suppose every pair had
intersection at most \(k+r-1\), and put
\(d_x=\#\{i:x\in S_i\}\).  Then
\[
 \sum_x\binom{d_x}{2}\le\binom{r+2}{2}(k+r-1),    \tag{MO5}
\]
whereas Cauchy--Schwarz and \(a^2/n\ge k+r\) give
\[
 \sum_x\binom{d_x}{2}
 \ge\frac{(r+2)^2a^2/n-(r+2)a}{2}
 \ge\frac{(r+2)^2(k+r)-(r+2)a}{2}.                \tag{MO6}
\]
The last lower bound exceeds the upper bound in \textup{(MO5)} by
\[
 \frac{r+2}{2}\bigl(k+3r-n+1\bigr).
\]
When \(3r>R\) this is positive.  When \(3r\le R\), the exact deep theorem
already gives the desired upper bound.  Thus in the remaining case some
pair satisfies \eqref{eq:large-overlap-collapse}, and
\cref{lem:large-overlap-collapse} contradicts \(\abs Z\ge r+2\).
The exact-support reduction and large-overlap lemma use only the MDS
information-set property, proving the stated extension.
\end{proof}

\begin{corollary}[Exact quadratic staircase]
\label{cor:mean-overlap-exact}
In the setting of \cref{thm:quadratic-mean-overlap}, put
\[
 r_{\rm quad}(n,k)=
 \floor{\frac{3n-\sqrt{n(5n+4k)}}2}.
 \tag{MO7}\label{eq:quadratic-radius}
\]
For every \(0\le r\le r_{\rm quad}(n,k)\) and every nonempty challenge
set \(\Gamma\),
\[
 B_{C,\Gamma}^{\rm MCA}(n-r)=\min\{\abs\Gamma,r+1\}.
\]
If \(k/n\to\rho\), then
\(r_{\rm quad}(n,k)/n\to
\alpha_\rho=(3-\sqrt{5+4\rho})/2\).
\end{corollary}

\begin{proof}
The smaller root of \(r^2-3nr+n(n-k)=0\) is the real number inside the
floor in \eqref{eq:quadratic-radius}.  Apply
\cref{thm:quadratic-mean-overlap}; the limiting formula follows by
division by \(n\).
\end{proof}

\begin{proposition}[Limit of pairwise-overlap information]
\label{prop:pairwise-overlap-limit}
Let \(k/n\to\rho\in(0,1)\), \(r/n\to\alpha\in(0,1)\), and suppose
\[
 \alpha^2<3\alpha-(1-\rho).
 \tag{MO8}\label{eq:pairwise-overlap-limit}
\]
For all sufficiently large \(n\), there are \(r+2\) distinct
\(r\)-subsets \(T_i\) such that, whenever
\(j=3r-(n-k)\le r\),
\(\abs{T_i\cap T_\ell}\le j-1\) for every \(i\ne\ell\).  If \(j>r\),
that intersection condition is automatic.  Hence pairwise support overlap
alone cannot extend the exact theorem past the quadratic boundary; further
progress must use the syndrome, determinant, or fiber incidence.  This is
a limitation of the method, not a construction of an MCA-bad line.
\end{proposition}

\begin{proof}
When \(j\le r\), choose \(r+2\) independent uniform \(r\)-subsets.
The intersection of a fixed pair is hypergeometric with mean
\(r^2/n=(\alpha^2+o(1))n\), whereas
\(j=(3\alpha-(1-\rho)+o(1))n\).  More explicitly,
\[
 \Pr(\abs{T_1\cap T_2}=s)
 =\frac{\binom rs\binom{n-r}{r-s}}{\binom nr}.
\]
After writing \(s=xn\), Stirling's formula gives the exponent
\[
 \alpha h(x/\alpha)+(1-\alpha)
 h((\alpha-x)/(1-\alpha))-h(\alpha),
\]
which is strictly concave and has its unique maximum zero at
\(x=\alpha^2\).  The strict gap in
\eqref{eq:pairwise-overlap-limit}, followed by summing at most \(n+1\)
terms, therefore gives probability \(e^{-\Omega(n)}\) that one fixed
intersection reaches \(j\).  A union bound over \(O(n^2)\) pairs is
smaller than one.  A
successful family is automatically distinct, because equal members have
intersection \(r\ge j\).  When \(j>r\), choose any \(r+2\) distinct
\(r\)-sets.
\end{proof}

The following complementary theorem can be stronger when the overlap
defect \(j=3r-R\) is small.  Its only new input is a packing inequality;
in particular it uses no smoothness or Fourier estimate.

\begin{theorem}[Exact overlap--packing criterion]
\label{thm:overlap-packing-exact}
Let \(C=\RS_\F(D,k)\), put \(R=n-k\), and fix an integer
\(0\le r\le R-1\).  Set \(a=n-r\) and \(j=3r-R\).  If either
\[
 j\le0,
 \tag{OP0}\label{eq:op-deep}
\]
or
\[
 1\le j\le r,
 \qquad
 \binom nj\le(r+1)\binom rj,
 \tag{OP}\label{eq:op-packing}
\]
then
\[
                    B_C^{\rm MCA}(n-r)=r+1.       \tag{OP$'$}
\label{eq:op-exact}
\]
The same upper bound holds for every linear \([n,k]\) MDS code; the
matching lower bound holds whenever the field contains at least \(r+1\)
distinct slope parameters.
\end{theorem}

\begin{proof}
Condition \eqref{eq:op-deep} is
\cref{cor:exact-deep-numerator}, so assume
\eqref{eq:op-packing}.  Fix an arbitrary received pair \((f,g)\), let
\(Z\) be its bad-slope set, and use
\cref{lem:exact-agreement-reduction} to choose, for every \(z\in Z\), an
exact witness support \(S_z\), \(\abs{S_z}=a\), and an explaining
polynomial \(p_z\in\F[X]_{<k}\).  A single noncommon support cannot carry
two distinct slopes, so the supports chosen for distinct elements of
\(Z\) are distinct.

For the stated MDS extension, the same exact-support reduction is
available: interpolate the unexplained component on any \(k\)-subset of
the original witness support; because every \(k\)-subset is an information
set, nonexplanation is detected at one further coordinate, and that
\((k+1)\)-set can be enlarged to size \(a\).

First suppose that two of them have a large intersection:
\[
 \abs{S_z\cap S_w}\ge n+k-a=k+r
 \qquad(z\ne w).
 \tag{OP1}\label{eq:op-large-overlap}
\]
Then \cref{lem:large-overlap-collapse} gives \(\abs Z\le r+1\).

It remains to consider the complementary case, in which every pair of
chosen supports satisfies
\(\abs{S_z\cap S_w}\le n+k-a-1\).  Put
\(T_z=D\setminus S_z\), so \(\abs{T_z}=r\).  For distinct slopes,
\[
 \abs{T_z\cap T_w}
 =n-2a+\abs{S_z\cap S_w}
 \le3r-R-1=j-1.                                  \tag{OP3}
\label{eq:op-small-overlap}
\]
No \(j\)-subset of \(D\) can therefore lie in two different \(T_z\)'s.
Counting pairs \((z,J)\) with \(J\in\binom{T_z}{j}\) yields
\[
                 \abs Z\binom rj\le\binom nj\le
                 (r+1)\binom rj,
\]
and hence \(\abs Z\le r+1\).  The universal tangent construction gives
the reverse inequality.  The proof used only linearity, the fact that
agreement of two codewords on \(k\) coordinates forces equality, and the
exact-support reduction; these are MDS properties.  The coordinate
tangent construction gives the stated MDS lower bound.
\end{proof}

\begin{proposition}[General overlap--packing upper bound]
\label{prop:general-overlap-packing-upper}
Let \(C\subseteq\F^n\) be a linear \([n,k]\) MDS code, let
\(\varnothing\ne\Gamma\subseteq\F\), let \(0\le r\le n-k-1\), and
suppose \(1\le j=3r-(n-k)\le r\).  Let \(A(n,r,j)\) be the largest size of a
family of \(r\)-subsets of an \(n\)-set whose distinct members intersect
in at most \(j-1\) points.  Then
\[
 B_{C,\Gamma}^{\rm MCA}(n-r)
 \le\min\left\{\abs\Gamma,
     \max\left(r+1,A(n,r,j)\right)\right\}          \tag{OP5}
\label{eq:general-overlap-packing-upper}
\]
and
\[
 A(n,r,j)\le
 \floor{\frac{\binom nj}{\binom rj}}.              \tag{OP6}
\label{eq:packing-number-bound}
\]
\end{proposition}

\begin{proof}
For a fixed received pair, either two chosen exact witness supports have
intersection at least \(k+r\), in which case
\cref{lem:large-overlap-collapse} gives at most \(r+1\) slopes, or every
pair of complementary \(r\)-sets intersects in at most \(j-1\) points.
The latter family has size at most \(A(n,r,j)\).  Finally, no \(j\)-subset
can lie in two members of such a packing, so counting contained
\(j\)-subsets proves \eqref{eq:packing-number-bound}.
\end{proof}

\begin{corollary}[Quadratic staircase at the official rates]
\label{thm:prize-rate-overlap}
Let \(D\subseteq\F\) have size \(n=2^e\), put
\(C=\RS_\F(D,k)\), let
\(\rho=k/n\in\{1/2,1/4,1/8,1/16\}\), and put
\[
 r_\rho(n)=\floor{
 n\frac{3-\sqrt{5+4\rho}}2}.
\]
For every \(n\)-point Reed--Solomon domain over every field,
\[
 B_C^{\rm MCA}(n-r)=r+1
 \qquad(0\le r\le r_\rho(n)).                     \tag{OP4}
\label{eq:prize-rate-staircase}
\]
\end{corollary}

\begin{proof}
This is \cref{cor:mean-overlap-exact}, since the displayed integer is
\(r_{\rm quad}(n,k)\).
\end{proof}

\begin{corollary}[Self-contained quadratic-range target compiler]
\label{cor:prize-window-compiler}
In the setting of \cref{thm:prize-rate-overlap}, put
\(q=\abs\F\) and \(B=\floor{q/2^{128}}\).  If
\[
 1\le B\le M_\rho(n)
 \defeq\min\{r_\rho(n)+1,n-k-1\},
\]
then
\[
 e_{\rm MCA}(C,\delta)\le2^{-128}
 \quad\Longleftrightarrow\quad
 \floor{\delta n}\le B-1.                        \tag{PW}
\label{eq:prize-window-iff}
\]
Thus every such row is solved throughout the field window
\(2^{128}\le q<2^{128}(M_\rho(n)+1)\): its largest safe grid radius is
\((B-1)/n\), while its real transition supremum is \(B/n\) and that
endpoint is unsafe.  If \(q<2^{128}\), even radius zero is unsafe.
\end{corollary}

\begin{proof}
For \(r\le B-1\), \cref{thm:prize-rate-overlap} gives the exact
numerator \(r+1\le B\le q/2^{128}\).  At \(r=B\), the tangent floor
applies because \(B\le n-k-1\), and gives at least
\(B+1>q/2^{128}\) slopes.  Monotonicity proves the
equivalence and the endpoint statements.  If \(B=0\), the tangent floor
at \(r=0\) gives the exact numerator one.
\end{proof}


\section{Prior Reed--Solomon completion below half distance}
\label{sec:literature-completion}

The quadratic theorem is self-contained and useful at finite parameters,
but it is not the largest known RS range at constant rate.  We now make
the relevant implication of \cite{BCHKS25} explicit at the integer
numerator level.  This section is deliberately separated from
\cref{sec:mean-overlap}: the input below belongs to the literature,
whereas \cref{thm:quadratic-mean-overlap} is the finite MDS argument of
this paper.

Put \(q=\abs\F\).  For a word \(u\) and a received pair \((u_0,u_1)\), write
\(d(u,C)\) and \(d_2((u_0,u_1),C^2)\) for their relative Hamming and
relative column distances from \(C\) and \(C^2\), respectively.  For
\(0\le\alpha<\beta<1\), define the two-radius correlated-agreement error
by
\[
 e_{\rm CA}(C,\alpha,\beta)
 =\frac1q\max_{\substack{u_0,u_1\in\F^D\\
                  d_2((u_0,u_1),C^2)>\beta}}
   \#\{\gamma\in\F:d(u_0+\gamma u_1,C)\le\alpha\}.       \tag{LC0}
 \label{eq:two-radius-ca}
\]
The one-radius notation \(e_{\rm CA}(C,\alpha)\) means that the far-pair
radius is also \(\alpha\).

\begin{theorem}[BCHKS two-radius bound]
\label{thm:bchks-two-radius}
Let \(C=\RS_\F(D,k)\), put \(q=\abs\F\),
and \(\Delta=(n-k+1)/n\).  If
\(\Delta/3\le\alpha\le\Delta/2-1/n\) and
\(\alpha<\beta<1\), then
\[
 e_{\rm CA}(C,\alpha,\beta)
 \le \frac1q\max\left\{
   \frac{\Delta-\alpha}{\alpha(\Delta-2\alpha)},
   \frac{\beta}{\beta-\alpha}
                         \right\}.              \tag{BCHKS}
 \label{eq:bchks-two-radius}
\]
\end{theorem}

\begin{proof}
BCHKS \cite[Thm.~1.3]{BCHKS25} uses degree at most \(k-1\), hence
dimension \(k\) and relative minimum distance \(\Delta\).  Fix a received
pair at column distance greater than \(\beta\), and let \(Z\) be its set of
\(\alpha\)-close slopes.  If
\[
 \abs Z\ge \frac{\Delta-\alpha}{\alpha(\Delta-2\alpha)},
\]
that theorem gives column distance at most
\((1+1/(\abs Z-1))\alpha\).  The assumed distance greater than \(\beta\)
therefore forces \(\abs Z<\beta/(\beta-\alpha)\).  Otherwise \(\abs Z\)
is already below the first displayed threshold.  Taking the larger bound,
then maximizing over the received pair and dividing by \(q\), proves
\eqref{eq:bchks-two-radius}.  This also verifies the dimension-\(k\)
normalization; all denominators are positive on the stated range.
\end{proof}

\begin{theorem}[Integral BCHKS completion]
\label{thm:literature-integral-completion}
Let \(C=\RS_\F(D,k)\), let \(q=\abs\F\), put \(R=n-k\), and let
\(r\in\Z\) satisfy
\[
 3r\ge R+1,\qquad 2r<R.                           \tag{LC1}
 \label{eq:literature-range}
\]
Then
\[
 B_C^{\rm MCA}(n-r)
 \le
 \floor{\max\left\{
       \frac{n(R+1-r)}{r(R+1-2r)},\ r+1
                  \right\}}.                    \tag{LC2}
 \label{eq:literature-integral-upper}
\]
Consequently, if
\[
             n(R-r)<r(r+2)(R-2r),                \tag{LC3}
 \label{eq:literature-integral-condition}
\]
then, for every nonempty \(\Gamma\subseteq\F\),
\[
 B_{C,\Gamma}^{\rm MCA}(n-r)=\min\{\abs\Gamma,r+1\}.
 \tag{LC4}\label{eq:literature-integral-exact}
\]
\end{theorem}

\begin{proof}
Put \(\delta_{\rm fld}=r/n\) and, for \(0<\theta<1\), put
\(\delta_{\rm int}=(r+\theta)/n\).  Hamming distances belong to
\(\{0,1/n,\ldots,1\}\), so replacing the far-pair radius
\(\delta_{\rm fld}\) by \(\delta_{\rm int}\) does not change the
far-pair condition; this is the grid observation of
\cite[Rem.~4.10]{ABF26}.  The hypotheses in
\eqref{eq:literature-range} give
\(\Delta/3\le\delta_{\rm fld}\le\Delta/2-1/n\), and
\(\delta_{\rm fld}<\delta_{\rm int}<1\).  Hence
\cref{thm:bchks-two-radius} gives
\[
 q\,e_{\rm CA}(C,r/n)
 \le \max\left\{
  \frac{n(R+1-r)}{r(R+1-2r)},\frac{r+\theta}{\theta}
               \right\}.                        \tag{LC5}
 \label{eq:bchks-specialized}
\]
On the integer grid,
\(q e_{\rm CA}(C,r/n)=B_C^{\rm CA}(n-r)\).  The second inequality in
\eqref{eq:literature-range} also places the radius below
\(d_{\min}/(2n)\), because \(d_{\min}=R+1\).  By
\cref{thm:exact-half-distance-sparse},
\[
 B_C^{\rm MCA}(n-r)
   =\max\{B_C^{\rm CA}(n-r),\min\{q,r\}\}.
\]
The right side of \eqref{eq:bchks-specialized} is strictly larger than
\(r\), because \((r+\theta)/\theta>r\).  It therefore bounds both
branches of this maximum and hence bounds the integer
\(B_C^{\rm MCA}(n-r)\).
For the full-field maximum one may alternatively invoke the known
below-half equality \cite[Lem.~4.6]{ABF26}; the sparse decomposition is
kept here because it is exact before maximizing over a restricted
challenge set.

Letting \(\theta\to1^-\) proves that this integer is at most the real
maximum in \eqref{eq:literature-integral-upper}, and taking the floor
proves \eqref{eq:literature-integral-upper}.  Under
\eqref{eq:literature-integral-condition}, positivity of the denominators
and \(r>0\) give
\[
 \frac{n(R+1-r)}{r(R+1-2r)}
 \le \frac{n(R-r)}{r(R-2r)}<r+2.
\]
Since \(r+1\) is integral, the floor in
\eqref{eq:literature-integral-upper} is at most \(r+1\), so the
full-field numerator is at most \(r+1\).
Restriction to \(\Gamma\) gives the upper bound
\(\min\{\abs\Gamma,r+1\}\), and
\cref{prop:universal-tangent-floor} gives the reverse bound.
\end{proof}

\begin{remark}[Priority]
The substance of the upper bound in
\cref{thm:literature-integral-completion} is prior work.  In particular,
\cite[Cor.~1.4]{BCHKS25} states a sharp lossless conclusion under its
explicit distance and radius hypotheses, and \cite[Rem.~2.5]{BCHKS25}
discusses the matching
perturbation construction.  Formula \eqref{eq:bchks-specialized} is
included to expose the integer limit and to make the challenge-subset and
target compilers below completely transparent.
\end{remark}

\begin{corollary}[Positive density of exact radii]
\label{cor:positive-radius-density}
Let \(C_n=\RS_{\F_n}(D_n,k_n)\) be an RS row sequence with
\(k_n/n\to\rho\in(0,1)\), let
\(\varnothing\ne\Gamma_n\subseteq\F_n\), and let \(r_n\in\Z\) satisfy
\(r_n/n\to\tau\), and suppose
\(0<\tau<(1-\rho)/2\).  Then
\[
 B_{C_n,\Gamma_n}^{\rm MCA}(n-r_n)
 =\min\{\abs{\Gamma_n},r_n+1\}                   \tag{LC6}
 \label{eq:positive-radius-exact}
\]
for all sufficiently large \(n\).  More uniformly, for every
\(\eta>0\), equality holds throughout
\[
 0\le r\le\left(\frac{1-\rho}{2}-\eta\right)n
\]
for every sufficiently large row whose rate tends to \(\rho\).
Consequently the exact staircase contains an initial interval of
normalized length \((1-\rho)/2-o(1)\), has lower density at least
\((1-\rho)/2\) in the full Hamming-radius interval, and has asymptotic
relative density one in the unique-decoding interval.
\end{corollary}

\begin{proof}
The quadratic theorem covers every sequence whose limiting radius lies
below \(\alpha_\rho=(3-\sqrt{5+4\rho})/2\).  In the remaining compact
subinterval of \((0,(1-\rho)/2)\), the hypotheses
\eqref{eq:literature-range} hold for large \(n\), while
\[
 \frac{n(R-r)}{r(R-2r)}=O_{\rho,\eta}(1)
 \quad\text{and}\quad r+1=\Theta_{\rho,\eta}(n).
\]
Thus \eqref{eq:literature-integral-condition} holds and
\cref{thm:literature-integral-completion} applies.  The uniform statement
follows by the same compactness argument, using the deep theorem near
radius zero.  Counting integer radii gives the density assertions.
\end{proof}

\begin{corollary}[Official-rate staircase]
\label{cor:official-rate-half-staircase}
Let \(D\subseteq\F\) have size \(n=2^e\ge128\), put
\(C=\RS_\F(D,k)\), let \(k=\rho n\), and
\(\rho\in\{1/2,1/4,1/8,1/16\}\).  Define
\[
 r_\rho^\sharp=\begin{cases}
 n/4-3,&\rho=1/2,\\
 3n/8-2,&\rho=1/4,\\
 7n/16-2,&\rho=1/8,\\
 15n/32-2,&\rho=1/16.
 \end{cases}                                      \tag{LC7}
 \label{eq:official-half-endpoints}
\]
Then for every nonempty \(\Gamma\subseteq\F\),
\[
 B_{C,\Gamma}^{\rm MCA}(n-r)=\min\{\abs\Gamma,r+1\}
 \qquad(0\le r\le r_\rho^\sharp).              \tag{LC8}
 \label{eq:official-half-staircase}
\]
\end{corollary}

\begin{proof}
Put \(F(r)=r(r+1)(R-2r)-n(R-r)\), where \(R=n-k\), and
\(L=\ceil{(R+1)/3}\).  The deep theorem covers \(r<L\).  On the
positive real axis, \(F'(r)=-6r^2+2(R-2)r+(R+n)\) has exactly one
positive zero, so the minimum of \(F\) on \([L,r_\rho^\sharp]\) occurs
at an endpoint.  Since \(R\ge n/2\ge64\),
\[
 F(L)\ge
 \frac{(R+1)(R+4)(R-8)-9n(2R-1)}{27}>0;
\]
indeed the numerator is at least
\(R\{R(R-4)-18n\}>0\).  Direct substitution at the other endpoint gives,
in the order of the four rates,
\[
 \frac{n^2-84n+288}{8},\quad
 \frac{3n^2-104n+128}{16},\quad
 \frac{21n^2-464n+512}{64},\quad
 \frac{105n^2-1952n+2048}{256},
\]
all positive for \(n\ge128\).  Therefore \(F(r)\ge0\) throughout the
remaining integer interval.  Since \(r>0\) and \(R-2r>0\) there,
\[
 r(r+2)(R-2r)-n(R-r)=F(r)+r(R-2r)>0,
\]
which is \eqref{eq:literature-integral-condition}.  Apply
\cref{thm:literature-integral-completion}.
\end{proof}

\begin{corollary}[Almost-half-distance target compiler]
\label{cor:half-distance-window-compiler}
In the setting of \cref{cor:official-rate-half-staircase}, let
\(0\le\eps\le1\) and put
\(B=\floor{\eps\abs\Gamma}\).  If
\[
 1\le B<\abs\Gamma,\qquad B\le r_\rho^\sharp+1,
\]
then the complete safe set is
\[
 \{\delta:e_{\rm MCA,\Gamma}(C,\delta)\le\eps\}
       =[0,B/n),                                   \tag{LC9}
 \label{eq:half-distance-safe-set}
\]
and the endpoint is unsafe.  For the full-field target \(2^{-128}\),
this applies throughout
\(2^{128}\le q<2^{128}(r_\rho^\sharp+2)\).
\end{corollary}

\begin{proof}
At every radius \(r\le B-1\),
\cref{cor:official-rate-half-staircase} gives numerator
\(r+1\le B\).  At radius \(B\), the tangent floor gives at least
\(B+1\) bad slopes.  Monotonicity and
\(r(\delta)=\floor{\delta n}\) prove the assertion.
\end{proof}


\section{Agreement-set shortening beyond Johnson}
\label{sec:agreement-shortening}

We record Jo's shortening mechanism because it preserves actual bad slopes
and therefore bypasses support-to-ray conversion.  The shortened evaluation
set need not retain any smoothness.

\begin{lemma}[Normalized shortening, Jo
  {\cite[Lem.~3.1]{Jo26}}]\label{lem:jo-shortening-map}
Let \(C=\RS_\F(D,k)\), let \(S\subseteq D\) have size \(t<k\), put
\(P_S(X)=\prod_{x\in S}(X-x)\), and let \(I_S(w)\) be the degree-less-than-
\(t\) polynomial interpolating \(w|_S\), with
\(P_\varnothing=1\) and \(I_\varnothing=0\).  Define, for
\(x\in D\setminus S\),
\[
 Q_S(w)(x)=\frac{w(x)-I_S(w)(x)}{P_S(x)}.          \tag{SH1}
\]
Then \(Q_S\) is linear and maps \(C\) into
\(C_S=\RS_\F(D\setminus S,k-t)\).  If a slope is MCA-bad for \(C\) on a
witness support \(A\supseteq S\), then it is MCA-bad for the shortened line
at agreement \(\abs A-t\) in \(C_S\).
\end{lemma}

\begin{proof}
We reproduce Jo's normalization argument.  For a code polynomial \(f\),
the polynomial \(f-I_S(f)\) vanishes on
\(S\), so division by \(P_S\) leaves degree less than \(k-t\).  If the
shortened received pair had explanations \(g_i\) on \(A\setminus S\), then
\(I_S(u_i)+P_Sg_i\) would explain each original word on all of \(A\),
contradicting the witness nontriviality.
\end{proof}

\begin{theorem}[Agreement-set shortening transfer, Jo
  {\cite[Thm.~3.2]{Jo26}}]\label{thm:jo-shortening-transfer}
Let \(C=\RS_\F(D,k)\) have length \(n\), let \(k+1\le a\le n\), and let
\(0\le t<k\) with \(t\le a\).  For every nonempty challenge set
\(\Gamma\subseteq\F\),
\[
 B_{C,\Gamma}^{\rm MCA}(a)
 \le \frac{\binom nt}{\binom at}
       \max_{\substack{S\subseteq D\\\abs S=t}}
       B_{C_S,\Gamma}^{\rm MCA}(a-t).              \tag{SH2}
 \label{eq:jo-shortening-transfer}
\]
\end{theorem}

\begin{proof}
This is Jo's double count, restricted to \(\Gamma\).  For a fixed received
line, choose a witness support \(A_\gamma\) of size at
least \(a\) for every bad \(\gamma\in\Gamma\).  Each \(A_\gamma\) contains
at least \(\binom at\) shortening sets, and
\cref{lem:jo-shortening-map} preserves the same slope for every such set.
Count pairs \((\gamma,S)\) first by slope and then over the \(\binom nt\)
choices of \(S\), and finally maximize over the received line.
\end{proof}

\begin{corollary}[Fixed integer steps beyond Johnson, Jo
  {\cite[Thm.~3.4]{Jo26}}]\label{cor:jo-fixed-johnson-steps}
Fix integers \(r\ge2\) and \(h\ge1\).  For all sufficiently large \(k\),
let \(n=rk\), let \(D\subseteq\F\) be arbitrary of size \(n\), and put
\[
 E_{r,h}=\floor{n-\sqrt{n(k-1)}}+h,\qquad
 t_{r,h}=\floor{\frac{2h\sqrt r}{(\sqrt r-1)^2}}+1. \tag{SH3}
\]
Then \(E_{r,h}\) is the \(h\)-th integer error budget strictly beyond the
exact Johnson boundary, and
\[
 B_C^{\rm MCA}(n-E_{r,h})=O_{r,h}(k^6).            \tag{SH4}
\]
The assertion is uniform over the field and evaluation set.
\end{corollary}

The key identity behind the corollary is
\[
 (a-t)^2-(n-t)(k-1-t)
 =a^2-n(k-1)+t(n+k-1-2a).                         \tag{SH5}
\]
The displayed \(t_{r,h}\) makes the shortened side strictly Johnson-safe;
the resulting normalized margin is \(\Omega_{r,h}(k^{-1})\), and the
BCHKS estimate is \(O_{r,h}(k^6)\).

\begin{proof}[Proof of \cref{thm:intro-shortening-exponent}]
Fix \(\tau\) with \(\tau_\rho(\delta)<\tau<\rho\) and put
\[
 t_n=\floor{\tau n},\quad E_n=\floor{\delta n},\quad
 N_n=n-t_n,\quad d_n=k_n-1-t_n,\quad
 A_n=a_n-t_n=N_n-E_n.
\]
Because \(\tau<\rho<\alpha\), while \(k_n/n\to\rho\) and
\(a_n/n\to\alpha\), eventually
\[
             0\le t_n\le\min\{a_n,k_n-2\}.          \tag{SH6}
\]
Thus both the shortening transfer and
\cref{thm:bchks-johnson-general} apply.  Moreover,
\[
 (\alpha-\tau)^2-(1-\tau)(\rho-\tau)
 =\alpha^2-\rho+\tau(1+\rho-2\alpha)>0,           \tag{SH9}
 \label{eq:shortened-density-surplus}
\]
so \(A_n^2>N_nd_n\) for all sufficiently large \(n\).  The limiting
reduced rate \((\rho-\tau)/(1-\tau)\) is positive, and the strict surplus
is bounded away from zero after normalization.  Consequently
\(\Jcal(N_n,d_n,E_n)=O_{\rho,\delta,\tau}(n)\), uniformly over the
arbitrary shortened domains \(D_n\setminus S\).
\Cref{thm:bchks-johnson-general} therefore makes the maximum shortened
numerator in
\eqref{eq:jo-shortening-transfer} polynomial in \(n\).  Stirling's formula
therefore gives
\[
 \limsup_{n\to\infty}\frac1n\log_2
 \bigl(1+B_{C_n,\Gamma_n}^{\rm MCA}(a_n)\bigr)
 \le H_2(\tau)-\alpha H_2(\tau/\alpha).
\]
Let \(\tau\downarrow\tau_\rho(\delta)\).  This order of limits is
essential: shortening at the critical density itself has no fixed Johnson
margin.

Solving equality in \eqref{eq:shortened-density-surplus} gives
\(\tau_\rho(\delta)\).  Indeed \(\rho<\alpha<\sqrt\rho\), so
\[
 1+\rho-2\alpha>0,\qquad
 \rho-\tau_\rho(\delta)
 =\frac{(\alpha-\rho)^2}{1+\rho-2\alpha}>0;
\]
hence \(0<\tau_\rho(\delta)<\rho\).  For fixed \(\alpha\),
\[
 \frac{d}{dx}\left[H_2(x)-\alpha H_2(x/\alpha)\right]
 =\log_2\frac{1-x}{\alpha-x}>0.
\]
Comparison at \(x=\tau_\rho(\delta)\) and \(x=\rho\) proves
\eqref{eq:shortening-beats-circuit}; direct substitution at the two
endpoints proves the final
claims.
\end{proof}

\begin{remark}[What shortening does not remove]
Splitting a total shortening \(t=t_1+t_2\) into stages gives the same cost,
because the two binomial ratios telescope to
\(\binom nt/\binom at\).  Thus constant relative improvement has the
positive exponent \(\Psi_\rho(\delta)\); the method does not by itself
produce a subexponential numerator beyond Johnson.
\end{remark}

\begin{proposition}[Four finite rows beyond Johnson]
\label{prop:paving-finite-rows}
Let
\[
 p=(2^{128}-255)2^{128}+1,\qquad
 B=\floor{p2^{-128}}=2^{128}-255,
\]
and let \(D\le\F_p^\times\) have order \(512\).  The number \(p\) is
prime.  At target \(2^{-128}\), agreement-set shortening and the
circuit-incidence envelope certify
\[
\begin{array}{c|c|c|c|c|c}
 \rho&k&a&(512-a)/512&\text{certificate}&\log_2 U<\\ \hline
 1/2&256&351&161/512&\text{shorten }123&126.234492\\
 1/4&128&233&279/512&\text{shorten }67&127.290843\\
 1/8&64 &143&369/512&\text{shorten }45&126.254265\\
 1/16&32&52&115/128&\text{circuit }b=33&127.310201.
\end{array}                                       \tag{F512}
 \label{eq:paving-finite-rows}
\]
In every row the rigorous numerator \(U\) is strictly smaller than \(B\),
so all four radii are safe and strictly beyond Johnson.  The decimal
logarithms are diagnostics only; the proof uses exact integer inequalities.
These are upper certificates, not claims that the MCA numerator attains
\(U\), and failure one agreement lower would not prove unsafety.
\end{proposition}

\begin{proof}
Write \(p=u2^{128}+1\) with odd \(u=2^{128}-255<2^{128}\).
The exact congruence
\(3^{(p-1)/2}\equiv-1\pmod p\) is a Proth certificate.  Indeed every
prime divisor \(\ell\) of \(p\) satisfies
\(\ell\ge2^{128}+1\), whereas
\(p<2^{256}+1<(2^{128}+1)^2\), so a composite \(p\) could not have a
prime divisor at most \(\sqrt p\).  Thus \(p\) is prime, and
\(512\mid p-1\) gives the subgroup.

For the first three rows put
\(E=512-a\), \(N=512-t\), \(d=k-1-t\), \(A=a-t\), and
\(x=m+1/2\), using
\[
 (t,m)=(123,162),(67,63),(45,15),\qquad
 C=\frac23x^5N^4+xdN^2(E+1).
\]
The exact inequalities
\[
 A^2>Nd,\qquad (m-1)^2A^2<dm^2N,\qquad
 d(m+1)^2N\le m^2A^2
\]
certify both the shortened Johnson condition and the ceiling parameter in
\cref{thm:bchks-johnson-general}.  Combining that theorem with
\cref{thm:jo-shortening-transfer} gives
\[
 U=\frac{\binom{512}{t}}{\binom at}\,
       \frac{C}{\sqrt{(dN)^3}}.
\]
After squaring positive quantities, direct integer arithmetic gives
\[
 \left(\frac{\binom{512}{t}}{\binom at}C\right)^2
 <B^2(dN)^3
\]
in all three cases, hence \(U<B\).  In the last row,
\cref{thm:intro-paving} with \(b=33\) gives the exact integer upper bound
\(\floor{\binom{512}{33}/\binom{51}{32}}\)
\(=210954686508560867421211382134708972291<B\).
Finally, the four exact inequalities \(a^2<512(k-1)\), whose positive
margins are respectively \(7359,10735,11807,13168\), prove that the four
displayed radii exceed the exact MDS Johnson radii
\(1-\sqrt{(k-1)/512}\), and hence also \(1-\sqrt\rho\).
\end{proof}

Jo's own challenge-scale post-Johnson rows and exact length-\(64\)
transitions are recorded in \cite[Thms.~5.1--5.2]{Jo26}; they are not
duplicated here.  The rows in \cref{prop:paving-finite-rows} instead compare
the shortening and circuit mechanisms on the common length-\(512\) field
already used elsewhere in this paper.

\section{Reproducible finite Prize instances}\label{sec:prize-rows}

\begin{table}[H]
\centering
\caption{Explicit maximal-dimension prize-admissible rows.  In each row
\(k=2^{40}\), \(B=\floor{p/2^{128}}=r_\rho(n)+1\), and \((s,a_0)\) is part of
the Proth certificate in \cref{prop:proth-row-check}.}
\label{tab:prize-proth-rows}
\small
\begin{tabular}{@{}ccrrc@{}}
\hline
rate & length & \(B\) & bits of \(p\) & \((s,a_0)\)\\
\hline
\(1/2\)  & \(2^{41}\) & \(389500552609\)  & 167 & \((92,3)\)\\
\(1/4\)  & \(2^{42}\) & \(1210584858040\) & 169 & \((93,13)\)\\
\(1/8\)  & \(2^{43}\) & \(2879806199253\) & 170 & \((95,5)\)\\
\(1/16\) & \(2^{44}\) & \(6233898019554\) & 171 & \((97,5)\)\\
\hline
\end{tabular}
\end{table}

For completeness, the four field orders are
\begin{align*}
p_{1/2}&=132540169958804033333249306710494641010898987122689,\\
p_{1/4}&=411940680852499481698306614369841346700408394874881,\\
p_{1/8}&=979947269755402568812854322316630667196565607677953,\\
p_{1/16}&=2121285573237585848299875619011192262679065433997313.
\end{align*}
Their odd Proth coefficients \(u\), in the same order, are
\begin{align*}
u_{1/2}&=26766274163673319604503,
&u_{1/4}&=41595378994516821279015,\\
u_{1/8}&=24737346889219389259839,
&u_{1/16}&=13387194060291799253121.
\end{align*}
For \(F_{n,k}(r)=r^2-n(3r-(n-k))\), the exact boundary checks are
\[
\begin{array}{c|r|r}
 \rho&F_{n,k}(B-1)&F_{n,k}(B)\\ \hline
 1/2&5154112775168&-663955886271\\
 1/4&7590647904465&-3182321912768\\
 1/8&13908181940112&-6720484728007\\
 1/16&19335616403905&-20973145690236.
\end{array}                                                    \tag{PC0}
\label{eq:quadratic-row-certificates}
\]

\begin{proposition}[Exact arithmetic and primality certificates]
\label{prop:proth-row-check}
For the four rows in \cref{tab:prize-proth-rows}, the displayed integer
\(p\) satisfies
\[
 p=u\,2^s+1,\qquad u\text{ odd},\qquad u<2^s,\qquad
 a_0^{(p-1)/2}\equiv-1\pmod p.                  \tag{PC1}
\label{eq:proth-certificates}
\]
Consequently every \(p\) is prime.  Moreover
\[
 n\mid p-1,\qquad p<2^{256},\qquad
 B\,2^{128}\le p<(B+1)2^{128},                  \tag{PC2}
\label{eq:prize-row-arithmetic}
\]
where \(n\) and \(B\) are the corresponding entries of the table.
They also satisfy \(F_{n,k}(B-1)\ge0>F_{n,k}(B)\).  Since
\(F_{n,k}'(r)=2r-3n<0\) on \([0,n]\), these two signs locate its smaller
root and give \(B-1=r_\rho(n)=r_{\rm quad}(n,k)\).
Thus \(\F_p^*\) contains a subgroup of order \(n\), and every row in
\cref{tab:prize-proth-rows} exists with the asserted slope budget.
\end{proposition}

\begin{proof}
All identities and inequalities in
\eqref{eq:quadratic-row-certificates}--\eqref{eq:prize-row-arithmetic}
are direct
integer calculations; the values of \((s,a_0)\), \(u\), \(p\), \(n\),
and \(B\) have all been printed above so that the calculation is
independently reproducible.  We recall why the certificate proves
primality.  If a prime \(\ell\) divides \(p\), then
\(a_0^{u2^{s-1}}\equiv-1\pmod\ell\), whereas
\(a_0^{u2^s}\equiv1\pmod\ell\).  The multiplicative order of \(a_0\)
modulo \(\ell\) therefore has exact \(2\)-part \(2^s\), so
\(2^s\mid\ell-1\).  Since \(u<2^s\),
\(p<2^{2s}+1<(2^s+1)^2\), so \(2^s+1>\sqrt p\).  A
composite \(p\) would have a prime divisor at most \(\sqrt p\), whereas
every prime divisor is at least \(2^s+1\), a contradiction.  This is Proth's theorem in the
special form used here.  Finally \(s\ge\log_2 n\) gives \(n\mid p-1\),
and cyclicity of \(\F_p^*\) gives the required subgroup.
\end{proof}

\begin{theorem}[Reed--Solomon endpoint, Jo
  {\cite[Thm.~4.3]{Jo26}}]\label{thm:jo-rs-endpoint}
Let \(C=\RS_{\F_q}(D,k)\), put \(M=\binom n{k+1}\), and use the
full-field challenge.  For every \(0\le a\le k+1\),
\[
 \ceil{\frac{Mq}{q+M-1}}
 \le B_C^{\rm MCA}(a)\le\min\{M,q\}.              \tag{EP}
\]
If \(q>\binom M2\), then
\[
                         B_C^{\rm MCA}(a)=M
                         \qquad(0\le a\le k+1).    \tag{EP$'$}
\]
\end{theorem}

\begin{corollary}[Exact first adjacent row]
\label{thm:exact-first-adjacent-row}
Let \(C=\RS_{\F_q}(D,k)\), assume \(k\le n-2\), and put
\[
 M=\binom n{k+1},\qquad U=\floor{\frac{M}{k+1}}.
\]
Suppose \(q>\binom M2\), use the full-field challenge, let
\(0\le\eps\le1\), put \(b=\floor{\eps q}\), and, when the set is nonempty,
define
\[
 a^*=\min\{a\in\{k+1,\ldots,n\}:B_C^{\rm MCA}(a)\le b\}.
\]
In either of the following cases \(a^*\) is defined, with
\[
 \begin{array}{ll}
 b\ge M &\Longrightarrow a^*=k+1,\\[1mm]
 U\le b<M&\Longrightarrow a^*=k+2.
 \end{array}                                      \tag{AD}
\]
In the first branch the complete official safe set is \([0,1]\); in the
second branch the transition occurs before capacity.
\end{corollary}

\begin{proof}
By \cref{thm:jo-rs-endpoint}, the numerator at \(k+1\) is exactly \(M\).
At \(a=k+2\), \cref{thm:intro-paving} with test size \(k+1\) gives
\(B_C^{\rm MCA}(k+2)\le\floor{M/(k+1)}=U\).  Monotonicity proves both
claims.
\end{proof}

\begin{theorem}[Subfield confinement]\label{thm:subfield-confinement-full}
Let \(\B\subseteq\F\), \(D\subseteq\B\), and let \(u_0,u_1\in\B^D\).  If
\(\gamma\in\F\setminus\B\) and \(u_0+\gamma u_1\) is explained on \(S\),
then the pair is explained on the same \(S\).  Hence every MCA-bad slope of
a \(\B\)-valued received line lies in \(\B\).
\end{theorem}

\begin{proof}
Choose a \(\B\)-basis \(1,\beta_1,\ldots\) of \(\F\).  Decompose the
coefficients of the explaining polynomial and of \(\gamma\) in this basis.
Because \(D\subseteq\B\), comparison of basis components on \(S\) gives a
codeword explaining \(u_1\) from any nonzero nonbase component of
\(\gamma\), and then a codeword explaining \(u_0\) from the base component.
\end{proof}

\section{Exact locator prefixes and collision-aware pole lines}
\label{sec:exact-prefix-pole}

Assume \(D\subseteq\B\subseteq\F\).  For an \(m\)-set \(S\subseteq D\),
write
\[
 Q_S(X)=\prod_{x\in S}(X-x)
       =X^m+c_1(S)X^{m-1}+\cdots+c_m(S),
 \qquad
 \Phi_w(S)=(c_1(S),\ldots,c_w(S)).
\]

The monic polynomial \(Q_S\), whose simple roots are precisely the points
of \(S\), is the \emph{support locator}.  Its \emph{depth-\(w\) locator
prefix} is the vector of the first \(w\) coefficients below the leading
coefficient.  For \(z\in\B^w\), the inverse image
\(\Phi_w^{-1}(z)\) is the \emph{prefix fiber over \(z\)}: the family of
all \(m\)-element supports having that same prefix.  Thus ``fiber'' in this
paper always refers to the inverse image of a displayed map.

\begin{proposition}[Exact prefix-list bijection]\label{prop:exact-prefix-list}
Let \(1\le K\le m\le n\), put \(w=m-K\), and set
\(U_z(X)=X^m+\sum_{i=1}^wz_iX^{m-i}\) for \(z\in\B^w\).  The codewords of
\(\RS_\F(D,K)\) agreeing with \(U_z|_D\) on at least \(m\) points are
exactly
\[
        (U_z-Q_S)|_D,\qquad S\in\Phi_w^{-1}(z).
 \tag{4.1}\label{eq:prefix-bijection}
\]
No such codeword agrees on more than \(m\) points.  In particular, some
received word has list size at least
\(\lceil\binom nm\abs\B^{-w}\rceil\).
\end{proposition}

\begin{proof}
For \(S\in\Phi_w^{-1}(z)\), the leading term and next \(w\) coefficients
cancel, so \(\deg(U_z-Q_S)<m-w=K\).  Conversely, if \(P\), \(\deg P<K\),
agrees with \(U_z\) at \(m\) points, then \(U_z-P\) is monic of degree
\(m\) with those \(m\) roots and equals their locator.  Coefficient
comparison gives the prefix.  A nonzero degree-\(m\) polynomial has no
additional root.  Pigeonhole the \(\binom nm\) supports over
\(\abs\B^w\) prefixes.
\end{proof}

The \emph{simple-pole conversion} chooses an auxiliary point
\(\alpha\notin D\) and divides coordinatewise by \(X-\alpha\).  This
turns one list of nearby polynomials into points on a received affine line
for the dimension-\(k\) code.  A \emph{collision} is an equality
\(P_i(\alpha)=P_j(\alpha)\) for two distinct listed polynomials; the next
formula pays such equalities rather than assuming all slopes are distinct.

\begin{theorem}[Collision-aware simple-pole conversion]
\label{thm:collision-aware-pole}
Let \(C=\RS_\F(D,k)\), \(C^+=\RS_\F(D,k+1)\), \(1\le k<n\), and
\(q=\abs\F>n\).  If a word has \(L\) distinct codewords of \(C^+\)
agreeing on at least \(m\ge k+1\) points, then one received line for \(C\)
has at least
\[
 M(L)=\left\lceil
 \frac{L(q-n)}{q-n+k(L-1)}
 \right\rceil
 \tag{4.2}\label{eq:collision-aware-pole}
\]
distinct MCA-bad slopes at agreement \(m\).
\end{theorem}

\begin{proof}
Let the received word be \(U\in\F^D\), and let the listed polynomials be
\(P_1,\ldots,P_L\), each of degree at most \(k\).  For
\(\alpha\in\F\setminus D\), define the word coordinatewise by
\(f_\alpha(x)=U(x)/(x-\alpha)\) and set
\(g_\alpha(x)=-1/(x-\alpha)\).  On an agreement set for \(P_i\), the line
point of slope \(P_i(\alpha)\) equals
\((P_i(x)-P_i(\alpha))/(x-\alpha)\), a degree-less-than-\(k\) evaluation.
The pair is not simultaneously explained on \(k+1\) points: an explanation
of \(g_\alpha\) by \(G\), \(\deg G<k\), would make
\((X-\alpha)G(X)+1\) vanish at \(k+1\) points although its degree is at
most \(k\) and its value at \(\alpha\) is \(1\).

For \(i\ne j\), the equality \(P_i(\alpha)=P_j(\alpha)\) holds at at most
\(k\) poles.  Some pole has at most \(k\binom L2/(q-n)\) colliding pairs.
If its value multiplicities are \(m_1,\ldots,m_M\), then
\(\sum m_i^2\le L+kL(L-1)/(q-n)\).  Cauchy--Schwarz gives
\(L^2\le M\sum m_i^2\), which yields \eqref{eq:collision-aware-pole}.
\end{proof}

\begin{corollary}[Exact good-pole reservoir]
\label{cor:exact-good-pole-reservoir}
In \cref{thm:collision-aware-pole}, let \(Q=q-n\), and for
\(\alpha\in\F\setminus D\) let \(C_\alpha\) be the number of unordered
pairs \(\{i,j\}\) with \(P_i(\alpha)=P_j(\alpha)\).  For every integer
\(T\ge0\), at least
\[
 \max\left\{0,Q-
   \floor{\frac{k\binom L2}{T+1}}\right\}          \tag{PR1}
\label{eq:good-pole-count}
\]
poles satisfy \(C_\alpha\le T\), and every such pole produces at least
\[
                     \ceil{\frac{L^2}{L+2T}}       \tag{PR2}
\label{eq:good-pole-slopes}
\]
distinct MCA-bad slopes.  In particular at least
\(\max\{0,Q-k\binom L2\}\) poles are fully separating.
\end{corollary}

\begin{proof}
Every nonzero difference \(P_i-P_j\) has at most \(k\) roots, so
\[
                 \sum_{\alpha\in\F\setminus D}C_\alpha
                    \le k\binom L2.
\]
Every pole with \(C_\alpha>T\) contributes at least \(T+1\) to this
sum, which gives \eqref{eq:good-pole-count}.  At a good
pole let \(m_1,\ldots,m_M\) be the multiplicities of the distinct list
values.  Then
\(\sum_im_i=L\) and
\(\sum_im_i^2=L+2C_\alpha\le L+2T\).  Cauchy--Schwarz yields
\(M\ge L^2/(L+2T)\), proving \eqref{eq:good-pole-slopes}; \(T=0\)
gives the last assertion.
\end{proof}

\begin{remark}[One pole, one received line]
The reservoir in \cref{cor:exact-good-pole-reservoir} fixes the same pole
when counting collisions, because different poles define different
received lines.  A cross-pole equation
\(P_i(\alpha)=P_j(\alpha')\) is one equation in two variables and
generically defines a curve; a Bezout degree bound does not turn it into a
finite collision set.  Such cross-line collisions are neither controlled
by the simple-pole argument nor needed by it.
\end{remark}

\begin{corollary}[Identity-prefix MCA floor]\label{cor:identity-prefix-floor}
Let \(C=\RS_\F(D,k)\), where
\(D\subseteq\B\subseteq\F\), \(\abs D=n\), and
\(q=\abs\F>n\).  For \(k+1\le m\le n\), put \(w=m-k-1\) and
\(L_m=\lceil\binom nm\abs\B^{-w}\rceil\).  Then
\[
 B_C^{\rm MCA}(m)\ge M(L_m).
 \tag{4.3}\label{eq:identity-prefix-floor}
\]
\end{corollary}


\begin{proposition}[Two deployed unsafe certificates]
\label{prop:deployed-prefix-attacks}
Put \(n=2^{21}\), \(k=2^{20}\), and define \(L_a\) and \(M(L_a)\) by
\cref{cor:identity-prefix-floor,eq:collision-aware-pole}.  The following
full-field rows have exact integer attack evaluations:
\[
\begin{array}{c|c|c|c|c}
\text{row and target}&a_0&B^*&M(L_{a_0})&M(L_{a_0+1})\\ \hline
p_{\rm KB}^6,\ 2^{-128}&1116047&
274980728111395087&138634741058327852652&57198030366\\
(2^{31}-1)^4,\ 2^{-100}&1116023&
16777215&4281388998575706&1752700 .
\end{array}                                       \tag{DP1}
 \label{eq:deployed-prefix-attacks}
\]
Here \(p_{\rm KB}=2^{31}-2^{24}+1\); the first domain is its
multiplicative subgroup of order \(n\), and the second is a
Mersenne--31 Chebyshev \(x\)-domain of order \(n\): the image of a disjoint
twin coset in a norm-one torus under \(u\mapsto(u+u^{-1})/2\), as formalized
in \cref{def:circle-twin-domain}.
Thus the first row is unsafe at radius \(981105/2097152\), by
\(8.9777\) bits beyond the 128-bit budget, and the second is unsafe at
radius \(981129/2097152\), by \(27.9270\) bits beyond its meaningful
100-bit budget.  The last column says only that this particular attack
falls below budget one agreement later; it is not a safety certificate.
\end{proposition}

\begin{proof}
For each row set \(w=a-k-1\) and evaluate
\[
 L_a=\ceil{\binom na p^{-w}},\qquad
 M(L_a)=\ceil{\frac{L_a(q-n)}{q-n+k(L_a-1)}}
\]
with \((p,q)=(p_{\rm KB},p_{\rm KB}^6)\) and
\((2^{31}-1,(2^{31}-1)^4)\), respectively.  Exact integer arithmetic
gives \eqref{eq:deployed-prefix-attacks}; the first two values exceed
their exact budgets.  \Cref{cor:identity-prefix-floor} proves unsafety,
and the radius and bit statements are the displayed integer differences
and base-two ratios.
\end{proof}

For equal-size sets, the \emph{Johnson distance} is
\(d_J(S,T)=\abs{S\setminus T}=\abs{T\setminus S}\).  It counts the
number of point replacements required to turn one support into the other.

\begin{proposition}[Prefix rigidity and its packing limit]
\label{prop:prefix-rigidity-full}
Two distinct \(m\)-sets in one depth-\(w\) prefix fiber have Johnson
distance at least \(w+1\).  Consequently, with
\(t=\floor{w/2}\), every fiber has size at most
\[
 \frac{\binom nm}{\sum_{i=0}^t\binom mi\binom{n-m}{i}}.
 \tag{4.4}\label{eq:packing-fiber-cap}
\]
At \(w\asymp n/\log\abs\B\), this packing loss is only
\(e^{o(n)}\) and does not supply the random factor \(\abs\B^{-w}\).
\end{proposition}

Prefix rigidity says that sharing many leading locator coefficients forces
distinct supports to be far apart in this metric.  Here the
\emph{polynomial-field regime} means \(\abs\B=n^{O(1)}\), equivalently
\(\log\abs\B=O(\log n)\).

\begin{proof}
If \(S,T\) have Johnson distance \(j\le w\), factor the locator of their
intersection, of degree \(m-j\), from \(Q_S-Q_T\).  The residual factor is
nonzero, so the difference has degree at least \(m-j\ge m-w\), whereas the
common prefix makes its degree at most \(m-w-1\), a contradiction.  Radius
\(t\) Johnson balls about the fiber members are disjoint and each has the
denominator in \eqref{eq:packing-fiber-cap}.  For
\(w=n/\log\abs\B\), the logarithm of that ball volume is
\(O(w\log(n/w))=o(n)\) in the polynomial-field regime.
\end{proof}

\begin{theorem}[Prefix fibers embed in received lines]
\label{thm:prefix-to-line-hardness}
Let \(\mathcal G\subseteq\Phi_w^{-1}(z)\) be a fiber of \(N\) many
\(m\)-sets, with \(0\le w\le m-2\), and put \(k=m-w-1\).  Every finite
extension \(\F/\B\) satisfying
\[
        \abs\F-n>k\binom N2
        \tag{4.5}\label{eq:extension-separation}
\]
admits one received line for \(\RS_\F(D,k)\) with at least \(N\) distinct
MCA-bad slopes furnished by \(\mathcal G\).
\end{theorem}

\begin{proof}
For distinct \(S,T\in\mathcal G\), the polynomial \(Q_S-Q_T\) is nonzero
of degree at most \(k\).  Condition \eqref{eq:extension-separation} leaves a
pole \(\alpha\in\F\setminus D\) avoiding the roots of all these
differences.  Thus the values \(Q_S(\alpha)\) are pairwise distinct.  By
\cref{prop:exact-prefix-list}, the polynomials \(U_z-Q_S\) form one list for
the dimension-\(k+1\) code; the pole line makes their values, and hence the
resulting MCA slopes, pairwise distinct.
\end{proof}

\begin{theorem}[Exact list--line bijection at a separating pole]
\label{thm:exact-list-line-bijection}
Let \(D\subseteq\B\subseteq\F\), let \(1\le k<n\), let \(m\ge k+1\),
and let \(U\in\B^D\).  Write
\[
 \Lcal_m(U)=\{P\in\B[X]:\deg P\le k,
       \ \abs{\{x\in D:U(x)=P(x)\}}\ge m\}.
\]
Suppose \(\alpha\in\F\setminus D\) and the evaluation map
\(P\mapsto P(\alpha)\) is injective on \(\Lcal_m(U)\).  Define
\[
 f_\alpha(x)=\frac{U(x)}{x-\alpha},\qquad
 g_\alpha(x)=-\frac1{x-\alpha}\qquad(x\in D).
\]
Then \(P\mapsto P(\alpha)\) is a bijection from \(\Lcal_m(U)\) onto
the set of finite MCA-bad slopes of the received line
\((f_\alpha,g_\alpha)\) for \(\RS_\F(D,k)\) at agreement \(m\).  The
corresponding explaining polynomial is
\[
       h_P(X)=\frac{P(X)-P(\alpha)}{X-\alpha},
\]
and its agreement set with
\(f_\alpha+P(\alpha)g_\alpha\) is exactly the agreement set of \(P\)
with \(U\).

If \(L=\abs{\Lcal_m(U)}\), such a pole exists in every finite extension
satisfying
\[
       \abs\F>n+k\binom L2.                         \tag{4.6}
\]
\end{theorem}

\begin{proof}
For \(P\in\Lcal_m(U)\), the polynomial \(h_P\) has degree less than
\(k\), and on every point where \(U=P\) one has
\[
 f_\alpha+P(\alpha)g_\alpha
   =\frac{P-P(\alpha)}{X-\alpha}=h_P.
\]
The word \(g_\alpha\) cannot be explained by a polynomial \(G\) of degree
less than \(k\) on \(k+1\) points, since then
\((X-\alpha)G+1\) would have degree at most \(k\), at least \(k+1\)
roots, and value one at \(\alpha\).  Thus every displayed slope is
MCA-bad.

Conversely, suppose \(f_\alpha+\gamma g_\alpha\) is explained by
\(h\in\F[X]_{<k}\) on at least \(m\) points.  Then
\(P(X)=(X-\alpha)h(X)+\gamma\) has degree at most \(k\) and agrees with
\(U\) there.  Interpolation on any \(k+1\) of those points uses a
Vandermonde system over \(\B\) with right-hand side in \(\B\), so
\(P\in\B[X]\).  Hence \(P\in\Lcal_m(U)\) and
\(\gamma=P(\alpha)\).  Injectivity proves the bijection, and the same
calculation proves equality of the agreement sets.

For distinct \(P,Q\in\Lcal_m(U)\), the nonzero polynomial \(P-Q\) has
degree at most \(k\), and therefore at most \(k\) roots.  The union of
\(D\) and the root sets of all \(\binom L2\) differences has size at most
\(n+k\binom L2\); a larger extension contains an admissible pole.
\end{proof}

\begin{corollary}[Exact prefix-fiber ray realization]
\label{cor:exact-prefix-ray-realization}
Let \(0\le w\le m-2\), put \(k=m-w-1\), and let
\(\Fcal_z=\Phi_w^{-1}(z)\) be a complete depth-\(w\) fiber of
\(m\)-subsets of \(D\).  If \(\alpha\) separates the values
\(Q_S(\alpha)\), \(S\in\Fcal_z\), then the pole line associated with
\(U_z=X^m+\sum_{i=1}^wz_iX^{m-i}\) has exactly
\(\abs{\Fcal_z}\) MCA-bad slopes at agreement \(m\).  The maps
\[
 S\longmapsto U_z(\alpha)-Q_S(\alpha)
 \longmapsto
 \left(U_z(\alpha)-Q_S(\alpha),S,
 \frac{U_z-Q_S-(U_z(\alpha)-Q_S(\alpha))}{X-\alpha}\right)
\]
are bijections between the prefix fiber, the bad-slope set, and the
exact-\(m\) witness incidence.  In particular, the complete prefix fiber,
its bad-slope set, and its exact-agreement witness incidence have the same
cardinality.
\end{corollary}

\begin{proof}
The exact prefix-list bijection identifies \(\Fcal_z\) with the complete
dimension-\((k+1)\) list for \(U_z\).  Apply
\cref{thm:exact-list-line-bijection}.  Since
\(U_z-(U_z-Q_S)=Q_S\) has precisely the roots in \(S\), every agreement
is exactly \(m\) and supplies exactly that support.
\end{proof}



\section{Further exact constructions}\label{sec:further-exact}

The main threshold theorem needs only the syndrome correspondence and the
mean-overlap argument.  This section records complementary exact
constructions that clarify the support-to-slope projection and the
structured domains used in the examples.  Every statement is a finite
identity, inequality, or construction; no equidistribution or inverse
theorem is assumed.

\subsection{Determinantal ray normalization}\label{drc:sec-determinantal-ray}

The final projection from supports to slopes can be made explicit without
postulating a support-pair multiplicity.  The construction below applies to
every Reed--Solomon row, in every characteristic.  It also identifies the
exact additional boundary maps that would govern extensions beyond the
proved overlap range.  No estimate for those maps is used in the main
threshold theorem.

The support counts above are stated for agreement supports
\(S\in\binom Da\), whereas this subsection uses the complementary error
supports \(E=D\setminus S\in\binom D{n-a}\).  Complementation is a
bijection.  If the domain is partitioned into fibers of a folding map, a
fiber containing \(j\) agreement positions contains its remaining positions
as errors; exceptional positions are complemented as well.  Thus every
fiber-occupancy census transfers exactly to the error-support formulation.

Let \(C=\RS_\F(D,k)\), put \(R=n-k\), and use the weighted Vandermonde
parity-check columns
\[
 h_x=\lambda_x(1,x,\ldots,x^{R-1})^{\mathsf T},
 \qquad \lambda_x\ne0,
\]
as in the syndrome--secant normal form.  For \(E\subseteq D\), write
\(V_E=\Span\{h_x:x\in E\}\).  Fix an exact agreement \(a\ge k+1\), and set
\[
             t=n-a,\qquad d=R-t=a-k.
             \tag{DR0}\label{drc:eq-depth}
\]
In particular, \(d-1=a-k-1\) is precisely the identity locator-prefix
depth.  For \(E\in\binom Dt\), write its error locator in ascending order as
\[
 Q_E(X)=\prod_{x\in E}(X-x)=\sum_{\ell=0}^{t}q_\ell(E)X^\ell,
 \qquad q_t(E)=1.
\]
For \(y=(\eta_0,\ldots,\eta_{R-1})\in\F^R\), define the
\emph{locator-recurrence boundary}
\[
 \mathcal A_E(y)=
 \left(\sum_{\ell=0}^{t}q_\ell(E)\eta_{j+\ell}\right)_{0\le j<d}
 \in\F^d.
 \tag{DR1}\label{drc:eq-recurrence-boundary}
\]
If \(f(X)=\sum_{i=0}^{R-1}f_iX^i\), we use the coefficient pairing
\(
 \langle y,f\rangle=\sum_{i=0}^{R-1}y_if_i
\).

\begin{lemma}[Locator recurrence and a secant span]
\label{drc:lem-locator-recurrence}
For every \(E\in\binom Dt\) and \(y\in\F^R\),
\[
                         y\in V_E
             \quad\Longleftrightarrow\quad
                         \mathcal A_E(y)=0.
             \tag{DR2}\label{drc:eq-recurrence-membership}
\]
\end{lemma}

\begin{proof}
Identify a polynomial of degree less than \(R\) with its coefficient
vector in \(\F^R\).  Since
\(
 \langle h_x,f\rangle=\lambda_xf(x)
\), every coefficient vector of a polynomial in
\(Q_E\F[X]_{<d}\) annihilates \(V_E\).  The MDS column property gives
\(\dim V_E=t\), while multiplication by the nonzero monic polynomial
\(Q_E\) gives
\(\dim Q_E\F[X]_{<d}=d=R-t\).  Consequently
\[
 V_E^\perp=
 \{\operatorname{coeff}(Q_EP):P\in\F[X]_{<d}\}.
\]
Testing against the basis \(1,X,\ldots,X^{d-1}\) of
\(\F[X]_{<d}\) gives exactly the \(d\) coordinates in
\eqref{drc:eq-recurrence-boundary}.
\end{proof}

For a syndrome pair \(y_0,y_1\in\F^R\), and a support family
\(\Omega\subseteq\binom Dt\), define its support-restricted finite bad-slope
set by
\[
 Z_\Omega(y_0,y_1)=
 \left\{\gamma\in\F:\begin{array}{l}
   \text{there is \(E\in\Omega\) with }
       y_0+\gamma y_1\in V_E,\\[-1mm]
   \{y_0,y_1\}\not\subseteq V_E
 \end{array}\right\}.
 \tag{DR3}\label{drc:eq-restricted-rays}
\]
For \(E\in\Omega\), abbreviate
\[
              u_E=\mathcal A_E(y_0),\qquad
              v_E=\mathcal A_E(y_1).
\]
Let \(\mathcal P_j(y_1)\) be the universal pivot locus on which \(j\) is
the first nonzero coordinate of \(v_E\), and put
\(\Omega_j=\Omega\cap\mathcal P_j(y_1)\):
\[
 \begin{aligned}
 \mathcal P_j(y_1)
   &=\{E\in\tbinom Dt:
          v_{E,i}=0\ \text{for }0\le i<j,\quad v_{E,j}\ne0\},\\
 \Omega_j&=\Omega\cap\mathcal P_j(y_1),
       &&0\le j<d.
 \end{aligned}
 \tag{DR4}\label{drc:eq-pivot-chart}
\]
On this chart define the \emph{determinantal ray boundary}
\[
 \Delta_j(E)=
 \bigl(v_{E,j}u_{E,\ell}-u_{E,j}v_{E,\ell}\bigr)_
       {\substack{0\le\ell<d\\ \ell\ne j}}
 \in\F^{d-1}.
 \tag{DR5}\label{drc:eq-pivot-boundary}
\]
When \(d=1\), the target in \eqref{drc:eq-pivot-boundary} is the trivial
group and the displayed tuple is empty.

\begin{theorem}[Exact determinantal ray compiler]
\label{drc:thm-exact-compiler}
With the notation above, an error support \(E\in\Omega\) supplies a finite
transverse slope if and only if
\[
 v_E\ne0,
 \qquad
 \rank\begin{pmatrix}u_E\\v_E\end{pmatrix}\le1.
 \tag{DR6}\label{drc:eq-rank-one-condition}
\]
On the pivot chart \(\Omega_j\), this is equivalent to
\(\Delta_j(E)=0\), and its unique slope is
\(
 \gamma_E=-u_{E,j}/v_{E,j}
\).  Consequently
\[
 Z_\Omega(y_0,y_1)=
 \bigcup_{j=0}^{d-1}
 \left\{-\frac{u_{E,j}}{v_{E,j}}:
       E\in\Omega_j,\ \Delta_j(E)=0\right\},
 \tag{DR7}\label{drc:eq-exact-ray-union}
\]
and
\[
 \abs{Z_\Omega(y_0,y_1)}
 \le
 \sum_{j=0}^{d-1}
 \abs{\Omega_j\cap\Delta_j^{-1}(0)}.
 \tag{DRC}\label{drc:eq-direct-compiler}
\]
The construction uses \(d\le n\) pivot charts and contains no unnamed
finite or asymptotic multiplicity.
\end{theorem}

\begin{proof}
By \cref{drc:lem-locator-recurrence}, \(E\) supports the line parameter
\(\gamma\) precisely when
\[
                         u_E+\gamma v_E=0.
                         \tag{DR8}\label{drc:eq-vector-slope}
\]
If \(v_E=0\), a solution of \eqref{drc:eq-vector-slope} forces \(u_E=0\),
so \(y_0,y_1\in V_E\) and the intersection is common rather than
transverse.  If \(v_E\ne0\), equation
\eqref{drc:eq-vector-slope} has a solution if and only if \(u_E\) is a
scalar multiple of \(v_E\), which is
\eqref{drc:eq-rank-one-condition}; the scalar, and therefore the finite
slope, is unique.  The charts \(\Omega_j\) partition the locus
\(v_E\ne0\).  On \(\Omega_j\), vanishing of all \(2\)-by-\(2\) minors
using the pivot column \(j\) is equivalent to proportionality, and the
proportionality factor is \(-u_{E,j}/v_{E,j}\).  This proves
\eqref{drc:eq-exact-ray-union}.  Taking the cardinality of a union and then
forgetting the slope projection gives \eqref{drc:eq-direct-compiler}.
\end{proof}

\begin{proposition}[Exact slope-incidence linearization]
\label{drc:prop-exact-slope-linearization}
With the notation of \cref{drc:thm-exact-compiler}, put
\[
 \Ical_\Omega=
 \{(\gamma,E)\in\F\times\Omega:
                 u_E+\gamma v_E=0,\ v_E\ne0\}.
 \tag{DRC1}\label{drc:eq-slope-incidence}
\]
Projection \(\Ical_\Omega\to\Omega\) is injective, its slope projection
is \(Z_\Omega(y_0,y_1)\), and
\[
 \abs{Z_\Omega(y_0,y_1)}
 \le\abs{\Ical_\Omega}
 =\sum_{\gamma\in\F}
   \abs{\{E\in\Omega:u_E+\gamma v_E=0,\ v_E\ne0\}}. \tag{DRC2}
\label{drc:eq-slope-sum}
\]
Thus a fixed-slope zero-fiber estimate must be aggregated over all slopes.
If every fixed-slope zero fiber were bounded by \(\abs\Omega/q^d\),
summing over the \(q\) slopes would give \(\abs\Omega/q^{d-1}\), not
\(\abs\Omega/q^d\).
\end{proposition}

\begin{proof}
For \(v_E\ne0\), the vector equation determines at most one
\(\gamma\), proving injectivity.  The exact determinantal compiler says
that its slope image is precisely the transverse bad-slope set.  Counting
the incidence first by \(E\) and then by \(\gamma\) gives
\eqref{drc:eq-slope-sum}; the final scale is the corresponding uniform
fiber calculation.
\end{proof}

\begin{proposition}[The recurrence boundary is not pointwise additive]
\label{drc:prop-recurrence-nonadditive}
The exact proportionality condition above does not, by itself, turn the
locator-recurrence boundary into a subset-sum map.  Already for
\(E=\{x,y\}\), one coordinate of \(\mathcal A_E(\eta)\) is
\[
 F(x,y)=\eta_jxy-\eta_{j+1}(x+y)+\eta_{j+2}.       \tag{DRC3}
\label{drc:eq-mixed-boundary}
\]
If \(\eta_j\ne0\), this function cannot be written as
\(c+g(x)+g(y)\) whenever the four admissible two-sets
\(\{x,y\},\{x',y\},\{x,y'\},\{x',y'\}\) lie in the chart and
\((x-x')(y-y')\ne0\).
Consequently, a proposed Fourier reduction to additive maps
\(E\mapsto\sum_{t\in E}g_\gamma(t)\) requires a separate coordinate
theorem; it is not a consequence of slope linearization.
\end{proposition}

\begin{proof}
For a two-set,
\(Q_E(X)=X^2-(x+y)X+xy\), which gives
\eqref{drc:eq-mixed-boundary}.  Its mixed difference is
\[
 F(x,y)-F(x',y)-F(x,y')+F(x',y')
   =\eta_j(x-x')(y-y').
\]
Every function \(c+g(x)+g(y)\) has zero mixed difference, proving the
claim whenever the displayed product is nonzero.
\end{proof}

The maps in \eqref{drc:eq-pivot-boundary} are explicit quadratic maps in
the locator coefficients.  They depend on the received syndrome pair, as
they must in a worst-line theorem, but their number and degree are literal.
The next corollary shows how they replace the abstract ray interface.


\begin{theorem}[Fixed-core determinacy ray compiler]
\label{thm:fixed-core-determinacy-ray}
Let \(C_0\subseteq D\) have size \(c\le k\), let \(a\ge k+1\), and
fix a received line \(r=(r_0,r_1)\).  Suppose that for every
\(\gamma\) in a set \(Z\) of distinct finite MCA-bad slopes we have
chosen an exact-\(a\) noncommon witness
\((\gamma,S_\gamma,h_\gamma)\) with \(C_0\subseteq S_\gamma\).  Then
\[
             \abs Z\le \binom{n-c}{k-c+1}.          \tag{RC$_{\rm core}$}
\]
In particular, a fixed-core cell with \(k-c=o(n)\) has a direct
subexponential ray payment.  This statement is uniform in the received
line and valid in every characteristic.
\end{theorem}

\begin{proof}
Put \(d=k-c\), let
\(Q_{C_0}(X)=\prod_{x\in C_0}(X-x)\), and choose
\(g_0,g_1\in\F[X]_{<k}\) interpolating \(r_0,r_1\), respectively, on
\(C_0\).  Write \(e_i=r_i-\operatorname{ev}_D(g_i)\).  Since
\(h_\gamma-g_0-\gamma g_1\) vanishes on \(C_0\), there is a unique
\(P_\gamma\in\F[X]_{<d}\) such that
\[
 h_\gamma=g_0+\gamma g_1+Q_{C_0}P_\gamma .
\]
For \(x\in D\setminus C_0\), the witness equation is
\[
 e_0(x)+\gamma e_1(x)-Q_{C_0}(x)P_\gamma(x)=0.     \tag{*}
\]
Homogenize the parameters as
\([z:\gamma:p_0:\cdots:p_{d-1}]\in\Pj^{d+1}\), and let the row
indexed by \(x\in D\setminus C_0\) be the linear functional
\[
 z e_0(x)+\gamma e_1(x)
   -Q_{C_0}(x)\sum_{j=0}^{d-1}p_jx^j.
\]
The rows indexed by \(S_\gamma\setminus C_0\) have rank exactly
\(d+1\).  Indeed, the displayed witness parameter is a nonzero kernel
vector, so the rank is at most \(d+1\).  If the rank were at most \(d\),
the kernel would contain an independent vector.  Subtracting a multiple
of the witness vector gives a nonzero kernel vector with \(z=0\).  Its
\(\gamma\)-coordinate cannot vanish: otherwise a nonzero polynomial of
degree less than \(d\) would vanish at all
\(a-c\ge d+1\) points of \(S_\gamma\setminus C_0\).  Adding a nonzero
multiple of this kernel vector to the witness vector therefore produces a
second finite slope explained on the same support.  By fixed-support slope
uniqueness that support would be common, contrary to the witness
hypothesis.

Choose, in a fixed order, the first \((d+1)\)-subset
\(T_\gamma\subseteq S_\gamma\setminus C_0\) on which the rows have rank
\(d+1\).  Their common projective kernel is the single witness parameter,
so \(T_\gamma\) determines \(\gamma\).  Hence
\(\gamma\mapsto T_\gamma\) is injective into
\(\binom{D\setminus C_0}{d+1}\), proving the bound.  If
\(d=o(n)\), then
\(\log\binom{n-c}{d+1}=o(n)\).
\end{proof}

\begin{remark}[Cores are internal unless their census is small]
The theorem pays one fixed core.  It does not permit an atlas to use every
positive-density subset \(C_0\) as a separate profile: that census can be
exponential.  A complete atlas either routes only an algebraically
specified subexponential family of cores or keeps the actual core internal
to a coarse intersection-size profile.
\end{remark}

\subsection{Sharpness of support-only ray bounds}
A support \(S\) in a multiplicative group \(D\) is
\emph{multiplicatively aperiodic} when its stabilizer
\(\{g\in D:gS=S\}\) is trivial.
\begin{theorem}[Aperiodic ray obstruction below the profile frontier]
\label{thm:aperiodic-ray-obstruction}
There is a sequence of smooth multiplicative rows
\[
 D_n=\B_n^\times\subseteq\F_n^\times,
 \qquad \log\abs{\B_n}=O(\log n)=o(n),
\]
in characteristic two, with rates tending to \(1/2\), and agreements
\(m_n=k_n+O(n/\log n)\), for which one received line has
\(2^{\Omega(n)}\) distinct finite MCA-bad slopes.  Every threshold-
\(m_n\) witness support of that line has trivial multiplicative
stabilizer, and after enlarging \(\F_n\) with only
\(\log\abs{\F_n}=O(n)\), the bad slopes and exact-agreement supports are
in bijection.  Thus smoothness, quotient-aperiodicity, or one-parameter
support-profile data,
and exact agreement do not force retained-support occupancy greater than
one.  This is a sharpness result for occupancy-based compilers, not a
counterexample to a direct distinct-slope theorem at the correctly
normalized profile scale.
\end{theorem}

\begin{proof}
Let \(q_r=2^r\), put \(\B_r=\F_{q_r}\),
\(D_r=\B_r^\times\), and \(n_r=q_r-1\).  Set
\[
 m_r=2^{r-1}-1=\frac{n_r-1}{2}.
\]
Fix \(0<\eta<\log2\), take
\(w_r=\floor{\eta n_r/\log q_r}\), and put
\(k_r=m_r-w_r-1\).  Then \(k_r/n_r\to1/2\) and
\(m_r-k_r=w_r+1=O(n_r/\log n_r)\).

Pigeonholing the \(m_r\)-set locators over their \(q_r^{w_r}\) prefixes
gives a fiber \(\Fcal_r\) satisfying
\[
 \log\abs{\Fcal_r}
 \ge \log\binom{n_r}{m_r}-w_r\log q_r
 \ge(\log2-\eta)n_r-o(n_r).
\]
Moreover
\[
 \gcd(m_r,n_r)=
 \gcd(2^{r-1}-1,2^r-1)=2^{\gcd(r-1,r)}-1=1.
\]
If an \(m_r\)-support were invariant under a nontrivial subgroup of
\(D_r\) of order \(c>1\), it would be a union of \(c\)-element orbits,
forcing \(c\mid m_r\) and \(c\mid n_r\), a contradiction.  Thus every
member of \(\Fcal_r\) is aperiodic.

Choose a finite extension \(\F_r/\B_r\) with
\[
 \abs{\F_r}>n_r+k_r\binom{\abs{\Fcal_r}}2.
\]
Since \(\abs{\Fcal_r}\le2^{n_r}\), the least extension degree meeting
this inequality has \(\log\abs{\F_r}=O(n_r)\).  By
\cref{cor:exact-prefix-ray-realization}, a separating pole gives exactly
\(\abs{\Fcal_r}=\exp(\Omega(n_r))\) distinct bad slopes, one for each
aperiodic support, and every occupancy equals one.
\end{proof}

\begin{remark}[Scope of the aperiodic obstruction]
The construction lies below the entropy--profile crossing: its natural
prefix scale is itself exponential.  It therefore does not contradict a
profile-envelope upper bound proved at the correct scale.  It does rule
out a field-independent polynomial ray bound, division of a one-parameter
support count by the full support-slice size without a retained multiplicity
theorem, and the assertion
that quotient-aperiodicity alone forces ray collisions.
\end{remark}

The opposite projection extreme occurs over the coefficient field itself.
It shows that a large aperiodic prefix fiber need not supply a lower
reserve after first-match saturation.

\begin{theorem}[Aperiodic one-ray saturation]
\label{thm:aperiodic-one-ray-saturation}
For an integer \(\nu\ge5\), put
\[
 q=2^\nu,\quad D=\F_q^\times,\quad n=q-1,\quad
 m=2^{\nu-1}-1,\quad w=\floor{n/\nu^2},\quad k=m-w-1.
\]
Refine the depth-\(w\) prefix by the locator constant coefficient:
\[
 \Xi_\nu(S)=\bigl(c_1(S),\ldots,c_w(S),c_m(S)\bigr)
 \qquad\left(S\in\binom Dm\right).
\]
There are \(z\in\F_q^w\) and \(c\in\F_q^\times\) for which
\[
 \Gcal_{z,c}=\{S:\Xi_\nu(S)=(z,c)\},
 \qquad
 \abs{\Gcal_{z,c}}
 \ge\ceil{\binom nm q^{-(w+1)}}
 =\exp\bigl((\log2-o(1))n\bigr).                  \tag{SAT1}
\label{eq:one-ray-large-fiber}
\]
Every support in \(\Gcal_{z,c}\) has trivial multiplicative stabilizer,
but all witnesses arising from the exact prefix list project at the pole
\(\alpha=0\) in the coefficient field \(\F_q\) to the single MCA slope \(-c\) for
\(\RS_{\F_q}(D,k)\).

More precisely, if \(\Fcal_z=\Phi_w^{-1}(z)\) and
\(\Ccal_d=\{S\in\Fcal_z:c_m(S)=d\}\), then the nonempty
\(\Ccal_d\), \(d\in\F_q^\times\), partition the supports of all
exact-\(m\) witnesses of this pole line, and the witnesses supported on
\(\Ccal_d\) have slope image exactly \(\{-d\}\).  Thus at most
\(q-1=\exp(o(n))\) explicit saturation cells
pay the complete line even though one cell contains exponentially many
aperiodic supports.
\end{theorem}

\begin{proof}
The map \(\Xi_\nu\) has at most \(q^{w+1}\) values, so pigeonholing
proves the first inequality in \eqref{eq:one-ray-large-fiber}.  Moreover
\[
 \log\binom nm=n\log2-O(\log n),\qquad
 (w+1)\log q\le(n/\nu+\nu)\log2=o(n),
\]
which proves its asymptotic form.  Also
\[
 \gcd(m,n)=\gcd(2^{\nu-1}-1,2^\nu-1)
           =2^{\gcd(\nu-1,\nu)}-1=1.
\]
An invariant support would be a union of orbits of a nontrivial subgroup
whose order divides both \(m\) and \(n\), so every \(m\)-support is
multiplicatively aperiodic.

Let \(U_z=X^m+\sum_{i=1}^wz_iX^{m-i}\) and
\(P_S=U_z-Q_S\).  Prefix cancellation gives \(\deg P_S\le k\), and
\(U_z-P_S=Q_S\) shows that \(P_S\) agrees with \(U_z\) exactly on \(S\).
At the only pole of the coefficient field outside \(D\), namely zero,
\[
                     P_S(0)=U_z(0)-Q_S(0)=-c_m(S).
\]
The pole-line calculation therefore maps every \(\Ccal_d\) to \(-d\).
Conversely, an exact-\(m\) explanation on this line gives a polynomial
\(P=Xh+\gamma\) of degree at most \(k\) agreeing with \(U_z\) on
\(m\) points.  The exact prefix-list bijection makes it the unique
\(P_S=U_z-Q_S\) for some \(S\in\Fcal_z\), and then
\(\gamma=P(0)=-c_m(S)\).  This proves coverage and the asserted
one-slope images.
\end{proof}

\begin{remark}[Two incompatible aperiodic inferences]
\label{rem:two-aperiodic-ray-extremes}
\Cref{thm:aperiodic-ray-obstruction} makes a separating extension pole
inject an exponentially large aperiodic fiber into distinct slopes, while
\cref{thm:aperiodic-one-ray-saturation} makes a coefficient-field pole collapse
such a refined fiber to one slope.  Aperiodicity alone therefore implies
neither ray expansion nor ray collision.  A lower compiler needs a
separating-pole estimate or a quantitative exclusion of earlier saturation
cells; an upper compiler needs an actual determinant or ray-image bound.
\end{remark}

\subsection{Smooth and circle evaluation domains}\label{sec:smooth-circle-domains}

\paragraph{Characteristic two.}
The support reduction, tangent construction, MDS paving bound, and
all-radius saturation theorem use only linear algebra and the MDS property,
and therefore hold over every finite field.  The domain transports require
more care: \(\abs{\F_{2^s}^{\times}}=2^s-1\) is odd, so characteristic two
has no nontrivial power-of-two multiplicative subgroup, while
\(x^2+y^2-1=(x+y-1)^2\) is singular, so the circle equivalences used here
are necessarily odd-characteristic statements.  Characteristic-two analogues
based on additive FFT domains or a nonsingular conic would require separate
agreement-preserving maps.  Thus the representation-free MDS results apply
to characteristic-two Reed--Solomon rows, but no multiplicative-subgroup or
circle transfer in that characteristic is asserted.

This subsection records the domain geometry behind the exact quotient
constructions.  The multiplicative and circle cases are treated simultaneously through
uniform-fiber folding maps.  The point is not merely that such maps exist,
but that their fibers are complete inside the evaluation domain and their
locators have a \emph{lacunary} top part, meaning that prescribed gaps occur
among the highest-degree coefficients.  These two facts make quotient
families visible in the witness moduli and keep their locator prefixes
stable under planted cores.

\begin{definition}[Complete-fiber folding map and smooth multiplicative coset]
\label{def:structured-folding}
Let \(D\subseteq\B\) be finite.  A polynomial map
\(\varphi:D\to Q:=\varphi(D)\) is a \emph{\(c\)-fold
complete-fiber folding map} if every \(y\in Q\) has exactly \(c\)
preimages in \(D\).  A support is \emph{complete at this folding scale}
if it is a union of whole fibers, equivalently if it is
\(\varphi^{-1}(E)\) for some \(E\subseteq Q\).  Passing from \(D\) to
\(Q\), and from a complete support to \(E\), is the corresponding
\emph{quotient} or \emph{folding}.

A \emph{multiplicative coset} is a set \(D=\theta H\subseteq\B^\times\),
where \(H\le\B^\times\) is a subgroup and \(\theta\in\B^\times\).
For every divisor \(c\mid\abs H\), the restriction
\[
             \pi_c:D\longrightarrow D^{(c)},\qquad
             \pi_c(x)=x^c,\qquad
             D^{(c)}=\{x^c:x\in D\},
\]
is a \(c\)-fold complete-fiber folding map.  We call \(D\)
\emph{smooth at scale \(c\)} when this power folding is among the
quotients retained by the row, and \emph{smooth} for a family of scales
when it is smooth at every scale in that family.  Thus the adjective
\emph{smooth} here is an explicit uniform-fiber property, not a claim
about nonsingularity of an algebraic variety.  The locator of a complete
support is \(V_E(X^c)\); its nonleading terms occur only in degree gaps
divisible by \(c\).  This controlled pattern of missing top coefficients
is what \emph{lacunary} means in this section.
\end{definition}

\begin{definition}[Circle domain and circle code]
\label{def:circle-domain-code}
Let \(\B_0\) have odd characteristic and let
\(\mathcal C/\B_0\) be the affine conic
\(\mathcal C:x^2+y^2=1\).  A \emph{standard circle domain} is a finite
set \(E\subseteq\mathcal C(\B_0)\).  If \(\F\supseteq\B_0\), the
degree-\(w\) \emph{circle code} on \(E\) is
\[
 \mathcal C_w(\F,E)
 =\{(f(P))_{P\in E}:
      f\in\operatorname{im}(\F[x,y]_{\le w}
          \to\F[x,y]/(x^2+y^2-1))\}.
\]
The degree is the total-degree filtration before passage to the quotient
ring.  A \emph{circle \(x\)-domain} is instead a set of first
coordinates obtained from the torus construction below; an RS code on
that \(x\)-domain and a bivariate circle code are different presentations
until a diagonal equivalence such as
\cref{lem:circle-uniformization} has been exhibited.
\end{definition}

\begin{definition}[Norm-one torus, twin coset, and Chebyshev domain]
\label{def:circle-twin-domain}
Following the circle-domain geometry of \cite{HLP24}, let
\(p\equiv3\pmod4\), choose \(i\in\F_{p^2}\) with \(i^2=-1\), and
write
\[
 \UU=\ker\!\left(N_{\F_{p^2}/\F_p}:\F_{p^2}^{\times}
                      \to\F_p^\times\right)
     =\{u\in\F_{p^2}^{\times}:u^{p+1}=1\}.
\]
The algebraic kernel of the norm morphism is the
\emph{norm-one torus}; \(\UU\) is its cyclic group of
\(\F_p\)-rational points.  Put
\(\chi(u)=(u+u^{-1})/2\in\F_p\).  If \(H\le\UU\) and
\(g\in\UU\) satisfies \(g^2\notin H\), then
\[
 \Dcal(g,H)=gH\sqcup g^{-1}H
\]
is a disjoint union exchanged by \(u\mapsto u^{-1}\); it is called a
\emph{twin coset}.  Its projected set
\[
             D(g,H)=\chi(\Dcal(g,H))\subseteq\F_p
\]
has cardinality \(\abs H\) and is the \emph{Chebyshev twin-coset
\(x\)-domain}.

For \(c\ge1\), let \(T_c\) be the Chebyshev polynomial normalized by
\[
       T_c\!\left(\frac{u+u^{-1}}2\right)
       =\frac{u^c+u^{-c}}2 .
\]
Thus \(T_c\circ\chi=\chi\circ(u\mapsto u^c)\).  Whenever
\(c\mid\abs H\) and the two image cosets remain disjoint, the restriction
of \(T_c\) to \(D(g,H)\) is the \emph{Chebyshev folding map}; a
Chebyshev twin-coset domain equipped with these complete-fiber maps is
the circle analogue of a smooth multiplicative coset.
\end{definition}

\begin{definition}[Map-smooth domain]\label{def:map-smooth}
Let \(\B\subseteq\F\) be finite fields, let \(\varphi\in\B[X]\) have
degree \(c\ge1\), and let \(D\subseteq\B\) have size \(n\).  The set
\(D\) is \((\varphi,c)\)-smooth if every nonempty fiber
\(\varphi^{-1}(y)\cap D\), \(y\in\varphi(D)\), has exactly \(c\)
elements.  We write \(Q=\varphi(D)\) and \(N=n/c\).
\end{definition}

A multiplicative coset \(D=\theta H\subseteq\B^\times\) is \((X^a,a)\)-smooth for every divisor \(a\mid\abs H\).  The circle case uses Chebyshev maps rather than power maps; see \cref{lem:cheb-smooth}.  The following lemma is the common quotient-support construction.  It is included because it produces explicit large lists and because the same lacunarity controls their locator prefixes.

\begin{lemma}[Graded complete-fiber prefix]
\label{lem:map-smooth-fiber}
Let \(D\) be \((\varphi,c)\)-smooth, put \(Q=\varphi(D)\) and
\(N=\abs Q=n/c\), and let \(C^+=\RS_\F(D,K)\).  For an integer
\[
 d=\ceil{K/c}\le m\le N,\qquad w=m-d,
\]
there is a word \(u\in\B^D\) having at least
\[
                         L=\ceil{\binom Nm\abs\B^{-w}}       \tag{MS1}
 \label{eq:graded-fiber-list}
\]
distinct codewords of \(C^+\) that agree with it on \(mc\) coordinates.
\end{lemma}

\begin{proof}
For each \(m\)-set \(E\subseteq Q\), write
\[
 P_E(Y)=\prod_{b\in E}(Y-b)
 =Y^m+\sum_{j=1}^{w}z_j(E)Y^{m-j}+R_E(Y),
 \qquad \deg R_E\le d-1.
\]
Pigeonholing the \(\binom Nm\) sets by their vector
\((z_1(E),\ldots,z_w(E))\in\B^w\) gives a fiber of size at least \(L\);
fix its value \(z\).  The word
\(u_z=\varphi^m+\sum_{j=1}^wz_j\varphi^{m-j}\) agrees on the \(mc\)
points above \(E\) with the codeword \(-R_E(\varphi)\).  This is a
codeword because
\(\deg R_E(\varphi)\le c(d-1)<K\).  Distinct \(E\) give distinct
codewords: each composed remainder has degree less than \(K\le n\), so
equality on the \(n\)-point domain first implies equality as polynomials.
Composition with the nonconstant \(\varphi\) is injective, and the fixed
leading prefix then gives
\(P_E=P_{E'}\).
\end{proof}

\begin{proposition}[Collision-aware map-smooth attack]
\label{prop:map-smooth-cap}
Let \(D\) be \((\varphi,c)\)-smooth, let \(q=\abs\F>n\), and put
\(C=\RS_\F(D,k)\).  For
\[
 d=\ceil{(k+1)/c}\le m\le N,\quad w=m-d,\quad
 L=\ceil{\binom Nm\abs\B^{-w}},
\]
every agreement \(k+1\le A\le mc\) satisfies
\[
 B_C^{\rm MCA}(A)\ge
 \left\lceil\frac{L(q-n)}{q-n+k(L-1)}\right\rceil.          \tag{MS2}
 \label{eq:map-smooth-attack}
\]
In particular, take \(m=d+1\).  If
\[
                         \binom Nm\ge\abs\B\,(q/k+1),       \tag{MS3}
 \label{eq:map-smooth-cap-condition}
\]
then at every official radius \(1-mc/n\le\delta\le1\),
\[
 e_{\rm MCA}(C,\delta)>
 \frac1{2k}\left(1-\frac nq\right).                         \tag{MS4}
 \label{eq:map-smooth-security-cap}
\]
\end{proposition}

\begin{proof}
Apply \cref{lem:map-smooth-fiber} with \(K=k+1\), then apply the
collision-aware pole conversion \cref{thm:collision-aware-pole}.  Its
received pair is not simultaneously explained on any \(k+1\) points, so
the same \(mc\)-agreement list witnesses every displayed threshold.
For the last assertion, \eqref{eq:map-smooth-cap-condition} gives
\(L\ge q/k+1\).  With \(x=q-n\),
\[
 \frac{Lx}{x+k(L-1)}>\frac{x}{2k},
\]
because the cross-multiplied inequality is \(kL+k>x\), which follows
from \(kL+k\ge q+2k>q-n\).  Divide by \(q\).
\end{proof}

We now verify the complete-fiber property for the construction in
\cref{def:circle-twin-domain}.  Thus
\(\Dcal=gH\sqcup g^{-1}H\subseteq\UU\),
\(D=\chi(\Dcal)\), and \(T_a(\chi(u))=\chi(u^a)\).

\begin{lemma}[Chebyshev smoothness]\label{lem:cheb-smooth}
Let \(\Dcal=gH\cup g^{-1}H\subseteq\UU\) be a twin coset, let \(D=\chi(\Dcal)\), and let \(a\mid\abs H\).  If \(g^{2a}\notin H^{(a)}=\{h^a:h\in H\}\), then \(D\) is \((T_a,a)\)-smooth.  If \(\ord(g)\nmid2\abs H\), this holds for all \(a\mid\abs H\).
\end{lemma}

\begin{proof}
The map \(u\mapsto u^a\) has kernel of size \(a\) on \(\UU\), and that kernel lies in \(H\).  Hence it maps \(gH\) and \(g^{-1}H\) onto \(g^aH^{(a)}\) and \(g^{-a}H^{(a)}\), respectively, with fibers of size \(a\).  The hypothesis says these image cosets are disjoint, so the image is again a twin coset.  Since \(\chi\) is exactly two-to-one on both twin cosets and identifies only inverse points, the equation \(T_a(x)=y\) on \(D\) has exactly \(a\) solutions for each \(y\in T_a(D)\).  Taking \(a=\abs H\) shows that the all-scale condition is equivalent to \(g^{2\abs H}\ne1\), and the implication to every smaller divisor follows by raising to \(\abs H/a\).
\end{proof}

\begin{lemma}[Diagonal uniformization of circle codes]\label{lem:circle-uniformization}
Assume \(\F\supseteq\F_{p^2}\).  Let
\(E\subseteq\{(x,y)\in\F_p^2:x^2+y^2=1\}\), define
\(\iota(x,y)=x+iy\), put \(E'=\iota(E)\subseteq\UU\), and let
\(w\ge0\) satisfy \(2w+1<\abs E\).  Let
\(\Ccal_w(\F,E)\) be the circle code of
\cref{def:circle-domain-code}.  Under the coordinate bijection \(\iota\)
and the diagonal multiplier \(t_u=u^{-w}\), the code
\(\Ccal_w(\F,E)\) is \(t\circ\RS_\F(E',2w+1)\).
Consequently list sizes, ordinary correlated-agreement (CA) numerators,
and mutual-correlated-agreement (MCA) numerators are preserved.
\end{lemma}

\begin{proof}
The map \(u=x+iy\) identifies \(\F[x,y]/(x^2+y^2-1)\) with \(\F[u,u^{-1}]\).  The subspace of total degree at most \(w\) maps to \(u^{-w}\F[u]_{\le2w}\), and both spaces have dimension \(2w+1\).  Multiplication by the nonzero diagonal vector \((u^{-w})_{u\in E'}\) preserves coordinatewise equality and hence every agreement notion used here.  Therefore it preserves list, CA, and MCA counts exactly.
\end{proof}

The preceding uniformization is convenient when \(i\) belongs to the
challenge field.  The following form works over every odd-characteristic
field and will also be used for scalar extensions of odd degree.

\begin{lemma}[Invariance under monomial equivalence]
\label{lem:monomial-mca-invariance}
Let \(C\subseteq\F^D\), let \(\pi:D'\to D\) be a bijection, and let
\((d_x)_{x\in D'}\in(\F^\times)^{D'}\).  Define
\[
 C'=\{(d_xc_{\pi(x)})_{x\in D'}:c\in C\}.
\]
The same coordinate permutation and diagonal multiplication give a
bijection between received pairs for \(C\) and \(C'\), preserving the
bad-slope set on every support and for every challenge set.  In
particular \(B_{C,\Gamma}^{\rm MCA}(a)=B_{C',\Gamma}^{\rm MCA}(a)\).
\end{lemma}

\begin{proof}
For every word \(u\), codeword \(c\), and support \(S\), equality
\(u|_S=c|_S\) is equivalent after the displayed coordinate change to
equality on \(\pi^{-1}(S)\).  The transformation is \(\F\)-linear, so it
commutes with \(r_0+\gamma r_1\).  It therefore preserves both the
single-line-point explanation and simultaneous explanation of the pair
on the corresponding support, slope by slope.
\end{proof}

\begin{corollary}[Generalized Reed--Solomon rows]
\label{cor:grs-transfer}
Every generalized Reed--Solomon code obtained from \(\RS_\F(D,k)\) by
a coordinate permutation and nonzero column multipliers has exactly the
same CA and MCA numerators as that RS code.  In particular all exact
staircases and target compilers in this paper transfer verbatim.
\end{corollary}

\begin{proof}
This is \cref{lem:monomial-mca-invariance} applied to the defining
monomial equivalence.
\end{proof}

\begin{proposition}[Stereographic circle equivalence]
\label{prop:stereographic-circle-equivalence}
Let \(\F\) have odd characteristic, let
\(E\subseteq\{(x,y)\in\F^2:x^2+y^2=1\}\) avoid \((-1,0)\), and put
\[
 t=\frac{y}{1+x},\qquad T=t(E).
\]
If \(2w+1\le\abs E\), then
\[
 \Ccal_w(\F,E)
  =\diag\bigl((1+t^2)^{-w}\bigr)_{t\in T}
       \RS_\F(T,2w+1)                             \tag{SC1}
 \label{eq:stereographic-code-equivalence}
\]
after the coordinate bijection \(E\to T\).  Hence every list, CA, and
support-wise MCA bad-slope set is preserved exactly.
\end{proposition}

\begin{proof}
The inverse parametrization is
\[
 x=\frac{1-t^2}{1+t^2},\qquad
 y=\frac{2t}{1+t^2},
\]
and \(1+t^2=2/(1+x)\ne0\) on \(E\).  Modulo the circle equation, the
degree-at-most-\(w\) filtered space is
\[
 \{f_0(x)+yf_1(x):\deg f_0\le w,\ \deg f_1\le w-1\},
\]
with dimension \(2w+1\).  After substitution and multiplication by
\((1+t^2)^w\), every such function becomes a polynomial of degree at
most \(2w\).  The substitution is injective in the localized coordinate
ring and the dimensions agree, so all of \(\F[t]_{\le2w}\) occurs.
Equation \eqref{eq:stereographic-code-equivalence} follows, and
\cref{lem:monomial-mca-invariance} proves the last assertion.
\end{proof}

\begin{corollary}[Exact circle staircase]
\label{cor:exact-circle-staircase}
Let \(E,w\) be as in \cref{prop:stereographic-circle-equivalence}, put
\(n=\abs E\) and \(k=2w+1<n\).  For every nonempty
\(\Gamma\subseteq\F\) and every \(0\le r\le n-k-1\),
\[
 (n-r)^2\ge n(k+r)
 \quad\Longrightarrow\quad
 B_{\Ccal_w(\F,E),\Gamma}^{\rm MCA}(n-r)
      =\min\{\abs\Gamma,r+1\}.                  \tag{SC2}
\]
The same equality holds under the hypotheses of
\cref{thm:literature-integral-completion}.  Consequently standard circle
codes admitting such a projection chart have the positive-relative-radius staircase of
\cref{cor:positive-radius-density}.
\end{corollary}

\begin{proof}
Apply \cref{lem:monomial-mca-invariance} and then
\cref{thm:quadratic-mean-overlap,thm:literature-integral-completion} to
the RS presentation in
\eqref{eq:stereographic-code-equivalence}.
\end{proof}


\begin{proposition}[A 128-bit circle certificate]
\label{prop:paving-circle-certificate}
Let \(p_0=2^{127}-1\), \(q=p_0^2\), and work over \(\F_q\).  Then
\[
 \floor{q2^{-128}}
 =85070591730234615865843651857942052863.          \tag{CP1}
 \label{eq:circle-budget}
\]
There are two length-\(512\) circle presentations certified beyond
Johnson:
{\small
\[
\begin{array}{c|c|c|c|c}
\text{presentation}&k&a&(512-a)/512&
 U_{512,k,\F_q}(a)\\ \hline
\text{Chebyshev \(x\)-domain RS}&32&53&459/512&
81136417887908025931235146974888066266\\
\text{degree-\(15\) circle code}&31&50&462/512&
60827291905480389917403158602781524185 .
\end{array}                                       \tag{CP2}
 \label{eq:circle-paving-rows}
\]
}
Both upper numerators are at most \eqref{eq:circle-budget}, giving
\(128.068311\) and \(128.483942\) proved bits, respectively.
\end{proposition}

\begin{proof}
The Lucas--Lehmer criterion \cite[Sec.~4.2]{Riesel} proves that \(p_0\) is
prime, and its norm-one
torus \(\UU\) over \(\F_{p_0}\) has order \(p_0+1=2^{127}\).  Choose a
generator \(z\), put \(H=\langle z^{2^{118}}\rangle\), and set \(g=z\).
Then \(\abs H=512\) and \(g^2\notin H\).  For the first row use the
projection of the twin coset \(gH\sqcup g^{-1}H\) and
\cref{lem:cheb-smooth}; for the second take the \(512\) circle points
\(E=\iota^{-1}(gH)\), where \(\iota(x,y)=x+iy\), and use
\cref{lem:circle-uniformization}.  The two displayed
bounds are exactly
\(\floor{\binom{512}{33}/\binom{52}{32}}\) and
\(\floor{\binom{512}{32}/\binom{49}{31}}\).
Their comparison with the exact budget proves safety.  Both radii exceed
their Johnson radii, and direct base-two logarithms give the printed bits.
\end{proof}

For a set \(\Pcal\) of primes, its \emph{relative prime density} is
\[
\operatorname{dens}_{\rm pr}(\Pcal)=
 \lim_{X\to\infty}\frac{\#\{p\le X:p\in\Pcal\}}{\pi(X)},
\]
when the limit exists.  The prime number theorem in arithmetic
progressions gives density \(1/\varphi(M)\) to every reduced residue
class modulo \(M\) \cite{Davenport}.
We write \(\overline{\operatorname{dens}}_{\rm pr}\) for the same
expression with \(\limsup\) in place of the limit.

\begin{theorem}[Fixed-length positive-density prime families]
\label{thm:fixed-length-prime-density}
Fix \(n=2^m\), \(m\ge2\).
\begin{enumerate}[label=\textup{(\alph*)},leftmargin=*]
\item The primes \(p\equiv1\pmod n\), of relative prime density
\(1/\varphi(n)=2/n\), support the unique subgroup
\(H\le\F_p^\times\) of order \(n\).  Every exact RS conclusion of
\cref{thm:quadratic-mean-overlap,thm:literature-integral-completion}
holds on \(H\) and on each of its multiplicative cosets.
\item The primes \(p\equiv-1\pmod {2n}\), of density
\(1/\varphi(2n)=1/n\), support a standard circle coset \(E\) of size
\(n\).  For every odd \(k=2w+1<n\), the circle code
\(\Ccal_w(\F_{p^2},E)\) has both exact staircases.
\item The primes \(p\equiv-1\pmod {4n}\), of density
\(1/\varphi(4n)=1/(2n)\), support a Chebyshev-smooth circle
\(x\)-domain \(D\subseteq\F_p\) of size \(n\), and
\(\RS_{\F_p}(D,k)\) has both exact staircases for every \(k<n\).
\end{enumerate}
\end{theorem}

\begin{proof}
Part~\textup{(a)} follows from cyclicity of \(\F_p^\times\).  For
part~\textup{(b)}, let \(\UU\le\F_{p^2}^\times\) be the norm-one group,
of order \(p+1\), choose \(g\in\UU\) of order \(2n\), and put
\(H=\langle g^2\rangle\).  The coset \(E'=gH\) has size \(n\), is
stable under inversion, and avoids \(-1\).  Its inverse image under
\(u=x+iy\) consists of \(n\) points of the circle over \(\F_p\).
\Cref{lem:circle-uniformization} identifies its circle code with a
diagonal translate of \(\RS_{\F_{p^2}}(E',2w+1)\).

For part~\textup{(c)}, choose \(g\in\UU\) of order \(4n\), put
\(H=\langle g^4\rangle\), and use the twin coset
\(gH\sqcup g^{-1}H\).  Its \(\chi\)-image has size \(n\), and
\(\ord(g)=4n\nmid2n\), so \cref{lem:cheb-smooth} proves smoothness at
every dyadic scale.  The density assertions are the prime number theorem
in the three displayed reduced residue classes.
\end{proof}

\begin{corollary}[Finite smooth and circle families beyond the prior window]
\label{cor:circle-32-family}
Let \(n=32,k=15,r=6\).  For every prime
\(p\equiv1\pmod {32}\), let \(H\le\F_p^\times\) have order \(32\)
and put \(C_p^+=\RS_{\F_p}(H,15)\).  For every prime
\(p\equiv-1\pmod {64}\), construct the circle coset \(E\) of
\cref{thm:fixed-length-prime-density}\textup{(b)} and put
\(C_p^-=\Ccal_7(\F_{p^2},E)\).  Then
\[
 B_{C_p^+}^{\rm MCA}(26)=B_{C_p^-}^{\rm MCA}(26)=7. \tag{SC3}
\]
At targets \(6/p\) and \(6/p^2\), respectively, both complete safe
sets are \([0,3/16)\), with unsafe endpoint.  The two supporting prime
sets have relative densities \(1/16\) and \(1/32\).
\end{corollary}

\begin{proof}
Here \(R=17\) and
\[
 (n-r)^2=26^2=676>672=32(15+6).
\]
Thus \cref{thm:quadratic-mean-overlap,cor:exact-circle-staircase} give
the numerator seven for the smooth and circle rows.  Radius five is in
the deep range and has numerator six; the tangent floor at radius six and
monotonicity give the stated safe sets.  This example lies
beyond the elementary one-third range because \(3r=18>R\).  It is not
certified by the specialized BCHKS estimate
\eqref{eq:bchks-specialized}, whose exact first term is
\[
 \frac{32(18-6)}{6(18-12)}=\frac{32}{3}>7.
\]
The older conservative \(R\)-based comparator is \(176/15\).  The respective densities
are \(1/\varphi(32)=1/16\) and \(1/\varphi(64)=1/32\).
\end{proof}

\begin{corollary}[Concrete Mersenne-31 circle rows]
\label{cor:mersenne-circle-rows}
Let \(p=2^{31}-1\), the Mersenne prime used in \cite{HLP24}, and let
\(\UU\) be its norm-one circle group.
Choose a generator \(g\), put \(H_0=\langle g^2\rangle\), and let
\(E\subseteq\{(x,y)\in\F_p^2:x^2+y^2=1\}\) be the circle domain
corresponding to \(gH_0\), so \(\abs E=2^{30}\).  This domain avoids
\((-1,0)\).  Over \(\F_{p^5}\), the circle code with
\(w=2^{28}\) and \(k=2^{29}+1\) has, at target \(2^{-128}\),
the exact safe set
\[
\left[0,\frac18-2^{-30}\right).                 \tag{SC4}
\label{eq:mersenne-circle-safe}
\]
Moreover, if \(H=\langle g^4\rangle\), the Chebyshev twin-coset
\(gH\sqcup g^{-1}H=gH_0\) has \(x\)-domain
\(D=\chi(gH_0)\) of size \(2^{29}\), and
\(\RS_{\F_{p^5}}(D,2^{27})\) has exact safe set
\[
\left[0,\frac14-2^{-29}\right).                 \tag{SC5}
\label{eq:mersenne-x-safe}
\]
The first code has rate slightly above \(1/2\); the second has exact
rate \(1/4\) but is a circle-smooth \(x\)-domain rather than a literal
multiplicative-domain Prize row.
\end{corollary}

\begin{proof}
The group \(\UU\) has order \(p+1=2^{31}\).  Since
\(-1=g^{2^{30}}\in H_0\), neither \(1\) nor \(-1\) belongs to the
disjoint coset \(gH_0\); in particular \(E\) avoids the pole
\((-1,0)\) of the stereographic chart.  Also
\(gH_0=gH\sqcup g^{-1}H\).  Inversion acts freely on this coset and
\(\chi(u)=(u+u^{-1})/2\) identifies precisely the pairs
\(\{u,u^{-1}\}\), proving \(\abs D=2^{29}\).  The relation
\(T_a\circ\chi=\chi\circ(u\mapsto u^a)\) and
\cref{lem:cheb-smooth} give the asserted dyadic smoothness.

Binomial expansion gives
\(B=\floor{p^5/2^{128}}=2^{27}-1\).  At radius \(B-1\), the quadratic
margin \((n-r)^2-n(k+r)\) for the bivariate row is
\(162129591417176068>0\), so
\cref{cor:exact-circle-staircase} and the tangent floor give
\eqref{eq:mersenne-circle-safe}.  Here \(p^5\equiv3\pmod4\), so
\(i\notin\F_{p^5}\); the applicable equivalence is the stereographic
one in \cref{prop:stereographic-circle-equivalence}, not torus
uniformization.  For the \(x\)-domain row the same budget lies inside
the quadratic staircase at rate \(1/4\), since its margin at radius
\(B-1\) is \(18014401193836548>0\).  The identical target argument gives
\eqref{eq:mersenne-x-safe}.
\end{proof}

\begin{corollary}[Every odd-characteristic extension-field tower]
\label{cor:density-one-extension-towers}
Fix an odd prime \(p\), put \(n_e=2^e\) and \(q_e=p^{n_e}\), and assume
\(e\ge7\).  Then \(4n_e\mid q_e-1\).  Over \(\F_{q_e}\) there exist
both a multiplicative subgroup RS row of length \(n_e\) and a
Chebyshev-smooth circle \(x\)-domain row of that length.  At every
official rate and every \(0\le r\le r_\rho^\sharp\), both satisfy
\[
 e_{\rm MCA}(C_e,r/n_e)=\frac{r+1}{q_e}.           \tag{SC6}
\]
Thus the construction exists in every odd characteristic and beats every
fixed error target throughout every compact subinterval below half the
minimum distance for large \(e\).
\end{corollary}

\begin{proof}
The lifting-the-exponent identity gives
\[
 v_2(p^{2^e}-1)=v_2(p-1)+v_2(p+1)+e-1\ge e+2.
\]
Cyclicity produces the multiplicative subgroup and an element of order
\(4n_e\).  Since \(n_e\) is even, \(q_e\equiv1\pmod4\); choose
\(i\in\F_{q_e}\) with \(i^2=-1\).  The map
\[
 u\longmapsto
 \left(\frac{u+u^{-1}}2,\frac{u-u^{-1}}{2i}\right)
\]
identifies \(\F_{q_e}^\times\) with the split circle over
\(\F_{q_e}\).  If \(g\) has order \(4n_e\) and
\(H=\langle g^4\rangle\), the twin coset
\(gH\sqcup g^{-1}H\) projects to an \(n_e\)-point \(x\)-domain.
The proof of \cref{lem:cheb-smooth}, which uses only cyclicity and
\(T_a\circ\chi=\chi\circ(u\mapsto u^a)\), applies verbatim and proves
Chebyshev smoothness at every dyadic scale.
Apply \cref{cor:official-rate-half-staircase}.  Since
\(q_e=p^{n_e}\), the normalized error tends to zero exponentially.
\end{proof}

\begin{proposition}[No growing positive-density prime-field family]
\label{prop:no-growing-prime-density}
Let \(\Pcal\) be a set of primes and assign a dyadic integer
\(n_p\to\infty\) along \(\Pcal\).  If \(n_p\mid p-1\) for every
\(p\in\Pcal\), or \(n_p\mid p+1\) for every \(p\in\Pcal\), then
\(\Pcal\) has upper relative prime density zero.  If in addition
\(p/n_p\) is bounded, there are only \(O(\log X)\) possible primes of
the indicated form below \(X\).
\end{proposition}

\begin{proof}
For every fixed \(E\), all sufficiently large assigned lengths are
divisible by \(2^E\).  Hence the primes eventually lie in the single
class \(1\pmod {2^E}\), or the single class \(-1\pmod {2^E}\), whose
relative prime density is \(1/\varphi(2^E)=2^{1-E}\).  Letting
\(E\to\infty\) proves the first assertion.  Under the bounded-ratio
hypothesis, \(p\mp1=a2^e\) with \(a\) in a fixed finite set; there are
only \(O(\log X)\) choices of \((a,e)\) below \(X\).
\end{proof}

\begin{remark}[Prize scope]
The prime-density statements fix \(n\); their densities tend to zero as
\(n\) grows.  The extension towers have exponential field size and soon
leave the \(q<2^{256}\) Prize box.  Finally the target \(6/p^2\) in
\cref{cor:circle-32-family} varies with \(p\).  None of these facts may be
rephrased as a positive-density growing prime-field solution at the fixed
target \(2^{-128}\).
\end{remark}

\subsection{Quotient--remainder profiles and partial occupancy}
\label{sec:quotient-obstruction}

The quotient and remainder variables can be separated exactly at the
locator-prefix level.  This removes a spurious factor coming from summing
over marked remainder sets, but it also exposes a genuine scaled quotient-
coefficient-field term which prevents an unrestricted comparison with the
identity profile.

\begin{definition}[Quotient/remainder profile and field drop]
\label{def:quotient-remainder-profile}
Let \(\phi:D\to Q\) be a \(c\)-fold complete-fiber folding map in the
sense of \cref{def:structured-folding}.  A
\emph{complete-fiber-plus-remainder support} is
\[
                  S=\phi^{-1}(E)\sqcup R,
\]
where \(E\subseteq Q\) and
\(R\subseteq D\setminus\phi^{-1}(E)\).  The set \(R\) is the
\emph{support remainder}: it records the selected points in fibers that
are not declared complete.  A \emph{quotient/remainder profile} fixes
\(\phi\), its degree \(c\), the size \(r=\abs R\), the relevant
coefficient field, and any planted remainder data, while \(E\) and the
unfixed part of \(R\) range subject to the displayed disjointness.
A fixed-\(R\) profile is the subprofile in which the remainder support
itself is part of the label.  The adjective \emph{Euclidean} used below
refers to separating the locator factor of degree \(r<c\) from the
composed quotient locator; it does not mean that an arbitrary support
remainder has size less than \(c\).

We call \(\B_\phi\subseteq\B\) a \emph{scaled quotient coefficient
field} when, for some \(\eta\in\B^\times\),
\[
                         Q=\phi(D)\subseteq\eta\B_\phi .
\]
When this field is recorded in a profile, it is taken minimal among the
subfields with this property.
Then the \(j\)-th elementary quotient coefficient lies in
\(\eta^j\B_\phi\).  A \emph{field drop} is the use of a proper such
subfield.  It is a \emph{positive-rate field drop} along a sequence when
\[
       1-\frac{\log\abs{\B_\phi}}{\log\abs\B}
\]
stays bounded away from zero and a positive-density number on the
entropy scale of quotient-prefix slots is live.  A field drop changes the number of available coefficient values and is therefore recorded explicitly.
\end{definition}

\begin{definition}[Balanced quotient core]
\label{def:balanced-quotient-core}
A \emph{split locator} is a monic squarefree polynomial all of whose roots
lie in the evaluation domain \(D\).
After common locator factors, fixed quotient data, and planted remainder
points have been removed, let \(A\) and \(B\) be two monic split
locators of the same degree \(e\).  They form a
\emph{depth-\(w\) balanced core} if they have the same first \(w\)
nonleading coefficients, equivalently
\[
                         \deg(A-B)\le e-w-1.
\]
It is a \emph{balanced quotient core} when the remaining locator factors
are pullbacks through the same folding map \(\phi\).  The word
\emph{balanced} records equality of degree and leading normalization;
the \emph{core} consists of the roots still allowed to move after common
and planted factors are removed.  If the resulting projective locator
family is one-dimensional it is a split pencil
\(Q_0+\lambda Q_1\); in higher dimension this definition by itself
supplies no distinct-ray bound.  The exact determinant-pivot compiler
still applies, but controlling its boundary maps requires a separate
quantitative upper bound on their fiber sizes.
\end{definition}

For a monic polynomial
\[
 A(X)=X^d+a_1X^{d-1}+\cdots+a_d
\]
and an integer $0\le u\le d$, write
\(\operatorname{pref}_u(A)=(a_1,\ldots,a_u)\).  Equivalently, if
\(A^\vee(Z):=Z^dA(Z^{-1})\), then
\(\operatorname{pref}_u(A)\) is the coefficient vector of
\(A^\vee(Z)\bmod Z^{u+1}\), with its constant coefficient omitted.

\begin{theorem}[Exact quotient--remainder prefix normal form]
\label{thm:exact-quotient-remainder-normal-form}
Let $\B\subseteq\F$ be finite fields, let $D\subseteq\B$ have $n$
distinct elements, and let $\phi\in\B[X]$ be monic of degree $c\ge2$.
Assume that the restriction
\[
             \phi:D\longrightarrow Q:=\phi(D)
\]
has complete fibers of size $c$; in particular, $N:=|Q|=n/c$.
Fix integers $m,r$ with $0\le r<c$ and $cm+r\le n$.  For
\[
 E\in\binom Qm,
 \qquad
 R\in\binom{D\setminus\phi^{-1}(E)}r,
\]
put
\[
 \begin{split}
 V_E(Y)&:=\prod_{y\in E}(Y-y)
       =Y^m+\sum_{j=1}^m v_j(E)Y^{m-j},\\
 P_R(X)&:=\prod_{x\in R}(X-x)
       =X^r+\sum_{i=1}^r p_i(R)X^{r-i},\\
 L_{E,R}(X)&:=P_R(X)V_E(\phi(X)).
 \end{split}                                      \tag{QR1}\label{eq:qr-locators}
\]
Then $L_{E,R}$ is the monic locator of the support
\(\phi^{-1}(E)\sqcup R\), of degree $A=cm+r$.

Let $0\le w\le A$ and put $d=\lfloor w/c\rfloor$.

\begin{enumerate}[label=\textup{(\roman*)},leftmargin=*]
\item If $w\ge r$, every nonempty fiber of
      \((E,R)\mapsto\operatorname{pref}_w(L_{E,R})\) determines $R$
      uniquely and, after $R$ is recovered, is exactly one quotient-prefix
      fiber.  More precisely, for every prefix value $z$ there are unique
      data $R_z$ and $\eta_z\in\B^d$, whenever the fiber is nonempty, such
      that
      \[
      \begin{split}
       &\{(E,R):\operatorname{pref}_w(L_{E,R})=z\}\\
       &\quad=
       \{(E,R_z): E\in\tbinom{Q\setminus\phi(R_z)}m,
           \ \operatorname{pref}_d(V_E)=\eta_z\}.
      \end{split}                                  \tag{QR2}\label{eq:qr-fiber-exact}
      \]
      Thus no factor counting possible remainder sets occurs in a fixed
      global prefix fiber.

\item If $w<r$, then necessarily $w<c$, and the quotient set $E$ is
      invisible to the first $w$ coefficients.  There is a fixed
      unitriangular bijection $T_{\phi,m,w}:\B^w\to\B^w$ such that
      \[
          \operatorname{pref}_w(L_{E,R})
          =T_{\phi,m,w}(\operatorname{pref}_w(P_R))                 \tag{QR3}
      \]
      for every admissible $(E,R)$.  Consequently, with the convention
      \(\binom uv=0\) for $v>u$,
      \[
      \begin{split}
       &\#\{(E,R):\operatorname{pref}_w(L_{E,R})=z\}\\
       &\quad=
       \sum_{\substack{R\in\binom Dr\\
          \operatorname{pref}_w(P_R)=T_{\phi,m,w}^{-1}(z)}}
          \binom{N-|\phi(R)|}{m}.
      \end{split}                                  \tag{QR4}\label{eq:qr-shallow-remainder}
      \]
      Thus this case is a remainder-prefix problem, multiplied by the
      number of quotient sets still available after $R$ is fixed.
\end{enumerate}
\end{theorem}

\begin{proof}
For a monic degree-$e$ polynomial $A$, set
\(A^\vee(Z)=Z^eA(Z^{-1})\).  Also set
\[
       \Phi(Z):=Z^c\phi(Z^{-1})
       =1+\phi_1Z+\cdots+\phi_cZ^c.
\]
Taking reciprocal polynomials in \eqref{eq:qr-locators} gives the exact
formal-power-series identity
\[
 L_{E,R}^\vee(Z)=P_R^\vee(Z)
 \left(
   \Phi(Z)^m+
   \sum_{j=1}^m v_j(E)Z^{cj}\Phi(Z)^{m-j}
 \right).                                          \tag{QR5}\label{eq:qr-infinity}
\]
All factors on the right have constant coefficient one except for the
displayed scalar coefficients $v_j(E)$.

Suppose first that $w\ge r$.  Since $r<c$, no term containing a $v_j$
occurs below degree $c$.  For $1\le i\le r$, the coefficient of $Z^i$
on the right of \eqref{eq:qr-infinity} is
\(p_i(R)\) plus an expression depending only on
\(p_1(R),\ldots,p_{i-1}(R)\), on $m$, and on $\Phi$.  The first $r$
coefficients therefore recover $p_1(R),\ldots,p_r(R)$ successively.
They recover the monic polynomial $P_R$, hence its root set $R$, uniquely.

We may now divide \(L_{E,R}^\vee\) by the unit
\(P_R^\vee\) modulo $Z^{w+1}$.  In the remaining factor of
\eqref{eq:qr-infinity}, the coefficient of $Z^{cj}$ is
\(v_j(E)\) plus an expression depending only on
\(v_1(E),\ldots,v_{j-1}(E)\), on $m$, and on $\Phi$.  Hence the
coefficients at degrees $c,2c,\ldots,dc$ recover
\(v_1(E),\ldots,v_d(E)\) successively.  Conversely, $P_R$ and these $d$
coefficients determine the right side of \eqref{eq:qr-infinity} modulo
$Z^{w+1}$, because every term containing $v_j$ with $j>d$ starts in degree
$cj>w$.  This proves \eqref{eq:qr-fiber-exact}; admissibility of $E$ is
exactly the condition $E\cap\phi(R)=\varnothing$.

If $w<r<c$, every term involving a $v_j$ starts in degree at least
$c>w$, so modulo $Z^{w+1}$ the bracket in \eqref{eq:qr-infinity} is the
fixed unit $\Phi(Z)^m$.  Multiplication by this unit sends the first $w$
coefficients of $P_R^\vee$ to the first $w$ coefficients of
$L_{E,R}^\vee$ by a unitriangular transformation.  It is independent of
$E$ and invertible.  For a fixed $R$, the admissible $E$ are precisely the
$m$-subsets of $Q\setminus\phi(R)$, of which there are
\(\binom{N-|\phi(R)|}{m}\).  Summing over the indicated remainder-prefix
fiber proves \eqref{eq:qr-shallow-remainder}.
\end{proof}

\begin{proposition}[Complete support factorization and its limitation]
\label{prop:complete-support-factorization}
Under the complete-fiber hypotheses of
\cref{thm:exact-quotient-remainder-normal-form}, every support
\(S\subseteq D\) has the unique decomposition
\[
 E(S)=\{y\in Q:\phi^{-1}(y)\subseteq S\},\qquad
 R(S)=S\setminus\phi^{-1}(E(S)),
 \tag{QR5a}\label{eq:complete-support-factorization}
\]
and
\(Q_S(X)=P_{R(S)}(X)V_{E(S)}(\phi(X))\).  The remainder meets every
nonfull \(\phi\)-fiber in at most \(c-1\) points, but its total size may
be \(\Theta(n)\).  The reciprocal identity \eqref{eq:qr-infinity} remains
valid for this arbitrary remainder.  When \(\abs{R(S)}\ge c\), however,
the quotient coefficients and the remainder coefficients begin to
interlace at degree \(c\), so the triangular recovery theorem above does
not give a comparison without an additional partial-occupancy atlas.
\end{proposition}

\begin{proof}
A fiber is put into \(E(S)\) exactly when all of its points lie in \(S\);
all other selected points form \(R(S)\).  This proves existence,
disjointness, the occupancy bound, and uniqueness.  Multiplying the
corresponding linear factors gives the locator identity.  The reciprocal
calculation used for \eqref{eq:qr-infinity} is purely multiplicative and
does not require \(\deg P_R<c\).  If \(\deg P_R\ge c\), both
\(p_c(R)\) and \(v_1(E)\) occur in degree \(c\), which proves the stated
limitation of the triangular separation.
\end{proof}

\begin{theorem}[Exact partial-occupancy disintegration]
\label{thm:exact-partial-occupancy}
Let \(D=D_0\sqcup X\), where \(\abs X=b\), and let
\(\phi:D_0\twoheadrightarrow Q\) have exactly \(N=\abs Q\) fibers, each
of cardinality \(c\).  For \(S\subseteq D\), put
\[
 X_S=S\cap X,\qquad
 E(S)=\{y\in Q:\phi^{-1}(y)\subseteq S\},\qquad
 R(S)=(S\cap D_0)\setminus\phi^{-1}(E(S)).
\]
Write
\[
 t=\abs{X_S},\qquad m=\abs{E(S)},\qquad
 p=\abs{\phi(R(S))},\qquad r=\abs{R(S)}.
\]
Then this datum is unique, every fiber meeting \(R(S)\) contains between
one and \(c-1\) points of \(R(S)\), and \(\abs S=t+cm+r\).  If
\(\Omega_{t,m,p,r}\) denotes the family of supports with these four
parameters, then
\[
 \abs{\Omega_{t,m,p,r}}
 =\binom bt\binom Np\binom{N-p}m
   [x^r]\bigl((1+x)^c-1-x^c\bigr)^p.                 \tag{PO1}
\]
Consequently the nonempty families \(\Omega_{t,m,p,r}\) partition
\(\binom Da\), and their exact add-back is
\[
 \sum_{\substack{t,m,p,r\\t+cm+r=a}}\abs{\Omega_{t,m,p,r}}
 =\binom{cN+b}{a}.                                   \tag{PO2}
\]
If \(\phi\) is the restriction of a monic degree-\(c\) polynomial and
every fiber \(\phi^{-1}(y)\) is the complete simple root set in \(D_0\)
of \(\phi(X)-y\), then their locators satisfy
\[
             Q_S=Q_{X_S}P_{R(S)}V_{E(S)}(\phi).
\]
\end{theorem}

\begin{proof}
Full fibers, partial fibers, and empty fibers are mutually exclusive, so
the displayed decomposition is canonical.  Choose the \(t\) exceptional
points, the \(p\) partial fibers, and the \(m\) full fibers in
\(\binom bt\binom Np\binom{N-p}m\) ways.  On one partial fiber the
weight enumerator of the allowed nonempty proper subsets is
\((1+x)^c-1-x^c\); hence the coefficient in \textup{(PO1)} counts the
choices having total partial weight \(r\).  Summing \textup{(PO1)} and
extracting the coefficient of \(x^a\) gives
\[
 [x^a](1+x)^b
 \left(1+x^c+\bigl((1+x)^c-1-x^c\bigr)\right)^N
 =[x^a](1+x)^{b+cN}=\binom{b+cN}{a}.
\]
Multiplying the linear factors belonging to the three disjoint parts of
the support proves the locator identity.
\end{proof}

\begin{proposition}[Partial-occupancy Fourier factorization]
\label{prop:partial-occupancy-fourier}
In the setting of \cref{thm:exact-partial-occupancy}, let \(G\) be a
finite abelian group, let \(g:D\to G\), and put
\(\Phi(S)=\sum_{x\in S}g(x)\).  For a character
\(\chi\in\widehat G\) and a regular fiber \(F_y=\phi^{-1}(y)\), define
\[
 A_y(\chi)=\prod_{x\in F_y}\chi(g(x)),\qquad
 P_y(\chi;z)=
 \prod_{x\in F_y}\bigl(1+z\chi(g(x))\bigr)-1-z^cA_y(\chi).
\]
Then
\[
 \sum_{S\in\Omega_{t,m,p,r}}\chi(\Phi(S))
 =[u^tv^m\zeta^pz^r]
 \prod_{x\in X}\bigl(1+u\chi(g(x))\bigr)
 \prod_{y\in Q}
 \bigl(1+vA_y(\chi)+\zeta P_y(\chi;z)\bigr).          \tag{PO3}
\]
Consequently, for every \(s\in G\),
\[
 \abs{\{S\in\Omega_{t,m,p,r}:\Phi(S)=s\}}
 =\frac1{\abs G}\sum_{\chi\in\widehat G}
   \overline{\chi(s)}
   \sum_{S\in\Omega_{t,m,p,r}}\chi(\Phi(S)).         \tag{PO4}
\]
At the trivial character, \textup{(PO3)} is exactly the support count
\textup{(PO1)}.  Thus support add-back is unconditional, while a finite
flatness theorem is exactly a bound for the nontrivial-character sum in
\textup{(PO4)}.
\end{proposition}

\begin{proof}
For a fixed regular fiber the three terms in the second product encode,
respectively, an empty fiber, a full fiber, and a nonempty proper subset.
The variables \(v,\zeta,z\) record the numbers \(m,p,r\), while \(u\)
records the exceptional weight \(t\).  Multiplicativity of \(\chi\) turns
the character of the support sum into the product of its selected-point
characters, proving \textup{(PO3)} by coefficient extraction.
Formula \textup{(PO4)} is character orthogonality on \(G\).
\end{proof}

\begin{remark}[Effective normalization of a partial-occupancy slice]
Fix a nonempty slice \(\Omega_\lambda=\Omega_{t,m,p,r}\), choose
\(S_0\in\Omega_\lambda\), and let
\[
 G_\lambda=
 \left\langle\Phi(S)-\Phi(S_0):S\in\Omega_\lambda\right\rangle
 \le G.
\]
The attained boundary values lie in the affine coset
\(\Phi(S_0)+G_\lambda\), and the effective inversion is
\[
 \abs{\{S\in\Omega_\lambda:\Phi(S)=s\}}
 =\frac1{\abs{G_\lambda}}
   \sum_{\chi\in\widehat{G_\lambda}}
   \overline{\chi(s-\Phi(S_0))}
   \sum_{S\in\Omega_\lambda}
      \chi(\Phi(S)-\Phi(S_0)).                       \tag{PO5}
\]
Every character \(\chi\) of \(G_\lambda\) extends to a character
\(\widetilde\chi\) of the ambient finite abelian group, and
\[
 \sum_{S\in\Omega_\lambda}\chi(\Phi(S)-\Phi(S_0))
 =\overline{\widetilde\chi(\Phi(S_0))}
   \sum_{S\in\Omega_\lambda}\widetilde\chi(\Phi(S)).
\]
The right side may therefore be evaluated from \textup{(PO3)} and is
independent of the chosen extension, because different extensions agree on
all support differences.  Equivalently, ambient characters group by their restriction
to \(G_\lambda\), and the annihilator multiplicity cancels exactly against
the ambient denominator.  This identifies the correct Fourier denominator;
it does not assert that the realized image fills the affine group.
\end{remark}

\begin{remark}[What partial-occupancy add-back does and does not prove]
The exponentially many local proper-subset choices are already contained
inside the coarse slices \(\Omega_{t,m,p,r}\); they incur no separate
profile-count loss in \textup{(PO2)}.  The theorem does not bound the
nontrivial Fourier terms in \textup{(PO4)}, the size of a realized
boundary image, or the projection of a witness cell to distinct slopes.
Those are separate analytic and ray payments.
\end{remark}


\section{Primitive inverse from a Sidon payment at accessible moment order}
\label{sec:primitive-inverse}

We now isolate the exact part of the primitive-fiber argument supplied by
additive combinatorics.  It is intentionally stated at effective image
scale.  In this section a \emph{primitive residual leaf} means a family
remaining after specified algebraic first-match cells have been removed;
the adjective is bookkeeping and does not itself assert Fourier flatness,
image coverage, or a slope bound.

\begin{definition}[Image-normalized fiber moments]
\label{def:primitive-fiber-moments}
For every \(N\), let \(T_N\) be a set of size \(N\), let
\(\Omega_N^0\subseteq\{0,1\}^{T_N}\) be a nonempty reference profile
slice---the full family of support incidence vectors sharing the fixed
profile parameters before any first-match deletion---let
\(\Omc_N\subseteq\Omega_N^0\) be a primitive residual leaf,
and let \(\Phi_N:\Omega_N^0\to V_N\) be a map to a finite set.  Write
\[
 I_N=\Phi_N(\Omega_N^0),\qquad L_N=\abs{I_N},\qquad
 \barN_N=\frac{\abs{\Omega_N^0}}{L_N},
\]
so \(L_N\le\abs{\Omega_N^0}\) and \(\barN_N\ge1\).
For \(s\in I_N\), put
\(F_{N,s}=\Omc_N\cap\Phi_N^{-1}(s)\) and
\(f_{N,s}=\abs{F_{N,s}}\).
For every finite \(F\subseteq\Z^{T_N}\), its additive energy is
computed in the torsion-free group \(\Z^{T_N}\) by
\[
 E(F)=\#\{(x_1,x_2,x_3,x_4)\in F^4:
                     x_1+x_2=x_3+x_4\}.
\]
Put \(\Delta_{N,s}=E(F_{N,s})/f_{N,s}^3\) when \(f_{N,s}>0\), and
\(\Delta_{N,s}=0\) when \(f_{N,s}=0\).
For an integer \(q\ge2\) and \(\sigma>0\), define the ordinary and
low-energy, or \emph{Sidon-heavy}, normalized moments by
\begin{align}
 \Gord_{N,q}
   &=\frac1{L_N}\sum_{s\in I_N}
        \left(\frac{f_{N,s}}{\barN_N}\right)^q,
                                                        \tag{PI1}\label{eq:pi-ordinary}\\
 \Gsid_{N,q,\sigma}
   &=\frac1{L_N}\sum_{\substack{s\in I_N\\
                         \Delta_{N,s}<e^{-\sigma N}}}
        \left(\frac{f_{N,s}}{\barN_N}\right)^q.       \tag{PI2}\label{eq:pi-sidon}
\end{align}
The terminology records that a large fiber with exponentially small
\(\Delta_{N,s}\) can carry large ordinary moment mass while remaining too
additively sparse for a high-energy inverse theorem.
\end{definition}

\begin{lemma}[Exact ambient/image conversion]
\label{lem:pi-ambient-image}
Suppose \(I_N\subseteq W_N\), put \(A_N=\abs{W_N}\), extend
\(f_{N,s}\) by zero on \(W_N\setminus I_N\), and set
\(\barN_N^{\rm amb}=\abs{\Omega_N^0}/A_N\).  If the ambient moments are
defined by replacing \(L_N,\barN_N,I_N\) in
\eqref{eq:pi-ordinary}--\eqref{eq:pi-sidon} with
\(A_N,\barN_N^{\rm amb},W_N\), then
\[
 (\Gord_{N,q})^{\rm amb}
   =\left(\frac{A_N}{L_N}\right)^{q-1}\Gord_{N,q},
 \qquad
 (\Gsid_{N,q,\sigma})^{\rm amb}
   =\left(\frac{A_N}{L_N}\right)^{q-1}
      \Gsid_{N,q,\sigma}.                            \tag{PI3}\label{eq:pi-normalization}
\]
Thus an ambient-scale moment payment transfers to image scale, whereas
the reverse transfer requires \(A_N/L_N=e^{o(N)}\) at the relevant
orders.
\end{lemma}

\begin{proof}
For every attained value,
\[
 \frac{f_{N,s}}{\barN_N^{\rm amb}}
 =\frac{A_N}{L_N}\frac{f_{N,s}}{\barN_N}.
\]
Raise to the \(q\)-th power and combine the factor \(A_N^{-1}\) with
\(L_N^{-1}\).  The same calculation applies after restricting the sum by
the energy condition.
\end{proof}

\begin{theorem}[BSG and Boolean-cube difference growth]
\label{thm:bsg-cube-package}
There is an absolute constant \(C\ge1\) with the following property.  If
\(A\) is a nonempty finite subset of an abelian group, \(K\ge1\), and
\[
 E(A)\defeq
 \#\{(a_1,a_2,a_3,a_4)\in A^4:a_1+a_2=a_3+a_4\}
 \ge\frac{\abs A^3}{K},
\]
then some \(A'\subseteq A\) satisfies
\[
 \abs{A'}\ge K^{-C}\abs A,\qquad
 \abs{A'-A'}\le K^C\abs{A'} .                       \tag{PI4}\label{eq:pi-bsg}
\]
Moreover, every \(U\subseteq\{0,1\}^d\), viewed in \(\Z^d\), satisfies
\[
                         \abs{U-U}\ge\abs U^{3/2}.  \tag{PI5}\label{eq:pi-cube}
\]
\end{theorem}

The first assertion is the energy form of the
Balog--Szemer\'edi--Gowers theorem; polynomial dependence on \(K\) follows
from Gowers's refinement \cite{BalogSzemeredi,GowersSzemeredi}.  The second
is the Boolean-quasicube specialization of the sumset inequality of
Green--Matolcsi--Ruzsa--Shakan--Zhelezov
\cite{GreenMatolcsiRuzsaShakanZhelezov}: their quasicube inequality,
with its two auxiliary sets specialized to \(-U\) and \(\{0\}\), gives
\(\abs{U-U}\ge\abs U^{3/2}\).  These are the only external
additive-combinatorics inputs below.

\begin{proposition}[High-energy Boolean fibers are small]
\label{prop:high-energy-boolean-small}
There is an absolute constant \(C_0\) such that, for every
finite set \(T\), every nonempty \(F\subseteq\{0,1\}^T\), and every
\(K\ge1\),
\[
                 E(F)\ge\frac{\abs F^3}{K}
                 \quad\Longrightarrow\quad
                 \abs F\le K^{C_0}.                 \tag{PI6}\label{eq:pi-high-energy}
\]
\end{proposition}

\begin{proof}
Apply \cref{thm:bsg-cube-package} to obtain \(F'\subseteq F\) satisfying
\eqref{eq:pi-bsg}.  The cube bound gives
\(\abs{F'}^{3/2}\le\abs{F'-F'}\le K^C\abs{F'}\), so
\(\abs{F'}\le K^{2C}\).  Since
\(\abs F\le K^C\abs{F'}\), we get \(\abs F\le K^{3C}\); take
\(C_0=3C\).
\end{proof}

\begin{theorem}[Primitive flatness from a Sidon payment at accessible moment order]
\label{thm:primitive-q-after-sidon}
In the setting of \cref{def:primitive-fiber-moments}, suppose there are
integers \(q_N\ge2\) such that the moment order is \emph{accessible}, by
which we mean
\[
              \frac{\log L_N}{q_N}=o(N)             \tag{PI7}\label{eq:pi-access}
\]
and, for every fixed \(\sigma>0\),
\[
              \Gsid_{N,q_N,\sigma}\le e^{o(Nq_N)}.   \tag{PI8}\label{eq:pi-payment}
\]
Then the primitive fibers are flat at effective image scale:
\[
              \max_{s\in I_N}f_{N,s}
                 \le e^{o(N)}\barN_N.               \tag{PI9}\label{eq:pi-q}
\]
\end{theorem}

\begin{proof}
Suppose \eqref{eq:pi-q} fails.  Along a subsequence there are
\(\eta>0\) and \(s_N\in I_N\) such that
\(f_{N,s_N}\ge e^{\eta N}\barN_N\).  By
\eqref{eq:pi-ordinary} and \eqref{eq:pi-access},
\[
 \Gord_{N,q_N}
 \ge L_N^{-1}e^{\eta Nq_N}
 =\exp\!\bigl(\eta Nq_N-\log L_N\bigr)
 \ge e^{3\eta Nq_N/4}
\]
for all sufficiently large \(N\) in this subsequence.

Let \(C_0\) be the constant in
\cref{prop:high-energy-boolean-small} and fix
\(0<\sigma<\eta/(4C_0)\).  The payment
\eqref{eq:pi-payment} is at most \(e^{\eta Nq_N/4}\) for large \(N\).
After subtracting it from the ordinary moment, the contribution of image
values with \(\Delta_{N,s}\ge e^{-\sigma N}\) is at least
\(e^{\eta Nq_N/2}\).  Hence one such image value satisfies
\(f_{N,s}/\barN_N\ge e^{\eta N/2}\).  Since \(\barN_N\ge1\), this gives
\(f_{N,s}\ge e^{\eta N/2}\).  On the other hand,
\(\Delta_{N,s}\ge e^{-\sigma N}\) and
\cref{prop:high-energy-boolean-small}, with \(K=e^{\sigma N}\), give
\(f_{N,s}\le e^{C_0\sigma N}<e^{\eta N/4}\), a contradiction.
\end{proof}

\begin{remark}[What the Gowers/cube step proves]
\label{rem:gowers-cube-scope}
The theorem is an unconditional implication from a Sidon payment at an
accessible moment order, namely
\eqref{eq:pi-access}--\eqref{eq:pi-payment}; it is not an unconditional
verification of that payment.  In a
smooth/circle application, \eqref{eq:pi-payment} still requires an
effective-image Fourier theorem---for example aggregate minor- and
major-character bounds---or a direct low-energy moment estimate on every
primitive residual leaf.  Even after \eqref{eq:pi-q}, a determinant or ray
compiler is needed to turn support fibers into distinct MCA slopes, and a
profile-envelope comparison is needed to reach a threshold.

The phrase ``Gowers/cube step'' refers here to the polynomial
Balog--Szemer\'edi--Gowers theorem followed by Boolean-quasicube difference
growth.  No Gowers-uniformity norm or Gowers-cube counting theorem is used.
\end{remark}


\section{Asymptotic and finite regimes}\label{sec:complementary}


\begin{theorem}[Capacity-edge shallow-prefix exponent]
\label{thm:unconditional-asymptotic-rs-mca}
Let \(C_n=\RS_{\F_n}(D_n,k_n)\), where
\(D_n\subseteq\B_n\subseteq\F_n\), \(\abs{D_n}=n\), and
\(k_n/n\to\rho\in(0,1)\).  Let
\(a_n\in\{k_n+1,\ldots,n\}\), and put
\(w_n=a_n-k_n-1\ge0\).  Assume
\[
 w_n\log_2\abs{\B_n}=o(n),\qquad \abs{\F_n}>n,     \tag{AS1}
\]
and, for some \(\phi\ge0\),
\[
 \frac1n\log_2\abs{\F_n}\to\phi,\qquad
 \frac1n\log_2(\abs{\F_n}-n)\to\phi.              \tag{AS2}
\]
Then \(a_n/n\to\rho\), and for the full-field challenge
\[
 \begin{aligned}
 \frac1n\log_2 B_{C_n}^{\rm MCA}(a_n)
   &=\min\{H_2(\rho),\phi\}+o(1),\\
 \frac1n\log_2 e_{\rm MCA}
 \left(C_n,1-\frac{a_n}{n}\right)
   &=-\bigl(\phi-H_2(\rho)\bigr)_++o(1).
 \end{aligned}                                    \tag{AS3}
\]
No smoothness, Fourier estimate, or auxiliary ray hypothesis is assumed.
\end{theorem}

\begin{proof}
Since \(\abs{\B_n}\ge n\), \textup{(AS1)} gives \(w_n=o(n)\), hence
\(a_n/n\to\rho\).  Put
\(M_n=\binom n{a_n}\) and
\(L_n=\ceil{M_n\abs{\B_n}^{-w_n}}\).
Exact-support uniqueness gives
\(B_{C_n}^{\rm MCA}(a_n)\le\min\{\abs{\F_n},M_n\}\), while the exact
locator-prefix list and collision-aware pole compiler give
\[
 B_{C_n}^{\rm MCA}(a_n)\ge
 \left\lceil\frac{L_n(\abs{\F_n}-n)}
 {\abs{\F_n}-n+k_n(L_n-1)}\right\rceil.           \tag{AS4}
\]
Stirling and \textup{(AS1)} give
\(\log_2L_n=H_2(\rho)n+o(n)\).  For \(x,y>0\),
\(xy/(x+k_ny)\ge\frac12\min\{y,x/k_n\}\), so
\textup{(AS2)} shows that the lower bound in \textup{(AS4)} has exponent
\(\min\{H_2(\rho),\phi\}\), matching the upper bound.  Subtract
\(\log_2\abs{\F_n}\) to obtain the error exponent.
\end{proof}

\begin{corollary}[Sharp capacity-edge criterion]
\label{cor:unconditional-rs-mca-capacity}
Under the hypotheses of
\cref{thm:unconditional-asymptotic-rs-mca}, let
\(\eps_n=2^{-\lambda n+o(n)}\), where \(\lambda\ge0\).
If
\[
                         \lambda<\phi-H_2(\rho),   \tag{AS5}
\]
then eventually
\[
 \begin{aligned}
 a_n^*&=k_n+1,&
 \delta_{n,\mathrm{sub,grid}}^*&=1-\frac{k_n+1}{n}
       =1-\rho+o(1),&
 \delta_{n,\mathrm{sub,sup}}^*&=1-\frac{k_n}{n},\\
 \delta_{C_n,\mathrm{off,grid}}^*
 &=\delta_{C_n,\mathrm{off,sup}}^*=1.&&&
 \end{aligned}                                    \tag{AS6}
\]
If \(\lambda>(\phi-H_2(\rho))_+\), every agreement sequence satisfying
\((a_n-k_n-1)\log\abs{\B_n}=o(n)\) is unsafe.  The equality case is
constant-sensitive.  Thus the strict capacity branch of
\cref{thm:intro-exponential-frontier} is exponent-sharp at the first
agreement, but it does not supply a matching attack at a fixed radius
away from capacity.
\end{corollary}

\begin{proof}
At \(a=k_n+1\), the support atlas gives
\(B_{C_n}^{\rm MCA}(a)\le\binom n{k_n+1}
=2^{H_2(\rho)n+o(n)}\), whereas the target budget has exponent
\(\phi-\lambda\); \textup{(AS5)} makes the first agreement safe.
The official conclusion follows from \cref{thm:official-saturation}.
The converse is the strict reverse comparison in \textup{(AS3)}; when it
applies at \(a_n=k_n+1\), the MDS plateau makes the entire official tail
with \(a(\delta)\le k_n+1\) unsafe.
\end{proof}

\begin{corollary}[Exact binary additive-domain examples]
\label{cor:binary-additive-examples}
Let \(k=2^{40}\), let \(D\) be any \(\F_2\)-linear subspace of the
indicated field having the indicated size, and take the target
\(2^{-128}\).  In each row put \(C=\RS_{\F}(D,k)\).  The four rows
\[
\begin{array}{c|c|c|c}
 k/n&\F&n&\delta_{C,\mathrm{sub,sup}}^*(2^{-128})\\ \hline
 1/2&\F_{2^{166}}&2^{41}&1/8\\
 1/4&\F_{2^{168}}&2^{42}&1/4\\
 1/8&\F_{2^{169}}&2^{43}&1/4\\
 1/16&\F_{2^{170}}&2^{44}&1/4
\end{array}
\]
have exactly the displayed supremal MCA transitions, and every endpoint
is unsafe.  Since every transition occurs before capacity, the official
supremum has the same value.
\end{corollary}

\begin{proof}
For \(q=2^m\), the slope budget is
\(B=\floor{q2^{-128}}=2^{m-128}\).  In the four rows this is respectively
\(2^{38},2^{40},2^{41},2^{42}\), and
\(3(B-1)\le n-k\).  Thus
\cref{cor:exact-deep-numerator} gives numerator \(B\) at radius
\(B-1\), while \cref{prop:universal-tangent-floor} gives at least
\(B+1\) slopes at radius \(B\).  Translate the largest safe grid radius
to the real closed-ball supremum \(B/n\).
\end{proof}

\begin{proposition}[Exact unsafe test]\label{prop:simple-pole-lower}
Let \(C_n=\RS_{\F_n}(D_n,k_n)\), where
\(D_n\subseteq\B_n\subseteq\F_n\), \(\abs{D_n}=n\),
\(q_n=\abs{\F_n}>n\), and \(1\le k_n<n\).  Fix a nonempty challenge
set \(\Gamma_n\subseteq\F_n\), a target \(\eps_n\in[0,1]\), and an
agreement \(k_n+1\le a_n\le n\), and put
\[
 B_n^*=\floor{\eps_n\abs{\Gamma_n}},\qquad
 L_n=\left\lceil
 \binom n{a_n}\abs{\B_n}^{-(a_n-k_n-1)}
 \right\rceil.
\]
If
\[
 \left\lceil \frac{\abs{\Gamma_n}}{q_n}
 \left\lceil
 \frac{L_n(q_n-n)}{q_n-n+k_n(L_n-1)}
 \right\rceil\right\rceil>B_n^*,
 \tag{UT1}\label{eq:exact-unsafe-budget}
\]
then agreement \(a_n\) is unsafe at the target.  The same statement holds
with \(L_n\) replaced by any larger identity, quotient, Chebyshev, or
remainder-profile list proved for the dimension-\(k_n+1\) code.
\end{proposition}

\begin{proof}
The identity list is \cref{prop:exact-prefix-list}, and
\cref{thm:collision-aware-pole} produces a bad-slope set \(Z\subseteq\F_n\)
of the inner-ceiling size.  Replacing a line \((r_0,r_1)\) by
\((r_0+\delta r_1,r_1)\) translates \(Z\) by \(-\delta\).  Averaging over
\(\delta\in\F_n\) gives \(\abs Z\abs{\Gamma_n}/q_n\) challenge slopes, so
some translate has at least the outer-ceiling number.  The profile variants
use the same conversion after their list sizes have been established.
\end{proof}


\begin{remark}[The target-normalized identity crossing]
\label{rem:target-normalized-identity-crossing}
The source repository's entropy equation \cite{RSMCARepo} has a rigorous
role on the attack side.
Put
\[
 \rho_n=\frac{k_n+1}{n},\qquad
 \beta_n=\log_2\abs{\B_n},\qquad
 b_n=\frac1n\log_2\max\{1,B_n^*\}.
\]
For an admissible integer agreement \(a_n\), set
\(g_n=(a_n-k_n-1)/n\) exactly and distinguish the raw locator mean
\(\widetilde L_n=\binom n{a_n}\abs{\B_n}^{-(a_n-k_n-1)}\) from
\(L_n=\ceil{\widetilde L_n}\).  Stirling gives
\[
 \frac1n\log_2\widetilde L_n
 =H_2(\rho_n+g_n)-\beta_ng_n+o(1),\qquad
 \frac1n\log_2L_n
 =\bigl(H_2(\rho_n+g_n)-\beta_ng_n\bigr)_++o(1).  \tag{IC1}
 \label{eq:identity-attack-exponent}
\]
Whenever the pole conversion in \cref{prop:simple-pole-lower} has only
subexponential loss, the target-normalized raw attack crossing is the
right endpoint
\[
 g_n^\dagger=\sup\left(\{g\in[0,1-\rho_n]:
 H_2(\rho_n+g)-\beta_ng\ge b_n\}\cup\{0\}\right). \tag{IC2}
 \label{eq:identity-attack-crossing}
\]
For \(b_n>0\) this is the ordinary exponent comparison with the ceiled
list.  At \(b_n=0\) it deliberately uses the raw nonnegative root rather
than the positive-part exponent, and constants in the finite test decide
the boundary.
The finite unsafe test remains the literal inequality
\eqref{eq:exact-unsafe-budget}; \eqref{eq:identity-attack-crossing}
locates it only up to asymptotic losses.  If \(b_n=o(1)\) and
\(\beta_n\to\infty\), the isolated root has the precise expansion
\[
 g_n^\dagger=
 \frac{H_2(\rho_n)-b_n}{\beta_n-H_2'(\rho_n)}
 +O_\rho(\beta_n^{-3}).                            \tag{IC3}
 \label{eq:identity-attack-expansion}
\]
Using the zero-level equation is justified only when \(b_n=o(1)\).
None of these formulas identifies the true MCA threshold without a
matching safe upper envelope; the smooth and circle quotient cells in the
repository show why that additional comparison cannot be assumed.
\end{remark}

\begin{theorem}[Integrated finite frontier bracket]
\label{thm:unconditional-support-envelope-bracket}
Let \(C=\RS_\F(D,k)\), let \(D\subseteq\B\subseteq\F\), put
\(q=\abs\F>n\), fix \(\varnothing\ne\Gamma\subseteq\F\), let
\(0\le\eps\le1\), and let
\(B^*=\floor{\eps\abs\Gamma}\).  For \(k+1\le a\le n\), define
\[
 \begin{split}
 L(a)&=\ceil{\binom na\abs\B^{-(a-k-1)}},\\
 P(a)&=\left\lceil\frac{\abs\Gamma}{q}
       \left\lceil\frac{L(a)(q-n)}
       {q-n+k(L(a)-1)}\right\rceil\right\rceil,\\
 U(a)&=\min\left\{\abs\Gamma,
       \min_{k+1\le b\le a}
       \floor{\frac{\binom nb}{\binom{a-1}{b-1}}}\right\}.
 \end{split}                                       \tag{SB1}
\]
Whenever the defining set is nonempty, write
\[
 a^*=\min\{a\in\{k+1,\ldots,n\}:
               B_{C,\Gamma}^{\rm MCA}(a)\le B^*\}.
\]
If \(k+1\le a_-<a_+\le n\) satisfy
\[
                         P(a_-)>B^*,\qquad U(a_+)\le B^*, \tag{SB2}
\]
then the first safe nontrivial agreement exists and
\[
                         a_-<a^*\le a_+.            \tag{SB3}
\]
If \(a_+=a_-+1\), then \(a^*=a_+\) exactly.  In particular,
\[
 \min\left\{\abs\Gamma,\binom n{k+1}\right\}\le B^*
 \quad\Longrightarrow\quad
 a^*=k+1,\qquad
 \delta_{C,\Gamma,\mathrm{off,grid}}^*
 =\delta_{C,\Gamma,\mathrm{off,sup}}^*=1.         \tag{SB4}
\]
These are unconditional finite comparisons; \textup{(SB2)} is the
literal target reserve and contains no asymptotic placeholder.  Any
stronger upper bound whose hypotheses have been checked may be intersected
with \(U(a)\); \cref{thm:retained-degree-mca} qualifies only after
\cref{ass:retained-factor-lift} is discharged.  Any
stronger realized-list construction may replace \(L(a)\) in \(P(a)\).
\end{theorem}

\begin{proof}
The exact prefix list and pole-line construction give
\(B_{C,\Gamma}^{\rm MCA}(a)\ge P(a)\) by
\cref{prop:simple-pole-lower}.  The circuit-incidence MDS envelope gives
\(B_{C,\Gamma}^{\rm MCA}(a)\le U(a)\) by
\cref{thm:intro-paving}.  Thus \(a_-\) is unsafe and \(a_+\) is
safe; monotonicity proves \textup{(SB3)} and the adjacent conclusion.
For \textup{(SB4)}, the first permitted agreement is already safe, and
\cref{thm:official-saturation} gives the complete official safe set
\([0,1]\).
\end{proof}



\section{Conclusion}\label{sec:conclusion}

The exact threshold depends first on the integer slope budget, not on a
radius slogan.  Uniformly below half the relative minimum distance,
\(B^*=\floor{\eps\abs\Gamma}\) gives first safe agreement
\(n-B^*+1\), largest safe grid radius \((B^*-1)/n\), and unsafe
continuous endpoint \(B^*/n\).  This packages the proved BCHKS
two-radius input, the exact sparse MCA layer, and the tangent construction
with all floors exposed.  If the first safe nontrivial agreement is instead
\(k+1\), the MDS plateau changes the official answer: every radius in
\([0,1]\) is safe and the official endpoint is one.  In the distinct
exponential-budget regime, Jo's shortening transfer and the BCHKS Johnson
theorem give the explicit subcapacity certificate
\(D_{\rho,\mathrm{short}}^{\rm sub}(\theta)\), strictly improving the
smallest-test MDS exponent between Johnson and capacity.  For
\(\theta>H_2(\rho)\) the subcapacity grid tends to \(1-\rho\) while the
official endpoint equals one; at equality the official outcome remains
constant-sensitive.  Below that branch the result is one-sided without a
matching constant-radius attack.

Jo's all-test-size circuit bound subsumes the support and paving envelope.
The augmented-word interpretation recorded here describes its smallest-test
specialization in matroid language and helps expose a
finite security interval rather than a one-sided score.  Together with
shortening it gives multiplicative and circle certificates beyond Johnson
at length \(512\).
The deployed prefix calculations move the unsafe side at KoalaBear and Mersenne--31,
but their adjacent failure does not certify safety.

The primitive inverse theorem discharges the high-energy branch once an
image-normalized Sidon payment is available at an accessible moment order.
At energy threshold \(e^{-\sigma N}\), its cube estimate gives
\(\abs F\le e^{C_0\sigma N}\); because the payment is assumed for every
fixed \(\sigma>0\), an alleged excess rate \(\eta>0\) is contradicted by
choosing \(\sigma<\eta/(4C_0)\).  The payment and residual ray compiler
remain substantive row inputs.  Staged shortening telescopes to the same
binomial cost as one-stage shortening, so it does not remove those
interfaces in the subexponential-budget, positive-relative-radius regime.

All smooth and circle transfers are made through complete-fiber maps or
an explicit monomial equivalence.  Fixed-length residue classes have
positive prime density, while no growing positive-density dyadic
prime-field family is claimed.  The identity entropy crossing remains a
precise attack locator after target normalization; equality with the true
near-capacity frontier still requires the missing first-match profile
exhaustion, image-normalized Sidon/Fourier payment, residual ray compiler,
and profile add-back with a safe target reserve.  These open
requirements are why the remaining difference is reported as a
soundness gap, with proved-bit gains computed from
\eqref{eq:intro-security-interval}.

\appendix
\section{Conditional parameter-retained interpolation audit}
\label{app:retained-interpolation}

The circuit-incidence bound is strongest in short, low-rate rows over very large
fields.  Near Johnson at deployed length, a parameter-retained
interpolation audit is potentially much stronger, but its exact lift is
not a source-verbatim theorem.  This appendix therefore isolates the
single required assumption before retaining the three degrees in the
BCHKS--BCIKS argument.  Three finite details
matter: ceilings in the monomial count, inseparability in small
characteristic, and the top-degree inequality used by the Hensel lift.
We state all three explicitly.

Write \(q=\abs\F\) and \(C=\RS_\F(D,K)\), so \(K\) is the dimension and code
polynomials have degree less than \(K\).  To avoid any degree--dimension
ambiguity, the interpolation below uses the conservative weight
\(i+Kj\), rather than the slightly sharper \(i+(K-1)j\).  Let
\(A=n-r\ge K+2\) be the agreement.  For a positive integer \(m\) and
positive real parameters \(D_X,D_Y,D_Z\), put
\(U=\ceil{D_X}\), \(V=\ceil{D_Y}\), and \(W=\ceil{D_Z}\).

\begin{assumption}[Parameter-retained factor lift]
\label{ass:retained-factor-lift}
Assume
\begin{align}
 V&\ge m,\quad W\ge V,\quad U>K(V-1),\notag\\
 D_X&<mA,\qquad \operatorname{char}(\F)>V-1,          \tag{RF1}\\
 (A-K-1)(2U-1)&>(n-K-1)(2K+1),\notag\\
 q&>2UD_Y.                                           \tag{RF2}
\end{align}
Let \(u_0,u_1\in\F^D\), and let \(0\ne Q\in\F[X,Y,Z]\) have
\((1,K,0)\)-weighted degree less than \(D_X\), \(Y\)-degree less
than \(D_Y\), and \((0,1,1)\)-weighted degree less than \(D_Z\).
For distinct \(\gamma\) in a set \(S\), suppose there are polynomials
\(P_\gamma\), \(\deg P_\gamma<K\), satisfying
\(Q(X,P_\gamma(X),\gamma)=0\), and chosen \(A\)-sets
\(A_\gamma\subseteq D\) on which
\(P_\gamma=u_0+\gamma u_1\).  If
\[
 \abs S>2U D_Y^2D_Z+(r+1)D_Y,                     \tag{RF3}
 \label{eq:retained-factor-threshold}
\]
then for some \(\gamma_0\in S\) there are
\(v_0,v_1\in\F[X]\), \(\deg v_i<K\), such that
\(u_i=v_i\) on the chosen support \(A_{\gamma_0}\), for
\(i=0,1\).
\end{assumption}

\begin{remark}[Source and reduction of the retained-lift hypothesis]
The intended source is the constrained-support form of the function-field Hensel
factor-lifting argument in \cite[Secs.~5.2.5--5.2.7 and App.~A]{BCIKS20},
with the degree summation of \cite[Sec.~3.2]{BCHKS25}.  We spell out the
parameter map because it is where finite-rounding errors arise.
Factor the content-free part of \(Q\) into irreducibles separable in
\(Y\), specialize at a common regular point \(x_0\), and factor each
specialization into \(Y\)-separable factors.  The characteristic
condition in \((\mathrm{RF1})\) removes the purely inseparable branch.
The excluded starting points are the zeros of a nonzero discriminant-
content polynomial of degree less than \(2UD_Y\), so the last condition
in \((\mathrm{RF2})\) supplies \(x_0\).
For a factor pair with degrees \(a_i,b_{ij},c_i\), the cited lifting
argument either charges at most \(2U a_i b_{ij}c_i\) slopes or produces
a factor \(Y-v_0(X)-Zv_1(X)\); after the lift, at most
\(a_i b_{ij}c_i\) slopes are deleted.  The integer truncation length is
controlled by \(2U-1\), not by the real number \(2D_X\); condition
\((\mathrm{RF2})\) is precisely the top-degree comparison that forces
\(\deg v_i<K\).  We have made this comparison one unit more
conservative than the source degree convention.

Summing over the factor pairs uses
\(\sum_i a_i<D_Y\), \(\sum_j b_{ij}=a_i\), and the content-free
weighted-degree budget \(\sum_i a_ic_i<D_YD_Z\).  If no useful factor
exists, these charges together with the content roots and the
\(r+1\) constrained-collinearity allowance total at most the right side
of \eqref{eq:retained-factor-threshold}.  Otherwise the lifted linear
factor gives \(P_\gamma=v_0+\gamma v_1\) on a set
\(T\) of more than \(r+1\) slopes.  Put
\(B=\{x:(u_0(x),u_1(x))\ne(v_0(x),v_1(x))\}\).
For each \(x\in B\), the two affine combinations agree at at most one
slope, while for each \(\gamma\in T\) they may disagree only outside
\(A_\gamma\), at no more than \(r\) coordinates.  Hence
\(\abs B(\abs T-1)\le r\abs T\), and the integrality together with
\(\abs T>r+1\) gives \(\abs B\le r\).  If none of the chosen supports
simultaneously explained the pair, every \(A_\gamma\) would meet
\(B\); at such a point the two affine
combinations agree, so one point of \(B\) can serve at most one slope.
Thus \(\abs T\le\abs B\le r\), a contradiction.  This is the
constrained-agreement step behind \cite[Thm.~4.3]{BCHKS25}.
What remains for unconditional use is a line-by-line verification that the
cited Hensel claims permit this simultaneous arbitrary-parameter map,
including the displayed truncation and degree ledger.  Because that combined
statement is not printed verbatim in either source, we isolate it as an
assumption instead of silently promoting the synthesis to a theorem.
\end{remark}

\begin{theorem}[Conditional retained-degree MCA bound]
\label{thm:retained-degree-mca}
Assume \cref{ass:retained-factor-lift}.  In addition to
\((\mathrm{RF1})\)--\((\mathrm{RF2})\), suppose
\[
 \sum_{j=0}^{V-1}(U-Kj)(W-j)
 >n\sum_{s=0}^{m-1}(W-s)(m-s).                   \tag{RF4}
 \label{eq:retained-interpolation-count}
\]
Then, for every challenge set \(\Gamma\subseteq\F\),
\[
 B_{C,\Gamma}^{\rm MCA}(A)
 \le\min\left\{\abs\Gamma,
 \ceil{2UD_Y^2D_Z+(r+1)D_Y}\right\}.             \tag{RF5}
 \label{eq:retained-mca-bound}
\]
\end{theorem}

\begin{proof}
For a received pair \((u_0,u_1)\), seek
\(Q=\sum q_{ijh}X^iY^jZ^h\) with
\(i+Kj<D_X\), \(j<D_Y\), and \(j+h<D_Z\), vanishing to
multiplicity \(m\) at every \((x,u_0(x)+Zu_1(x))\), \(x\in D\).
The number of coefficients is exactly the left side of
\eqref{eq:retained-interpolation-count}.  The Hasse derivative of
\(Y\)-order \(s\) contributes \(m-s\) derivative conditions, each a
polynomial identity in \(Z\) with \(W-s\) coefficients; hence the
right side is an upper bound on the number of homogeneous linear
constraints.  The strict inequality produces a nonzero interpolant.

For each MCA-bad slope choose an exact \(A\)-support and an explaining
polynomial \(P_\gamma\).  Multiplicity gives at least \(mA\) zeros of
\(Q(X,P_\gamma(X),\gamma)\), while its degree is less than
\(D_X<mA\); it is therefore zero.  If the number of bad slopes exceeded
\eqref{eq:retained-mca-bound},
\cref{ass:retained-factor-lift} would give a chosen support on which
both received coordinates are explained, contradicting the defining
MCA nontriviality.  Intersecting with \(\Gamma\) proves the result.
\end{proof}

\begin{remark}[Proof tier]
The interpolation count and the MCA contradiction above are proved in
full conditional on \cref{ass:retained-factor-lift}.  The unresolved
import is precisely the parameter-retained combination of
\cite[Claims~5.7, 5.9--5.10, Lem.~A.1 and Claim~A.2]{BCIKS20}.
BCHKS records the same useful-factor step and degree summation in
\cite[Sec.~3.2]{BCHKS25}.  The arbitrary-parameter statement above is
not printed verbatim as one numbered theorem in either source.  The BCIKS
claim and appendix numbering here refers to the extended ePrint version.
The numerical rows below use the conservative dimension weight and are
conditional on this single exposed proof obligation, not on the
repository's experimental profile ledgers.
\end{remark}

\begin{corollary}[Four conditional near-Johnson KoalaBear rows]
\label{cor:retained-koalabear}
Let
\(p_{\rm KB}=2^{31}-2^{24}+1\),
\(\F=\F_{p_{\rm KB}^6}\), \(q=\abs{\F}=p_{\rm KB}^6\),
\(n=2^{21}\), and let \(D\le\F_{p_{\rm KB}}^*\) have order \(n\).
Take the full-field challenge \(\Gamma=\F\), and for each of the four
dimensions \(K=\rho n\) put \(C=\RS_{\F}(D,K)\).  Use
\(D_X=U-1+2^{-64}\), \(D_Y=V-1+2^{-64}\), and
\(D_Z=W-1+2^{-64}\), with the following convention-safe parameters.
\[
\begin{array}{c|r|r|r|r|r}
\rho&r&m&U&V&W\\ \hline
1/2&611982&119&176735230&169&27525\\
1/4&1045433&104&109378776&209&29028\\
1/8&1352390&90&67028580&256&31500\\
1/16&1569744&78&41137824&314&34101
\end{array}                                        \tag{RF6}
\label{eq:retained-parameters}
\]
Conditional on \cref{ass:retained-factor-lift}, all hypotheses of
\cref{thm:retained-degree-mca} hold.  If \(R_{\rm ret}\) denotes the
resulting integer upper numerator, the tangent lower numerator \(r+1\)
leaves a security interval of width at most
\(g_{\rm tan}=\log_2(R_{\rm ret}/(r+1))\).  The four conditional upper
numerators, security bits, and residual gaps are
\[
\begin{array}{c|r|c|c}
\rho&R_{\rm ret}&\log_2(q/R_{\rm ret})&g_{\rm tan}\\ \hline
1/2&274589064742726105&128.002056&38.706920\\
1/4&274721012201264929&128.001363&37.935074\\
1/8&274578888391530706&128.002110&37.562917\\
1/16&274861787390229386&128.000624&37.349385
\end{array}                                        \tag{RF7}
\label{eq:retained-results}
\]
In particular every numerator is at most
\(\floor{q2^{-128}}=274980728111395087\), subject to the retained-lift
assumption.
\end{corollary}

\begin{proof}
Here \(p_{\rm KB}=127\cdot2^{24}+1\) and
\(3^{(p_{\rm KB}-1)/2}\equiv-1\pmod {p_{\rm KB}}\), so Proth's
theorem proves primality.  Also \(p_{\rm KB}>V-1\) and
\(2^{21}\mid p_{\rm KB}-1\).
For the rows in table order, the strict interpolation margins in
\eqref{eq:retained-interpolation-count} are
\(4889934,13182624,11133440,4204064\), and the left-minus-right margins in
\((\mathrm{RF2})\) are
\[
152123705899212,\quad113730027157979,\quad
63736189920080,\quad32093320774290.
\]
The other conditions follow from \(U=mA\),
\[
 mA-D_X=1-2^{-64}>0,\qquad
 \begin{array}{c|rrrr}
 \rho&1/2&1/4&1/8&1/16\\ \hline
 U-K(V-1)&574462&326872&181860&112288.
 \end{array}
\]
Substitution in \eqref{eq:retained-mca-bound}, followed by exact ceiling,
gives the four integers in \eqref{eq:retained-results}.  The bit columns
are base-two logarithms, and the last column is
\(\log_2(R_{\rm ret}/(r+1))\).
\end{proof}

\begin{remark}[What the conditional bits do and do not close]
Conditional on the retained lift, the computation crosses 128 bits at
four near-Johnson radii.  The tangent lower attacks still leave intervals of
roughly \(37.3\)--\(38.7\) bits there.  These are therefore major
safe-side reductions of the soundness gap, not exact threshold
determinations and not beyond-Johnson certificates.  The circuit-incidence theorem
and the conditional retained bound occupy complementary regimes.
\end{remark}


\section{Verification details}\label{app:verification}

This appendix records elementary coding-theoretic and integer calculations
that can be checked independently of the geometric proofs.  The short
derivations fix the normalization conventions.  Exact integer checks
accompany this source release in two scripts:
\path{verify_paving_mca_v9_2.py} checks the
unconditional paving, all-radius saturation, circle, primality,
field-budget, and Johnson
comparisons, while \path{verify_retained_bchks_v9_2.py} checks only the
arithmetic consequences of the retained-lift assumption in
\cref{app:retained-interpolation}.  It does not discharge that assumption.
Commands and SHA-256 digests are recorded in
\path{REPRODUCIBILITY_v9.2.md}.  General proof-development materials are
pinned at
\href{https://github.com/przchojecki/rs-mca/tree/02728b208ea785d02115ea967236aebf653b31ec}
{repository commit \texttt{02728b2}}; the two scripts are ancillary files
of this source release and are not claimed to occur in that earlier snapshot.
The deployed prefix values are direct evaluations of
\eqref{eq:collision-aware-pole}.

\subsection{Reed--Solomon and MDS conventions}

Let \(D=\{x_1,\ldots,x_n\}\subseteq\F\) consist of distinct points and
let \(1\le k<n\).  Evaluation on \(D\) is injective on
\(\F[X]_{<k}\): a nonzero polynomial of degree below \(k\) cannot vanish
at the \(n\ge k\) distinct points of \(D\).  It follows that
\(\RS_\F(D,k)\) has dimension \(k\).  If two codewords are distinct,
their difference is the evaluation of a nonzero polynomial of degree at
most \(k-1\), hence has at most \(k-1\) zero coordinates and at least
\(n-k+1\) nonzero coordinates.  Conversely, the evaluation of
\(\prod_{i=1}^{k-1}(X-x_i)\) has exactly \(k-1\) zero coordinates.
Thus the minimum distance is
\[
                         d_{\min}=n-k+1.
\]
This proves directly that the code is MDS.

For \(x\in D\), put
\[
 \lambda_x=\left(\prod_{y\in D\setminus\{x\}}(x-y)\right)^{-1}.
\]
The Lagrange interpolation formula for a polynomial \(f\) of degree at
most \(n-1\) shows that its leading coefficient is
\(\sum_{x\in D}\lambda_x f(x)\).  Applying this to
\(f(X)=X^j\), \(0\le j\le n-2\), gives
\[
                         \sum_{x\in D}\lambda_xx^j=0.
\]
Let \(R=n-k\) and let \(H\) be the \(R\)-by-\(n\) matrix whose column at
\(x\) is
\[
                    h_x=\lambda_x(1,x,\ldots,x^{R-1})^{\mathsf T}.
\]
If \(P\in\F[X]_{<k}\), then the \(j\)-th entry of
\(H(P(x))_{x\in D}^{\mathsf T}\) is
\(\sum_x\lambda_xx^jP(x)=0\), because \(j+\deg P\le n-2\).
Thus \(H\) annihilates the code.  Every \(R\)-by-\(R\) submatrix of
\(H\) is a nonzero diagonal scaling of a Vandermonde matrix, so it is
invertible.  Hence \(H\) has rank \(R\), its kernel has dimension \(k\),
and its kernel is exactly \(\RS_\F(D,k)\).

The same Vandermonde calculation proves that every set of at most \(R\)
columns is linearly independent.  In particular, for
\(E\subseteq D\), \(|E|\le R\), the space
\(V_E=\Span\{h_x:x\in E\}\) has dimension \(|E|\).  A word \(u\) agrees
with a codeword outside \(E\) if and only if its syndrome lies in \(V_E\).
Indeed, writing \(u=c+e\) with \(e\) supported on \(E\) gives one
direction.  Conversely, if
\(Hu^{\mathsf T}=\sum_{x\in E}e_xh_x\), subtracting the word whose
coordinates on \(E\) are \(e_x\) leaves zero syndrome and therefore a
codeword.  This is the precise error-support form of the syndrome
normalization used in \cref{prop:syndrome-line-normal-form}.

For completeness, the fixed-support slope uniqueness used repeatedly in
the paper is also an MDS-free linear observation.  Suppose a support
\(S\) explains \(u_0+\gamma u_1\) and \(u_0+\gamma' u_1\) for distinct
\(\gamma,\gamma'\).  Subtracting the two explaining codewords and dividing
by \(\gamma-\gamma'\) explains \(u_1\) on \(S\); subtracting
\(\gamma\) times this explanation from the first explains \(u_0\) there.
Thus \(S\) is common.  Consequently a noncommon support contributes at
most one finite MCA-bad slope.

\begin{lemma}[Monotonicity and exact endpoint conversion]
\label{lem:appendix-endpoint}
For fixed \(C\) and nonempty \(\Gamma\), the integer numerator
\(B_{C,\Gamma}^{\rm MCA}(a)\) is nonincreasing in \(a\).  Let
\(0\le\eps\le1\), let
\(B^*=\floor{\eps|\Gamma|}\), and suppose \(a_0>k+1\) is the first safe
nontrivial agreement, while \(a_0-1\) is unsafe.  Then the official and
subcapacity transitions agree: the largest safe grid radius is
\(1-a_0/n\), whereas the safe real radii form
\[
                 \left[0,\frac{n-a_0+1}{n}\right),
\]
so their supremum is one grid step larger and is not attained.
If \(C\) is MDS and \(a_0=k+1\), the subcapacity safe set is
\([0,1-k/n)\), its largest grid point is \(1-(k+1)/n\), and the official
safe set is \([0,1]\).
\end{lemma}

\begin{proof}
Every witness at agreement at least \(a+1\) is also a witness at
agreement at least \(a\), proving monotonicity.  For a real radius
\(\delta\), the closed Hamming ball permits precisely
\(r=\floor{\delta n}\) errors and therefore requires agreement
\(n-r\).  The condition \(n-\floor{\delta n}\ge a_0\) is equivalent to
\(\floor{\delta n}\le n-a_0\), or
\(\delta<(n-a_0+1)/n\).  The largest grid point in that interval is
\((n-a_0)/n=1-a_0/n\).
In the final case, \cref{eq:subcapacity-constant} extends safety from
agreement \(k+1\) to every smaller threshold, which is exactly the official
tail from capacity through radius one.
\end{proof}

This lemma explains the half-open intervals in the main theorem.  In
particular, at target \(2^{-128}\), writing
\(B=\floor{q/2^{128}}\), exact numerator \(B\) at error radius \(B-1\)
and numerator at least \(B+1\) at radius \(B\) give the safe set
\([0,B/n)\), rather than a closed interval.

\subsection{The quadratic boundary}

Put \(R=n-k\) and
\[
                  F_{n,k}(r)=r^2-n(3r-R).
\]
The mean-overlap hypothesis is \(F_{n,k}(r)\ge0\).  Its two roots are
\[
       \frac{3n\pm\sqrt{9n^2-4nR}}2
       =\frac{3n\pm\sqrt{n(5n+4k)}}2.
\]
Since \(F_{n,k}(0)=nR>0\) and
\(F'_{n,k}(r)=2r-3n<0\) for \(0\le r\le n\), the admissible initial
interval ends at
\[
 r_{\rm quad}(n,k)=
 \floor{\frac{3n-\sqrt{n(5n+4k)}}2}.
\]
If \(k/n\to\rho\), division by \(n\) gives
\[
 \frac{r_{\rm quad}(n,k)}n
 \longrightarrow
 \alpha_\rho=\frac{3-\sqrt{5+4\rho}}2.
\]
The derivative observation also shows that checking
\(F_{n,k}(B-1)\ge0>F_{n,k}(B)\) proves the exact identity
\(r_{\rm quad}(n,k)=B-1\); no floating-point approximation is involved.

At the four challenge rates the limiting constants are, for orientation,
\[
\begin{array}{c|c|c}
 \rho&\alpha_\rho&\text{decimal value}\\ \hline
 1/2&(3-\sqrt7)/2&0.1771243444\ldots\\
 1/4&(3-\sqrt6)/2&0.2752551286\ldots\\
 1/8&(6-\sqrt{22})/4&0.3273960600\ldots\\
 1/16&(6-\sqrt{21})/4&0.3543560762\ldots
\end{array}
\]
The decimals are not used in any proof; the exact signs in
\eqref{eq:quadratic-row-certificates} locate the finite boundaries.

\subsection{Proth certificates and budget divisions}

We spell out the certificate format.  Suppose
\(p=u2^s+1\), where \(u\) is odd and \(u<2^s\), and suppose an integer
\(a_0\) satisfies
\[
                         a_0^{(p-1)/2}\equiv-1\pmod p.
\]
If a prime \(\ell\) divides \(p\), then the order of \(a_0\) modulo
\(\ell\) has exact \(2\)-part \(2^s\).  Hence \(2^s\mid\ell-1\), so
\(\ell\ge2^s+1\).  The inequality \(u<2^s\) gives
\(p=u2^s+1<2^{2s}+1\), and therefore \(2^s>\sqrt p-1\).
If \(p\) were composite it would have a prime divisor at most
\(\sqrt p\), contradicting the lower bound \(\ell\ge2^s+1>\sqrt p\).
Thus \(p\) is prime.  This is Proth's theorem in exactly the form needed
here; see \cite[Sec.~4.1]{Riesel}.

For a positive integer \(N\), write
\(\operatorname{bits}(N)=\lceil\log_2N\rceil\).  All certificate inputs
are collected below.  The equality in the last
column is a residue equality modulo the corresponding \(p\).
\[
\begin{array}{c|r|r|r|r|r}
 \rho&s&u&a_0&\operatorname{bits}(p)&
 a_0^{(p-1)/2}\bmod p\\ \hline
 1/2&92&26766274163673319604503&3&167&p-1\\
 1/4&93&41595378994516821279015&13&169&p-1\\
 1/8&95&24737346889219389259839&5&170&p-1\\
 1/16&97&13387194060291799253121&5&171&p-1
\end{array}
\]
Together with the decimal values of \(p\) printed before
\cref{prop:proth-row-check}, this table is a complete Proth certificate:
the reader checks the single exact identity \(p-1=u2^s\), the parity and
size of \(u\), and one modular exponentiation.  Modular exponentiation is
an integer operation and does not rely on a probabilistic primality test.

For a completely mechanical check, first compute
\(t=a_0^u\bmod p\) by binary square-and-multiply, then square \(t\)
exactly \(s-1\) times modulo \(p\).  The certificate accepts precisely
when the result is \(p-1\).  The algorithm uses
\(O(\log u+s)\) modular multiplications.  The remaining row checks use
only division with remainder.  Writing
\[
                         p=B2^{128}+r_p,
\]
the four exact remainders are
\[
\begin{array}{c|r}
 \rho&r_p\\ \hline
 1/2&1381541083842484386787422633985\\
 1/4&2921538492713497448761933168641\\
 1/8&2495687119199326634196634435585\\
 1/16&20440865928680199099134339186689
\end{array}
\]
and each lies strictly between \(0\) and
\(2^{128}=340282366920938463463374607431768211456\).  Therefore the
displayed \(B\) is exactly \(\floor{p/2^{128}}\).

Finally, \(n=2^e\) with \(e=41,42,43,44\), respectively, while
\(s=92,93,95,97\).  Thus \(n\mid2^s\mid p-1\), and the cyclic group
\(\F_p^\times\) contains a unique subgroup of order \(n\).  Each field
has the printed bit length, all smaller than \(256\), and
\(k=2^{40}\) gives the required four rates.  The exact boundary signs are
\[
\begin{array}{c|r|r}
 \rho&F_{n,k}(B-1)&F_{n,k}(B)\\ \hline
 1/2&5154112775168&-663955886271\\
 1/4&7590647904465&-3182321912768\\
 1/8&13908181940112&-6720484728007\\
 1/16&19335616403905&-20973145690236
\end{array}
\]
and hence \(B-1=r_{\rm quad}(n,k)\) in every row.

\begin{proposition}[Independent certificate verification]
\label{prop:appendix-certificate-verification}
The data in \cref{tab:prize-proth-rows} define four prime fields, four
power-of-two multiplicative subgroups of the stated lengths, and slope
budgets \(B=\floor{p2^{-128}}\).  In every row the exact safe set is
\([0,B/n)\), and the endpoint is unsafe.
\end{proposition}

\begin{proof}
The Proth calculation proves primality, the divisibility calculation
produces the subgroup, and the remainder calculation gives the target
budget.  The two signs locate the quadratic boundary at \(B-1\).
Apply \cref{cor:prize-window-compiler} and then
\cref{lem:appendix-endpoint}.  Every step is an equality or inequality
between the printed integers.
\end{proof}

\begin{thebibliography}{GGSW26}

\bibitem[ABF26]{ABF26}
G. Arnon, D. Boneh, and G. Fenzi,
\newblock Open problems in list decoding and correlated agreement,
\newblock IACR Cryptology ePrint Archive, Report 2026/680, 2026,
revised July 6, 2026.

\bibitem[ACFY25]{WHIR}
G. Arnon, A. Chiesa, G. Fenzi, and E. Yogev,
\newblock WHIR: Reed--Solomon proximity testing with super-fast verification,
\newblock in \emph{EUROCRYPT 2025}, 214--243;
IACR ePrint 2024/1586.

\bibitem[BS94]{BalogSzemeredi}
A. Balog and E. Szemer\'edi,
\newblock A statistical theorem of set addition,
\newblock \emph{Combinatorica} 14 (1994), no.~3, 263--268.

\bibitem[BCHKS25]{BCHKS25}
E. Ben-Sasson, D. Carmon, U. Hab\"ock, S. Kopparty, and S. Saraf,
\newblock On proximity gaps for Reed--Solomon codes,
\newblock IACR Cryptology ePrint Archive, Report 2025/2055, 2025.

\bibitem[BCIKS20]{BCIKS20}
E. Ben-Sasson, D. Carmon, Y. Ishai, S. Kopparty, and S. Saraf,
\newblock Proximity gaps for Reed--Solomon codes,
\newblock \emph{J. ACM} 70 (2023), no.~5, 1--57,
\newblock doi:10.1145/3614423; conference version in
\emph{61st IEEE Symposium on Foundations of Computer Science
(FOCS 2020)} and extended version, IACR ePrint 2020/654.

\bibitem[BCGM25]{BordageChiesaGuanManzur25}
S. Bordage, A. Chiesa, Z. Guan, and I. Manzur,
\newblock All polynomial generators preserve distance with mutual
correlated agreement,
\newblock IACR Cryptology ePrint Archive, Report 2025/2051, 2025,
revised May 2026; conference version in CCC 2026.

\bibitem[Cho26]{RSMCARepo}
P. Chojecki,
\newblock \texttt{rs-mca}: proof development and exact certificate ledgers,
\newblock source repository, pinned snapshot
\href{https://github.com/przchojecki/rs-mca/tree/02728b208ea785d02115ea967236aebf653b31ec}{\texttt{02728b208ea7}}, 2026,
\newblock \url{https://github.com/przchojecki/rs-mca}, accessed July 17, 2026.

\bibitem[CS25]{CS25}
E. Crites and A. Stewart,
\newblock On Reed--Solomon proximity gaps conjectures,
\newblock IACR Cryptology ePrint Archive, Report 2025/2046, 2025.

\bibitem[Dav00]{Davenport}
H. Davenport,
\newblock \emph{Multiplicative Number Theory}, third ed.,
\newblock Graduate Texts in Mathematics 74, Springer, New York, 2000.

\bibitem[PP26]{ProximityPrize}
{Ethereum Foundation},
\newblock The Proximity Prize: \$1M Reed--Solomon Challenge,
\newblock preliminary version, 2026,
\newblock \url{https://proximityprize.org/}, accessed July 17, 2026.

\bibitem[GYXK26]{GaoYangXuKan26}
Y. Gao, H. Yang, Y. Xu, and H. Kan,
\newblock List-decoding counterexamples yield lower bounds on mutual
correlated agreement error,
\newblock arXiv:2607.10572, version 1, July 2026.

\bibitem[Gow98]{GowersSzemeredi}
W. T. Gowers,
\newblock A new proof of Szemer\'edi's theorem for arithmetic progressions
of length four,
\newblock \emph{Geom. Funct. Anal.} 8 (1998), no.~3, 529--551.

\bibitem[GG25]{GoyalGuruswami25}
R. Goyal and V. Guruswami,
\newblock Optimal proximity gaps for subspace-design codes and (random)
Reed--Solomon codes,
\newblock IACR Cryptology ePrint Archive, Report 2025/2054, 2025,
revised March 2026.

\bibitem[GGSW26]{GoyalGuruswamiSunWootters26}
R. Goyal, V. Guruswami, Y. Sun, and M. Wootters,
\newblock Locality of curve-decoding and improved proximity gaps,
\newblock arXiv:2607.08516, 2026.

\bibitem[GMR+22]{GreenMatolcsiRuzsaShakanZhelezov}
B. Green, D. Matolcsi, I. Z. Ruzsa, G. Shakan, and D. Zhelezov,
\newblock A weighted Pr\'ekopa--Leindler inequality and sumsets with
quasicubes,
\newblock in \emph{Analysis at Large}, Springer, 2022, 125--129;
\newblock arXiv:2003.04077.

\bibitem[Hab25]{Habock25}
U. Hab\"ock,
\newblock A note on mutual correlated agreement for Reed--Solomon codes,
\newblock IACR Cryptology ePrint Archive, Report 2025/2110, 2025.

\bibitem[HLP24]{HLP24}
U. Hab\"ock, D. Levit, and S. Papini,
\newblock Circle STARKs,
\newblock IACR Cryptology ePrint Archive, Report 2024/278, 2024.

\bibitem[Jo26]{Jo26}
S. Jo,
\newblock Reed--Solomon mutual correlated agreement beyond the Johnson radius,
\newblock IACR Cryptology ePrint Archive, Report 2026/1432, 2026,
\newblock \url{https://eprint.iacr.org/2026/1432}.

\bibitem[Kam26]{Kambire26}
A. Kambir\'e,
\newblock Proximity gaps conjecture fails near capacity over prime fields,
\newblock arXiv:2604.09724, 2026.

\bibitem[KKH26]{KKH26}
D. Krachun, S. Kazanin, and U. Hab\"ock,
\newblock Failure of proximity gaps close to capacity,
\newblock IACR Cryptology ePrint Archive, Report 2026/782, 2026.

\bibitem[Rie94]{Riesel}
H. Riesel,
\newblock \emph{Prime Numbers and Computer Methods for Factorization},
second ed., Progress in Mathematics 126, Birkh{\"a}user, Boston, 1994.

\bibitem[YZ26]{YuanZhu26}
C. Yuan and R. Zhu,
\newblock A syndrome--space approach to proximity gaps and correlated
agreement for random linear codes and random Reed--Solomon codes,
\newblock arXiv:2605.07595, version 2, revised July 10, 2026.

\end{thebibliography}

\end{document}
